% Systèmes oscillants et propagation des ondes mécaniques — Physique Tle D — Cours signé OnBuch+
% Programme officiel MINESEC. Compiler avec : tectonic P14-systemes-oscillants-ondes-mecaniques.tex
\newcommand{\DOCMATIERE}{Physique}
\newcommand{\DOCNIVEAU}{Tle D}
\newcommand{\DOCMODULE}{Module 4 : Onde, matière et transformations nucléaires}
\newcommand{\DOCLECON}{P14}
\newcommand{\DOCTITRE}{Systèmes oscillants et propagation des ondes mécaniques}
% OnBuch+ — préambule « cours signés » ENRICHI (compilé avec tectonic / XeLaTeX)
% Système visuel « Pop » d'OnBuch+ : fond crème (jamais blanc), bordures encre
% épaisses, ombres « dures » décalées, titres Archivo Black en capitales.
% Version MATHÉMATIQUES : graphes (pgfplots/tikz), boîtes pédagogiques
% (définition, propriété, méthode, attention, le savais-tu, exercice résolu,
% exercice, TP, bilan), page de garde et signature OnBuch+ ancrée en bas.
%
% Macros attendues dans main.tex AVANT \input{preamble} :
%   \DOCMATIERE  \DOCNIVEAU  \DOCMODULE  \DOCLECON (numéro)  \DOCTITRE
\documentclass[11pt]{article}

\usepackage{fontspec}
\usepackage[french]{babel}
\usepackage[a4paper,margin=2cm,top=3.4cm,bottom=2.8cm,headheight=1.4cm,footskip=1.3cm]{geometry}
\usepackage{amsmath,amssymb,amsfonts}
\usepackage{mathtools} % \xmapsto, \coloneqq... souvent utilisés par les modèles
\usepackage{cancel} % simplifications barrées (bilans de forces)
% \abs{z} (module, valeur absolue) et \norm{u} (norme), utilisés par les modèles
\providecommand{\abs}{}\let\abs\relax\DeclarePairedDelimiter{\abs}{\lvert}{\rvert}
\providecommand{\norm}{}\let\norm\relax\DeclarePairedDelimiter{\norm}{\lVert}{\rVert}
\usepackage[table]{xcolor} % [table] : fournit aussi \rowcolors (alternance de lignes)
\usepackage{pagecolor}
\usepackage{tikz}
\usetikzlibrary{arrows.meta,positioning,calc,shapes.geometric,shapes.misc,shapes.arrows,decorations.pathmorphing,decorations.markings,patterns,shadows,mindmap,trees,fit,backgrounds,matrix,intersections,angles,quotes}
\usepackage{pgfplots}
\usepackage[version=4]{mhchem} % équations nucléaires \ce{^{235}_{92}U}
\usepackage[european,siunitx]{circuitikz} % schémas électriques
\pgfplotsset{compat=1.18}
\usepgfplotslibrary{fillbetween} % aires entre courbes (calcul intégral)
\usepackage{enumitem}
\usepackage{fancyhdr}
% [most] charge toutes les bibliothèques tcolorbox (listings, minted, raster,
% documentation...) et gonfle inutilement la mémoire TeX sur un document long
% (~300 boîtes) : on ne charge que ce qui est réellement utilisé.
\usepackage[skins,breakable]{tcolorbox}
\usepackage{graphicx}
\usepackage{colortbl}
\usepackage{booktabs}
\usepackage{tabularx}
\usepackage{diagbox} % case d'en-tête diagonale des tableaux à double entrée
% Colonnes extensibles centrée / gauche / droite (C, L, R), souvent utilisées par les modèles
\newcolumntype{C}{>{\centering\arraybackslash}X}
\newcolumntype{L}{>{\raggedright\arraybackslash}X}
\newcolumntype{R}{>{\raggedleft\arraybackslash}X}
\usepackage{array}
\usepackage{multicol}
\usepackage{siunitx}
% parse-numbers=false : accepte aussi \SI{2,5 \times 10^{-3}}{...}
\sisetup{locale=FR,output-decimal-marker={,},group-minimum-digits=4,parse-numbers=false}
\DeclareSIUnit\ohm{\ensuremath{\symup{\Omega}}} % U+2126 absent de la police
\DeclareSIUnit\division{div} % réglages d'oscilloscope (ms/div, V/div)
\usepackage{qrcode}
% \degree : commande fréquemment utilisée par les modèles pour un angle ou une
% température en dehors de \SI{}{} (non fournie par les paquets chargés).
\providecommand{\degree}{^{\circ}}
% \qed (fin de démonstration), utilisé par les modèles sans amsthm
\providecommand{\qed}{\hfill$\square$}
% \barre : double barre des valeurs interdites dans un tableau de variations (tkz-tab)
\providecommand{\barre}{\ensuremath{\|}}
% environnement proof (amsthm non chargé) utilisé par les modèles
\ifdefined\proof\else\newenvironment{proof}[1][Démonstration]{\par\noindent\textit{#1.}\ }{\qed\par}\fi
\usepackage{lastpage}
\usepackage{hyperref}
\hypersetup{hidelinks}

% ============================================================
% Palette « Pop » OnBuch+ — mêmes hex que la classe P (plus_ui.dart)
% ============================================================
\definecolor{poporange}{HTML}{F59321}
\definecolor{popdark}{HTML}{E07A0C}
\definecolor{popcream}{HTML}{FFF6E8}
\definecolor{popink}{HTML}{1C1714}
\definecolor{poppurple}{HTML}{7A5AE0}
\definecolor{popgreen}{HTML}{2E9E6A}
\definecolor{poppink}{HTML}{E0457B}
\definecolor{popblue}{HTML}{2F7DE1}
\definecolor{popgold}{HTML}{C98A12}
\definecolor{popmuted}{HTML}{8A7F76}
\definecolor{popline}{HTML}{EADFD2}
\definecolor{popgray}{HTML}{8A7F76}
\colorlet{popgrey}{popgray}
% Alias tolérants (le modèle nomme parfois les couleurs différemment)
\colorlet{popwhite}{white}
\colorlet{popblack}{popink}
\colorlet{popred}{poppink}
\colorlet{popyellow}{popgold}
% Teintes claires pour fonds de tableaux / zones
\colorlet{popblueL}{popblue!10!white}
\colorlet{poporangeL}{poporange!14!white}
\colorlet{popgreenL}{popgreen!12!white}
\colorlet{poppurpleL}{poppurple!12!white}
\colorlet{poppinkL}{poppink!12!white}
\colorlet{popgoldL}{popgold!14!white}

\pagecolor{popcream}

% ============================================================
% Typographie
% ============================================================
\defaultfontfeatures{Ligatures=TeX}
\setmainfont{PlusJakartaSans-Medium.ttf}[Path=../../../fonts/,BoldFont=PlusJakartaSans-Bold.ttf]
% Maths en sans-serif (Fira Math) avec chiffres et lettres droites de Plus
% Jakarta, pour que les formules se fondent dans le texte.
\usepackage[math-style=ISO,bold-style=ISO]{unicode-math}
\setmathfont{FiraMath-Regular.otf}[Path=../../../fonts/]
\setmathfont{PlusJakartaSans-Medium.ttf}[Path=../../../fonts/,range={up/num,up/{Latin,latin},\mathup}]
\usepackage{icomma} % virgule décimale sans espace en mode math
% Caractères mathématiques Unicode absents des polices de texte (Plus Jakarta,
% Archivo Black) : rendus par leur équivalent mathématique, en texte comme en
% formule — sinon ils s'affichent en carré vide (ex. « Arithmétique dans ℤ »).
\usepackage{newunicodechar}
\newunicodechar{ℝ}{\ensuremath{\mathbb{R}}}
\newunicodechar{ℤ}{\ensuremath{\mathbb{Z}}}
\newunicodechar{ℂ}{\ensuremath{\mathbb{C}}}
\newunicodechar{ℕ}{\ensuremath{\mathbb{N}}}
\newunicodechar{ℚ}{\ensuremath{\mathbb{Q}}}
\newunicodechar{𝔻}{\ensuremath{\mathbb{D}}}
\newunicodechar{α}{\ensuremath{\alpha}}
\newunicodechar{β}{\ensuremath{\beta}}
\newunicodechar{γ}{\ensuremath{\gamma}}
\newunicodechar{δ}{\ensuremath{\delta}}
\newunicodechar{ε}{\ensuremath{\varepsilon}}
\newunicodechar{θ}{\ensuremath{\theta}}
\newunicodechar{λ}{\ensuremath{\lambda}}
\newunicodechar{μ}{\ensuremath{\mu}}
\newunicodechar{σ}{\ensuremath{\sigma}}
\newunicodechar{τ}{\ensuremath{\tau}}
\newunicodechar{φ}{\ensuremath{\varphi}}
\newunicodechar{ψ}{\ensuremath{\psi}}
\newunicodechar{ω}{\ensuremath{\omega}}
\newunicodechar{Σ}{\ensuremath{\Sigma}}
% majuscules produites par \MakeUppercase dans les titres (α-aminés -> Α)
\newunicodechar{Α}{\ensuremath{\alpha}}
\newunicodechar{Β}{\ensuremath{\beta}}
\newunicodechar{Γ}{\ensuremath{\Gamma}}
\newunicodechar{Δ}{\ensuremath{\Delta}}
\newunicodechar{Ε}{\ensuremath{\varepsilon}}
\newunicodechar{Θ}{\ensuremath{\Theta}}
\newunicodechar{Λ}{\ensuremath{\Lambda}}
\newunicodechar{Μ}{\ensuremath{\mu}}
\newunicodechar{Τ}{\ensuremath{\tau}}
\newunicodechar{Φ}{\ensuremath{\Phi}}
\newunicodechar{Ψ}{\ensuremath{\Psi}}
\newunicodechar{Ω}{\ensuremath{\Omega}}
\newunicodechar{↦}{\ensuremath{\mapsto}}
\newunicodechar{∈}{\ensuremath{\in}}
\newunicodechar{∉}{\ensuremath{\notin}}
\newunicodechar{∧}{\ensuremath{\wedge}}
\newunicodechar{⇔}{\ensuremath{\Leftrightarrow}}
\newunicodechar{⟺}{\ensuremath{\Longleftrightarrow}}
\newunicodechar{⇒}{\ensuremath{\Rightarrow}}
\newunicodechar{⟹}{\ensuremath{\Longrightarrow}}
\newunicodechar{∘}{\ensuremath{\circ}}
\newunicodechar{∪}{\ensuremath{\cup}}
\newunicodechar{∩}{\ensuremath{\cap}}
\newunicodechar{≡}{\ensuremath{\equiv}}
\newunicodechar{⊂}{\ensuremath{\subset}}
\newunicodechar{⊥}{\ensuremath{\perp}}
\newunicodechar{∃}{\ensuremath{\exists}}
\newunicodechar{∀}{\ensuremath{\forall}}
\newunicodechar{∛}{\ensuremath{\sqrt[3]{\,}}}
\newunicodechar{∜}{\ensuremath{\sqrt[4]{\,}}}
\newunicodechar{⃗}{\ensuremath{{}^{\to}}}
\newunicodechar{ⁿ}{\ifmmode^{n}\else\textsuperscript{\ensuremath{n}}\fi}
\newunicodechar{ˣ}{\ifmmode^{x}\else\textsuperscript{\ensuremath{x}}\fi}
\newunicodechar{ᵇ}{\ifmmode^{b}\else\textsuperscript{\ensuremath{b}}\fi}
\newunicodechar{ʳ}{\ifmmode^{r}\else\textsuperscript{\ensuremath{r}}\fi}
\newunicodechar{ᵘ}{\ifmmode^{u}\else\textsuperscript{\ensuremath{u}}\fi}
\newunicodechar{ʸ}{\ifmmode^{y}\else\textsuperscript{\ensuremath{y}}\fi}
\newunicodechar{ᵃ}{\ifmmode^{a}\else\textsuperscript{\ensuremath{a}}\fi}
\newunicodechar{ᵉ}{\ifmmode^{e}\else\textsuperscript{\ensuremath{e}}\fi}
\newunicodechar{ᵏ}{\ifmmode^{k}\else\textsuperscript{\ensuremath{k}}\fi}
\newunicodechar{ᵖ}{\ifmmode^{p}\else\textsuperscript{\ensuremath{p}}\fi}
\newunicodechar{ᴺ}{\ifmmode^{N}\else\textsuperscript{\ensuremath{N}}\fi}
\newunicodechar{ᶜ}{\ifmmode^{c}\else\textsuperscript{\ensuremath{c}}\fi}
\newunicodechar{ʰ}{\ifmmode^{h}\else\textsuperscript{\ensuremath{h}}\fi}
\newunicodechar{ᵠ}{\ifmmode^{q}\else\textsuperscript{\ensuremath{q}}\fi}
\newunicodechar{ₙ}{\ifmmode_{n}\else\textsubscript{\ensuremath{n}}\fi}
\newunicodechar{ₐ}{\ifmmode_{a}\else\textsubscript{\ensuremath{a}}\fi}
\newunicodechar{ᵢ}{\ifmmode_{i}\else\textsubscript{\ensuremath{i}}\fi}
\newunicodechar{ᵣ}{\ifmmode_{r}\else\textsubscript{\ensuremath{r}}\fi}
\newunicodechar{ₖ}{\ifmmode_{k}\else\textsubscript{\ensuremath{k}}\fi}
\newunicodechar{ᵦ}{\ifmmode_{\beta}\else\textsubscript{\ensuremath{\beta}}\fi}
\newunicodechar{ₓ}{\ifmmode_{x}\else\textsubscript{\ensuremath{x}}\fi}
\newfontfamily\popdisplay{ArchivoBlack-Regular.ttf}[Path=../../../fonts/]
\newfontfamily\poplogo{SpaceGrotesk-Bold.ttf}[Path=../../../fonts/]
\newfontfamily\popsemi{PlusJakartaSans-SemiBold.ttf}[Path=../../../fonts/]
\newfontfamily\popxbold{PlusJakartaSans-ExtraBold.ttf}[Path=../../../fonts/]
\newfontfamily\popxboldls{PlusJakartaSans-ExtraBold.ttf}[Path=../../../fonts/,LetterSpace=3.0]

\setlist[enumerate]{leftmargin=1.7em,itemsep=5pt,topsep=4pt}
\setlist[itemize]{leftmargin=1.7em,itemsep=4pt,topsep=4pt,label={\color{poporange}\textbullet}}
\setlength{\parindent}{0pt}
\setlength{\parskip}{5pt}
\renewcommand{\arraystretch}{1.25}

% pgfplots : style Pop par défaut
\pgfplotsset{
  every axis/.append style={axis line style={popink,line width=1pt},tick style={popink},
    grid style={popline,line width=0.5pt},label style={font=\small},tick label style={font=\footnotesize}},
  popcurve/.style={poporange,line width=1.8pt,no markers,smooth},
}
% Même style disponible dans un \draw TikZ simple (hors pgfplots)
\tikzset{popcurve/.style={poporange,line width=1.8pt}}
\tikzset{popgrid/.style={popline,very thin}}
\colorlet{popcurve}{poporange} % le nom du style est parfois utilisé comme couleur

% ============================================================
% Pill / squiggle
% ============================================================
\newcommand{\poppill}[3]{%
  \tikz[baseline=-0.55ex]\node[draw=popink,line width=1.2pt,rounded corners=2.6pt,%
    fill=#2,inner xsep=6.5pt,inner ysep=3pt]{%
    \popxboldls\fontsize{7}{7}\selectfont\color{#3}\MakeUppercase{#1}};%
}
\newcommand{\popsquiggle}[1]{%
  \tikz[baseline=-0.4ex]{
    \draw[#1,line width=2.6pt,line cap=round]
      plot [smooth] coordinates {(0,0.09) (0.34,0.01) (0.68,0.21) (1.02,0.01) (1.36,0.10)};
  }%
}
% Mot-clé surligné (marqueur orange sous le texte)
\newcommand{\cle}[1]{{\popxbold\color{popdark}#1}}

% ============================================================
% En-tête / pied
% ============================================================
\pagestyle{fancy}
\fancyhf{}
\renewcommand{\headrulewidth}{0pt}
\renewcommand{\footrulewidth}{0pt}
\lhead{\raisebox{-0.15ex}{\poplogo\fontsize{13}{13}\selectfont\color{popink}OnBuch\color{poporange}+}}
\rhead{\poppill{\DOCMATIERE\ \textbullet\ \DOCNIVEAU\ \textbullet\ Leçon \DOCLECON}{white}{popink}}
\lfoot{\popxbold\fontsize{7.6}{7.6}\selectfont\color{popmuted}\MakeUppercase{Cours signé OnBuch+}\ \textbullet\ {\popsemi\DOCTITRE}}
\rfoot{\tikz[baseline=-0.6ex]\node[rounded corners=8.5pt,fill=popink,minimum height=17pt,minimum width=17pt,inner xsep=5pt,inner sep=0]{\popxbold\fontsize{8}{8}\selectfont\color{popcream}\thepage\,/\,\pageref*{LastPage}};}

% ============================================================
% Titres
% ============================================================
\newcommand{\chaptitre}[2]{%
  \begin{tcolorbox}[enhanced,colback=poporange,colframe=popink,boxrule=2.6pt,arc=11pt,
      drop shadow=popink,left=15pt,right=15pt,top=12pt,bottom=11pt,before skip=3pt,after skip=18pt]
    \poppill{#2}{popink}{popcream}\par\vspace{9pt}
    {\raggedright\hyphenpenalty=10000\exhyphenpenalty=10000\popdisplay\fontsize{22}{24}\selectfont\color{popcream}\MakeUppercase{#1}\par}
  \end{tcolorbox}
}
% Section numérotée : pastille orange avec le numéro + titre Archivo
\newcounter{popsec}
\newcommand{\coursec}[1]{%
  \stepcounter{popsec}\setcounter{popsub}{0}%
  \par\vspace{16pt}\noindent
  \begin{minipage}{\linewidth}
  \tikz[baseline=-0.7ex]\node[draw=popink,line width=1.6pt,fill=poporange,rounded corners=4pt,minimum size=22pt,inner sep=0]{\popdisplay\fontsize{12}{12}\selectfont\color{popcream}\thepopsec};%
  \hspace{8pt}{\popdisplay\fontsize{15}{17}\selectfont\color{popink}\MakeUppercase{#1}}\par
  \vspace{4pt}\noindent{\color{poporange}\rule{36pt}{3.6pt}}
  \end{minipage}\par\nopagebreak\vspace{8pt}
}
\newcounter{popsub}
\newcommand{\courssub}[1]{%
  \stepcounter{popsub}%
  \par\vspace{9pt}\noindent{\popxbold\fontsize{11.5}{13}\selectfont\color{popink}{\color{poporange}\thepopsec.\thepopsub}\hspace{6pt}#1}\par\nopagebreak\vspace{4pt}
}

% ============================================================
% Boîtes pédagogiques — même recette : fond blanc, bordure épaisse colorée,
% ombre dure encre, badge plein. Titre optionnel [#1].
% ============================================================
\tcbset{popbase/.style={breakable,enhanced,colback=white,boxrule=2.3pt,arc=9pt,
  shadow={3pt}{-3pt}{0pt}{popink},left=14pt,right=14pt,top=11pt,bottom=11pt,before skip=11pt,after skip=15pt,
  fontupper=\popsemi\fontsize{10.6}{14.5}\selectfont\color{popink},
  fontlower=\popsemi\fontsize{10.6}{14.5}\selectfont\color{popink}}}
% Coche dessinée (le glyphe ✓ n'existe pas dans les polices du document)
\newcommand{\popcheck}[1]{\tikz[baseline=-0.5ex]\draw[#1,line width=1.6pt,line cap=round,line join=round](-0.13,0.0)--(-0.04,-0.09)--(0.13,0.1);}
\providecommand{\coche}{\popcheck{popgreen}}
\providecommand{\resultat}[1]{\par\smallskip\noindent\fcolorbox{popgreen}{popgreenL}{\parbox{\dimexpr\linewidth-2\fboxsep-2\fboxrule\relax}{\textbf{Résultat.} #1}}\par\smallskip}
\newcommand{\popboxhead}[3]{\poppill{#1}{#2}{white}\ifx\relax#3\relax\else\hspace{6pt}{\popxbold\fontsize{10.6}{12}\selectfont\color{popink}#3}\fi\par\vspace{7pt}}

\newtcolorbox{aretenir}[1][]{popbase,colframe=popgreen,before upper={\popboxhead{À retenir}{popgreen}{#1}}}
\newtcolorbox{exemplebox}[1][]{popbase,colframe=poporange,before upper={\popboxhead{Exemple}{poporange}{#1}}}
\newtcolorbox{definition}[1][]{popbase,colframe=popblue,before upper={\popboxhead{Définition}{popblue}{#1}}}
\newtcolorbox{propriete}[1][]{popbase,colframe=poppurple,before upper={\popboxhead{Propriété}{poppurple}{#1}}}
\newtcolorbox{methode}[1][]{popbase,colframe=popgold,before upper={\popboxhead{Méthode}{popgold}{#1}}}
\newtcolorbox{attention}[1][]{popbase,colframe=poppink,before upper={\popboxhead{Attention}{poppink}{#1}}}
\newtcolorbox{experience}[1][]{popbase,colframe=popblue,colback=popblueL,before upper={\popboxhead{Expérience}{popblue}{#1}}}
% « Le savais-tu ? » : carte violette pleine (contexte, culture, Cameroun)
\newtcolorbox{savaistu}[1][]{popbase,colback=poppurple,colframe=popink,
  fontupper=\popsemi\fontsize{10.4}{14.2}\selectfont\color{white},
  before upper={\poppill{Le savais-tu\,?}{white}{poppurple}\ifx\relax#1\relax\else\hspace{6pt}{\popxbold\color{white}#1}\fi\par\vspace{7pt}}}
% Exercice résolu : énoncé en haut, \tcblower puis la solution sur fond crème
\newcounter{popexo}\newcounter{popexr}
\newtcolorbox{exoresolu}[1][]{popbase,colframe=poporange,colbacklower=popcream,
  segmentation style={popink,line width=1.2pt,dashed},
  before upper={\stepcounter{popexr}\popboxhead{Exercice résolu \thepopexr}{poporange}{#1}},
  before lower={\poppill{Solution}{popink}{popcream}\par\vspace{6pt}}}
% Exercice (non résolu en place) — la correction est dans la partie Corrigés
\newtcolorbox{exercice}[1][]{popbase,colframe=popink,
  before upper={\stepcounter{popexo}\popboxhead{Exercice \thepopexo}{popink}{#1}}}
% Corrigé
\newtcolorbox{corrige}[1][]{popbase,colframe=popgreen,colback=popgreenL,
  before upper={\popboxhead{Corrigé}{popgreen}{#1}}}
% Objectifs / prérequis / bilan
\newtcolorbox{objectifs}{popbase,colframe=popink,colback=poporangeL,
  before upper={\popboxhead{Objectifs de la leçon}{popink}{}}}
\newtcolorbox{prerequis}{popbase,colframe=popink,colback=popblueL,
  before upper={\popboxhead{Prérequis}{popblue}{}}}
\newtcolorbox{bilan}{popbase,colframe=popink,colback=white,boxrule=2.8pt,
  before upper={\popboxhead{Fiche bilan}{poporange}{L'essentiel à connaître pour l'examen}}}
% Figure encadrée avec légende
\newtcolorbox{popfigure}[1][]{enhanced,breakable=false,colback=white,colframe=popline,boxrule=1.4pt,arc=8pt,
  left=10pt,right=10pt,top=10pt,bottom=8pt,before skip=10pt,after skip=12pt,halign=center,
  lower separated=false,halign lower=center,
  fontlower=\popsemi\fontsize{9}{11.5}\selectfont\color{popmuted},
  #1}
\newcommand{\legende}[1]{\par\vspace{6pt}{\popsemi\fontsize{9}{11.5}\selectfont\color{popmuted}#1}}

% ============================================================
% Signature OnBuch+
% ============================================================
\newcommand{\poplogoinline}[1]{{\poplogo\fontsize{#1}{#1}\selectfont\color{popink}OnBuch\color{poporange}+}}
\newcommand{\signatureonbuch}{%
  \begin{tcolorbox}[enhanced,colback=popink,colframe=popink,boxrule=2.2pt,arc=10pt,
      drop shadow=poporange,left=14pt,right=14pt,top=10pt,bottom=10pt,before skip=0pt,after skip=0pt]
    \tikz[baseline=-0.7ex]\node[circle,fill=popgreen,draw=popcream,line width=1.3pt,minimum size=22pt,inner sep=0]{\popcheck{white}};%
    \hspace{9pt}\begin{minipage}[c]{0.62\linewidth}
      {\popxbold\fontsize{12}{13}\selectfont\color{popcream}Signé\ \textbullet\ {\poplogo OnBuch\color{poporange}+}}\par\vspace{2pt}
      {\raggedright\popsemi\fontsize{8}{10.5}\selectfont\color{popline}\MakeUppercase{Équipe pédagogique OnBuch+}\\\MakeUppercase{Conforme au programme officiel MINESEC}\par}
    \end{minipage}\hfill
    {\poplogo\fontsize{22}{22}\selectfont\color{popcream}OnBuch\color{poporange}+}
  \end{tcolorbox}%
}

% ============================================================
% Page de garde
% #1 titre · #2 matière · #3 niveau · #4 module · #5 n° leçon · #6 durée ·
% #7 description / objectifs (texte ou itemize)
% ============================================================
\newcommand{\pagegarde}[7]{%
  \thispagestyle{empty}
  \newgeometry{a4paper,margin=1.7cm,top=1.6cm,bottom=1.4cm}%
  \begin{tikzpicture}[remember picture,overlay]
    \fill[poporange!18] ([xshift=-3.2cm,yshift=-3.6cm]current page.north east) circle (4.6cm);
    \fill[poppurple!14] ([xshift=2.4cm,yshift=7.6cm]current page.south west) circle (3.8cm);
  \end{tikzpicture}%
  {\poplogo\fontsize{30}{30}\selectfont\color{popink}OnBuch\color{poporange}+}\hfill
  \raisebox{0.6ex}{\poppill{Cours signé \textbullet\ Premium}{popink}{popcream}}\par
  \vspace{1pt}\popsquiggle{poppurple}\par
  \vspace{8pt}
  \begin{tcolorbox}[enhanced,colback=poporange,colframe=popink,boxrule=2.8pt,arc=13pt,
      drop shadow=popink,left=18pt,right=18pt,top=12pt,bottom=13pt,after skip=10pt]
    \poppill{#2 \textbullet\ #3}{popink}{popcream}\hspace{5pt}\poppill{#4}{white}{popink}\par\vspace{8pt}
    {\popdisplay\fontsize{14}{14}\selectfont\color{popink}LEÇON #5}\par\vspace{4pt}
    {\raggedright\hyphenpenalty=10000\exhyphenpenalty=10000\popdisplay\fontsize{25}{27}\selectfont\color{popcream}\MakeUppercase{#1}\par}
    \vspace{8pt}
    {\popxbold\fontsize{9.5}{9.5}\selectfont\color{popink}\MakeUppercase{Durée conseillée : #6}}
  \end{tcolorbox}
  \begin{tcolorbox}[enhanced,colback=white,colframe=popink,boxrule=2.2pt,arc=10pt,
      drop shadow=popink,left=16pt,right=16pt,top=10pt,bottom=10pt,after skip=8pt]
    \poppill{Dans cette leçon}{poppurple}{white}\par\vspace{5pt}
    \popsemi\fontsize{9.3}{12.6}\selectfont\color{popink}#7
  \end{tcolorbox}
  \noindent\begin{minipage}[t]{0.487\linewidth}
    \begin{tcolorbox}[enhanced,colback=popgreen,colframe=popink,boxrule=2.2pt,arc=10pt,
        drop shadow=popink,left=11pt,right=11pt,top=10pt,bottom=10pt,height=3.1cm,valign=center]
      \tcbox[colback=white,colframe=popink,boxrule=1.3pt,arc=5pt,boxsep=0pt,nobeforeafter,
        left=3pt,right=3pt,top=3pt,bottom=3pt,valign=center]{\qrcode[height=1.9cm]{https://play.google.com/store/apps/details?id=com.luvvix.onbuch&hl=fr}}%
      \hspace{8pt}\begin{minipage}[c]{0.5\linewidth}
        {\popdisplay\fontsize{11}{12.5}\selectfont\color{white}\MakeUppercase{Télécharge OnBuch}}\par\vspace{4pt}
        {\raggedright\popsemi\fontsize{8.4}{11}\selectfont\color{white}Gratuit \textbullet\ Google Play\\Tuteur IA, annales, concours, résultats\par}
      \end{minipage}
    \end{tcolorbox}
  \end{minipage}\hfill
  \begin{minipage}[t]{0.487\linewidth}
    \begin{tcolorbox}[enhanced,colback=poppurple,colframe=popink,boxrule=2.2pt,arc=10pt,
        drop shadow=popink,left=11pt,right=11pt,top=10pt,bottom=10pt,height=3.1cm,valign=center]
      \tikz[baseline=-0.7ex]\node[circle,fill=white,draw=popink,line width=1.3pt,minimum size=1.6cm,inner sep=0]{\popdisplay\fontsize{24}{24}\selectfont\color{poppurple}+};%
      \hspace{8pt}\begin{minipage}[c]{0.6\linewidth}
        {\popdisplay\fontsize{11}{12.5}\selectfont\color{white}\MakeUppercase{Passe à OnBuch+}}\par\vspace{4pt}
        {\raggedright\popsemi\fontsize{8.4}{11}\selectfont\color{white}Tous les cours signés, Tuteur IA illimité, fiches PDF — onglet OnBuch+ de l'app\par}
      \end{minipage}
    \end{tcolorbox}
  \end{minipage}
  \vspace{8pt}
  \popcontenu
  \vfill
  \signatureonbuch
  \restoregeometry
  \newpage
}

% Rangée « Ce cours contient » (page de garde)
\newcommand{\poptile}[3]{% #1 couleur #2 gros texte #3 légende
  \begin{tcolorbox}[enhanced,colback=#1,colframe=popink,boxrule=2pt,arc=9pt,shadow={2.5pt}{-2.5pt}{0pt}{popink},
      left=6pt,right=6pt,top=9pt,bottom=9pt,halign=center,valign=center,height=1.75cm,width=\linewidth,nobeforeafter]
    {\popdisplay\fontsize{12.5}{13}\selectfont\color{white}\MakeUppercase{#2}}\par\vspace{3pt}
    {\centering\popsemi\fontsize{7.8}{9.5}\selectfont\color{white}#3\par}
  \end{tcolorbox}}
\newcommand{\popcontenu}{%
  \par\noindent{\popxboldls\fontsize{7.5}{7.5}\selectfont\color{popmuted}CE COURS SIGNÉ CONTIENT}\par\vspace{7pt}
  \noindent\begin{minipage}[t]{0.235\linewidth}\poptile{popblue}{Cours}{complet, illustré, pas à pas}\end{minipage}\hfill
  \begin{minipage}[t]{0.235\linewidth}\poptile{poporange}{Méthodes}{et exercices résolus}\end{minipage}\hfill
  \begin{minipage}[t]{0.235\linewidth}\poptile{poppink}{Exercices}{type examen + corrigés}\end{minipage}\hfill
  \begin{minipage}[t]{0.235\linewidth}\poptile{popgreen}{Fiche bilan}{l'essentiel à retenir}\end{minipage}\par}

% Sommaire
\newtcolorbox{sommaire}{enhanced,colback=white,colframe=popink,boxrule=2.2pt,arc=10pt,
  drop shadow=popink,left=18pt,right=18pt,top=13pt,bottom=13pt,before skip=6pt,after skip=20pt,
  before upper={\poppill{Sommaire}{poppurple}{white}\par\vspace{9pt}\popsemi\fontsize{10.6}{14.5}\selectfont\color{popink}}}

% Signature de fin de cours — poussée tout en bas de la dernière page
\newcommand{\findecours}{%
  \par\vfill\nopagebreak
  \begin{center}\popsquiggle{poporange}\end{center}\vspace{4pt}
  \signatureonbuch
}

%% FIGFIX-BEGIN v3 — correctifs globaux des figures (docs/onbuch-generation/tools/figfix)
\usepackage{adjustbox}\usepackage{etoolbox}
% toute figure TikZ/pgfplots plus large que la ligne est réduite à la largeur disponible
\BeforeBeginEnvironment{tikzpicture}{\begin{adjustbox}{max width=\linewidth}}
\AfterEndEnvironment{tikzpicture}{\end{adjustbox}}
% légendes pgfplots lisibles par-dessus les courbes
\pgfplotsset{every axis/.append style={legend style={fill=white,fill opacity=0.92,text opacity=1,draw=black!15,inner sep=3pt}}}
% étiquettes d'arêtes (syntaxe edge["..."]) avec fond blanc
\tikzset{every edge quotes/.append style={fill=white,inner sep=1.2pt}}
% bulles de cartes mentales : texte à la ligne dans la bulle, bulle un peu plus grande
\tikzset{every concept/.append style={text width=2.0cm,align=center,minimum size=2.3cm,inner sep=2pt}}
%% FIGFIX-END
\begin{document}

\pagegarde{Systèmes oscillants et propagation des ondes mécaniques}{Physique}{Tle D}{Module 4 : Onde, matière et transformations nucléaires}{P14}{5 h}{Cette leçon étudie les phénomènes périodiques sinusoïdaux (représentation temporelle et de Fresnel, stroboscopie) puis les ondes mécaniques progressives, transversales ou longitudinales. Elle établit la double périodicité temporelle et spatiale des ondes sinusoïdales et l'équation du mouvement d'un point du milieu.
\vspace{4pt}
\begin{multicols}{2}\small
\begin{itemize}
\item À la fin de cette leçon, je sais définir et mesurer l'amplitude, la période, la fréquence et la phase d'une grandeur sinusoïdale
\item À la fin de cette leçon, je sais exploiter un oscillogramme pour déterminer période, amplitude et déphasage
\item À la fin de cette leçon, je sais interpréter les résultats d'un éclairage stroboscopique (immobilité apparente, mouvement ralenti apparent)
\item À la fin de cette leçon, je sais représenter une grandeur sinusoïdale par un vecteur de Fresnel et additionner deux grandeurs de même fréquence
\item À la fin de cette leçon, je sais différencier une onde transversale d'une onde longitudinale et définir la célérité
\item À la fin de cette leçon, je sais exploiter la double périodicité λ = cT et la relation de retard θ = d/c
\end{itemize}\end{multicols}}

\begin{sommaire}
\begin{multicols}{2}
\begin{enumerate}
\item Généralités sur les phénomènes périodiques sinusoïdaux
\item Stroboscopie
\item Représentation de Fresnel
\item Ondes mécaniques progressives
\item Onde progressive sinusoïdale : double périodicité
\item Synthèse et méthodes de résolution
\item Activité d'intégration
\item Méthodes et exercices résolus
\item Exercices
\item Corrigés des exercices
\item Fiche bilan
\end{enumerate}
\end{multicols}
\end{sommaire}

\begin{objectifs}
\begin{itemize}
\item À la fin de cette leçon, je sais définir et mesurer l'amplitude, la période, la fréquence et la phase d'une grandeur sinusoïdale
\item À la fin de cette leçon, je sais exploiter un oscillogramme pour déterminer période, amplitude et déphasage
\item À la fin de cette leçon, je sais interpréter les résultats d'un éclairage stroboscopique (immobilité apparente, mouvement ralenti apparent)
\item À la fin de cette leçon, je sais représenter une grandeur sinusoïdale par un vecteur de Fresnel et additionner deux grandeurs de même fréquence
\item À la fin de cette leçon, je sais différencier une onde transversale d'une onde longitudinale et définir la célérité
\item À la fin de cette leçon, je sais exploiter la double périodicité λ = cT et la relation de retard θ = d/c
\item À la fin de cette leçon, je sais établir l'équation du mouvement d'un point du milieu à partir de celle de la source
\item À la fin de cette leçon, je sais comparer l'état vibratoire de deux points (points en phase, en opposition de phase)
\end{itemize}
\end{objectifs}

\begin{prerequis}
\begin{itemize}
\item Fonctions trigonométriques sinus et cosinus, cercle trigonométrique (Première C)
\item Mouvement rectiligne uniforme : relation v = d/Δt (Seconde)
\item Oscillateurs mécaniques : période propre, équation horaire x(t) = Xm cos(ωt + φ) (Tle, leçon sur les oscillations)
\item Lecture d'un graphique et exploitation d'une échelle (Première)
\end{itemize}
\end{prerequis}

\newpage

\coursec{Généralités sur les phénomènes périodiques sinusoïdaux}

À Douala, un technicien règle l'éclairage stroboscopique d'une salle de spectacle : sous certains flashs, les pales du ventilateur semblent immobiles ou tourner au ralenti. Plus loin, un sismologue de l'IRGM détecte les ondes d'un séisme enregistré à la station de Foumban et en déduit la distance de l'épicentre grâce au retard des vibrations. Ces deux phénomènes, en apparence très différents, reposent sur les mêmes mathématiques : les grandeurs sinusoïdales et la propagation des ondes mécaniques. Comment décrire, représenter et additionner ces vibrations ? Comment prévoir l'état vibratoire d'un point quelconque d'un milieu traversé par une onde ? Telle est la problématique de cette leçon, dont nous posons ici les fondations : le langage des phénomènes périodiques.

\courssub{Phénomènes périodiques : observation et définitions}

\textbf{Observation.} Observe autour de toi : le pendule d'une vieille horloge qui balance, la lame d'une règle coincée sur le bord d'une table que l'on fait vibrer, la tension du secteur ENEO qui alimente ta maison, le piston d'une mototaxi qui monte et descend. Tous ces mouvements ont un point commun : \emph{ils se répètent identiques à eux-mêmes à intervalles de temps égaux}.

\begin{definition}[Phénomène périodique]
Un phénomène est dit \cle{périodique} s'il se reproduit identique à lui-même à intervalles de temps réguliers. La plus petite durée au bout de laquelle le phénomène se reproduit à l'identique est appelée \cle{période}, notée $T$ ; elle s'exprime en secondes (s).
\end{definition}

\begin{definition}[Fréquence et pulsation]
La \cle{fréquence} $f$ d'un phénomène périodique est le nombre de périodes accomplies par seconde :
\[ f = \dfrac{1}{T} \qquad \text{avec } f \text{ en hertz (Hz) et } T \text{ en secondes (s).} \]
La \cle{pulsation} $\omega$ (lettre grecque oméga) est définie par :
\[ \omega = 2\pi f = \dfrac{2\pi}{T} \qquad \text{en radians par seconde (rad/s).} \]
\end{definition}

\begin{attention}[Unités et analyse dimensionnelle]
Le hertz est un nom donné à l'inverse de la seconde : $1\ \text{Hz} = 1\ \text{s}^{-1}$. Le radian étant sans dimension, la pulsation s'exprime aussi en $\text{s}^{-1}$, mais on écrit rad/s pour rappeler qu'il s'agit d'une vitesse angulaire. Vérifie toujours : si $T = 20\ \text{ms}$, alors $f = 1/0{,}020 = 50\ \text{Hz}$, et non $0{,}05\ \text{Hz}$ (erreur classique due à l'oubli de la conversion ms $\to$ s).
\end{attention}

\begin{exemplebox}[La tension du secteur camerounais]
La tension délivrée par ENEO sur le réseau basse tension est sinusoïdale de fréquence $f = 50\ \text{Hz}$. On en déduit :
\begin{align*}
T &= \dfrac{1}{f} = \dfrac{1}{50} = 0{,}020\ \text{s} = 20\ \text{ms} \\
\omega &= 2\pi f = 2\pi \times 50 \approx 314\ \text{rad/s}.
\end{align*}
Concrètement, la tension aux bornes d'une prise change de sens 100 fois par seconde : c'est pourquoi une lampe à incandescence semble briller en continu (notre œil ne perçoit pas ces variations rapides).
\end{exemplebox}

\courssub{Grandeur sinusoïdale : écriture générale}

\textbf{Intuition.} Fixe une craie sur la pointe d'une lame vibrante et fais défiler une feuille à vitesse constante devant elle : la trace obtenue est une \emph{sinusoïde}. De même, un point $M$ en mouvement circulaire uniforme à la vitesse angulaire $\omega$ a pour projeté sur un axe un mouvement dont l'élongation varie en cosinus du temps. C'est le lien profond entre rotation uniforme et oscillation sinusoïdale, que la représentation de Fresnel (section 3) exploitera.

\begin{propriete}[Expression d'une grandeur sinusoïdale (admise)]
Toute grandeur physique sinusoïdale du temps (élongation, tension, intensité, pression acoustique...) s'écrit :
\[ y(t) = Y_m \cos(\omega t + \varphi) \]
où :
\begin{itemize}
\item $Y_m$ est l'\cle{amplitude} (ou valeur maximale) : $Y_m > 0$, exprimée dans l'unité de $y$ ;
\item $\omega$ est la pulsation en rad/s ;
\item $\varphi$ est la \cle{phase à l'origine} (en rad) : elle fixe l'état de la grandeur à $t = 0$ ;
\item la quantité $(\omega t + \varphi)$ est la \cle{phase à l'instant} $t$ (en rad).
\end{itemize}
\end{propriete}

\textbf{Vocabulaire.} Pour un oscillateur mécanique, $y(t)$ s'appelle l'\cle{élongation} : c'est l'écart algébrique du point vibrant par rapport à sa position d'équilibre. La grandeur $y$ oscille entre $-Y_m$ et $+Y_m$ puisque le cosinus est compris entre $-1$ et $1$.

\begin{propriete}[Valeur efficace (admise)]
La \cle{valeur efficace} d'une grandeur sinusoïdale est définie à partir de sa valeur moyenne quadratique ; on admet le résultat :
\[ Y_{\text{eff}} = \dfrac{Y_m}{\sqrt{2}}. \]
C'est elle que mesurent les voltmètres et ampèremètres en régime sinusoïdal. Ainsi, la tension « $220\ \text{V}$ » du secteur est une valeur efficace : son amplitude vaut $U_m = 220\sqrt{2} \approx 311\ \text{V}$.
\end{propriete}

\begin{popfigure}
\begin{center}
\begin{tikzpicture}
\begin{axis}[
  width=0.92\linewidth, height=6cm,
  axis lines=middle,
  xlabel={$t$ (ms)}, ylabel={$y$},
  xmin=-2, xmax=45, ymin=-2.2, ymax=2.8,
  xtick={0,10,20,30,40}, ytick={-1.5,0,1.5},
  yticklabels={$-Y_m$,$0$,$Y_m$},
  tick label style={font=\small},
  samples=200
]
\addplot[popcurve,domain=0:42]{1.5*cos(360*x/20 - 60)};
% amplitude
\draw[<->,poporange,thick] (axis cs:1.667,0) -- (axis cs:1.667,1.5) node[midway,right,font=\small]{$Y_m$};
% période
\draw[<->,popblue,thick] (axis cs:1.667,1.9) -- (axis cs:21.667,1.9) node[midway,above,font=\small]{$T$};
\draw[dashed,popblue] (axis cs:1.667,1.5) -- (axis cs:1.667,2.1);
\draw[dashed,popblue] (axis cs:21.667,1.5) -- (axis cs:21.667,2.1);
% décalage dû à phi
\draw[<->,popgreen,thick] (axis cs:0,-1.9) -- (axis cs:3.333,-1.9) node[midway,below,font=\small]{$\Delta t = \dfrac{|\varphi|}{\omega}$};
\draw[dashed,popgreen] (axis cs:3.333,-1.5) -- (axis cs:3.333,-2.1);
\addplot[popmuted,dashed,domain=0:42]{1.5*cos(360*x/20)};
\end{axis}
\end{tikzpicture}
\end{center}
\legende{Grandeur sinusoïdale $y(t) = Y_m\cos(\omega t + \varphi)$ (trait plein) comparée à $Y_m\cos(\omega t)$ (pointillés). La phase à l'origine $\varphi$ se traduit par un décalage horaire $\Delta t = |\varphi|/\omega$ ; la période $T$ et l'amplitude $Y_m$ se lisent directement.}
\end{popfigure}

\begin{attention}[Signe de $\varphi$ et sens du décalage]
Si $\varphi > 0$, la courbe de $y(t) = Y_m\cos(\omega t + \varphi)$ est \emph{en avance} sur celle de $Y_m\cos(\omega t)$ : elle est décalée vers la gauche de $\Delta t = \varphi/\omega$. Si $\varphi < 0$, elle est en retard (décalée vers la droite). Retiens : « avance = plus tôt = vers la gauche ».
\end{attention}

\courssub{Déphasage entre deux grandeurs de même fréquence}

Considérons deux grandeurs sinusoïdales de \emph{même fréquence} (donc même pulsation $\omega$) :
\[ y_1(t) = Y_{1m}\cos(\omega t + \varphi_1) \qquad\text{et}\qquad y_2(t) = Y_{2m}\cos(\omega t + \varphi_2). \]

\begin{definition}[Déphasage]
Le \cle{déphasage} de $y_2$ par rapport à $y_1$ est la différence des phases à l'origine :
\[ \Delta\varphi = \varphi_2 - \varphi_1 \qquad \text{(en rad).} \]
\begin{itemize}
\item Si $\Delta\varphi > 0$ : $y_2$ est \cle{en avance de phase} sur $y_1$ ;
\item si $\Delta\varphi < 0$ : $y_2$ est \cle{en retard de phase} sur $y_1$ ;
\item si $\Delta\varphi = 0$ : les grandeurs sont \cle{en phase} ;
\item si $\Delta\varphi = \pi$ : elles sont \cle{en opposition de phase} ;
\item si $\Delta\varphi = \pm \pi/2$ : elles sont \cle{en quadrature de phase}.
\end{itemize}
\end{definition}

Le déphasage est lié au \cle{décalage horaire} $\Delta t$ entre deux passages homologues (par exemple deux maxima consécutifs des deux courbes) par la proportionnalité fondamentale :
\[ \boxed{\ \Delta\varphi = \dfrac{2\pi}{T}\,\Delta t = \omega\,\Delta t\ } \]
En effet, une période complète $T$ correspond à $2\pi$ rad ; par proportionnalité, un décalage $\Delta t$ correspond à $2\pi\Delta t/T$.

\begin{attention}[Piège majeur : $\Delta t$ n'est pas $\Delta\varphi$]
L'erreur la plus fréquente au Baccalauréat est de confondre le décalage horaire $\Delta t$, qui s'exprime en \emph{secondes} (et se lit sur l'axe des temps), avec le déphasage $\Delta\varphi$, qui s'exprime en \emph{radians}. Ils sont reliés par $\Delta\varphi = 2\pi\,\Delta t/T$. Exemple : si $T = 20\ \text{ms}$ et $\Delta t = 5\ \text{ms}$, alors $\Delta\varphi = 2\pi \times 5/20 = \pi/2\ \text{rad}$, et non $5\ \text{rad}$ !
\end{attention}

\courssub{Exploitation d'un oscillogramme}

L'oscilloscope est l'instrument roi pour étudier les grandeurs sinusoïdales. Sur son écran, deux réglages sont essentiels :
\begin{itemize}
\item le \cle{balayage} (ou base de temps) $b$, en ms/div : durée représentée par une division horizontale ;
\item la \cle{sensibilité verticale} $k$, en V/div : tension représentée par une division verticale.
\end{itemize}

\begin{popfigure}
\begin{center}
\begin{tikzpicture}[scale=0.62]
% écran
\draw[fill=popdark!8,rounded corners] (0,0) rectangle (12,8);
% grille
\draw[popline,step=1] (0.5,0.5) grid (11.5,7.5);
\draw[popmuted] (0.5,4) -- (11.5,4);
\draw[popmuted] (6,0.5) -- (6,7.5);
% voie 1
\draw[popblue,very thick,domain=0.5:11.5,samples=200] plot (\x,{4+2.2*cos(360*(\x-0.5)/11)});
% voie 2
\draw[poporange,very thick,domain=0.5:11.5,samples=200] plot (\x,{4+1.4*cos(360*(\x-0.5)/11 - 90)});
% double flèche décalage
\draw[<->,popgreen,very thick] (1.875,6.9) -- (4.625,6.9) node[midway,above,font=\small\bfseries]{$\Delta t$};
\draw[dashed,popgreen] (1.875,6.2) -- (1.875,7.1);
\draw[dashed,popgreen] (4.625,5.4) -- (4.625,7.1);
% amplitude voie 1
\draw[<->,popblue,thick] (9.6,4) -- (9.6,6.2) node[midway,right,font=\small]{$Y_{1m}$};
% légendes
\node[popblue,font=\small\bfseries] at (10.6,6.4) {voie 1};
\node[poporange,font=\small\bfseries] at (10.6,2.2) {voie 2};
\end{tikzpicture}
\end{center}
\legende{Écran d'oscilloscope bicourbe : deux tensions sinusoïdales de même fréquence, déphasées. La double flèche verte matérialise le décalage horaire $\Delta t$ entre deux maxima homologues ; la période $T$ se mesure entre deux maxima de la même voie.}
\end{popfigure}

\begin{methode}[Déterminer $T$, $Y_m$ et $\Delta\varphi$ sur un oscillogramme]
\begin{enumerate}
\item \textbf{Période} : compter le nombre $n$ de divisions horizontales entre deux maxima consécutifs d'une même voie, puis $T = n \times b$ (convertir en secondes).
\item \textbf{Fréquence et pulsation} : $f = 1/T$ puis $\omega = 2\pi/T$.
\item \textbf{Amplitude} : compter le nombre $m$ de divisions verticales entre l'axe et un maximum, puis $Y_m = m \times k$.
\item \textbf{Décalage horaire} : compter le nombre $d$ de divisions entre deux maxima homologues des deux voies, puis $\Delta t = d \times b$.
\item \textbf{Déphasage} : $\Delta\varphi = 2\pi\,\dfrac{\Delta t}{T}$ (en rad), avec le signe donné par l'observation : la voie dont le maximum apparaît \emph{en premier} (plus à gauche) est en avance.
\end{enumerate}
\end{methode}

\begin{exoresolu}[Oscillogramme d'un montage électrique]
Sur l'écran d'un oscilloscope, on observe deux tensions de même fréquence. Réglages : balayage $b = 5\ \text{ms/div}$, sensibilité voie 1 : $k_1 = 2\ \text{V/div}$, voie 2 : $k_2 = 1\ \text{V/div}$. La voie 1 culmine à 3,0 divisions au-dessus de l'axe ; la voie 2 culmine à 2,0 divisions. Une période occupe 4,0 divisions et le maximum de la voie 2 apparaît 1,0 division \emph{après} celui de la voie 1. Déterminer $T$, $f$, $\omega$, $U_{1m}$, $U_{2m}$, la valeur efficace $U_{1\text{eff}}$ et le déphasage $\Delta\varphi = \varphi_2 - \varphi_1$.
\tcblower
\textbf{Période, fréquence, pulsation.}
\begin{align*}
T &= 4{,}0 \times 5\ \text{ms} = 20\ \text{ms} = 2{,}0\times 10^{-2}\ \text{s} \\
f &= \dfrac{1}{T} = \dfrac{1}{2{,}0\times 10^{-2}} = 50\ \text{Hz} \\
\omega &= 2\pi f \approx 314\ \text{rad/s}.
\end{align*}
\textbf{Amplitudes.}
\[ U_{1m} = 3{,}0 \times 2 = 6{,}0\ \text{V} \qquad U_{2m} = 2{,}0 \times 1 = 2{,}0\ \text{V}. \]
\textbf{Valeur efficace de la voie 1.}
\[ U_{1\text{eff}} = \dfrac{U_{1m}}{\sqrt{2}} = \dfrac{6{,}0}{\sqrt{2}} \approx 4{,}2\ \text{V}. \]
\textbf{Déphasage.} Le décalage horaire vaut $\Delta t = 1{,}0 \times 5 = 5\ \text{ms}$. La voie 2 atteint son maximum \emph{après} la voie 1 : elle est en retard, donc $\Delta\varphi < 0$ :
\[ \Delta\varphi = -2\pi\,\dfrac{\Delta t}{T} = -2\pi\times\dfrac{5}{20} = -\dfrac{\pi}{2}\ \text{rad}. \]
Les deux tensions sont en quadrature de phase, la voie 2 étant en retard de $\pi/2$ sur la voie 1.
\end{exoresolu}

\begin{exercice}[À toi de jouer]
La tension d'un alternateur de groupe électrogène à Yaoundé s'écrit $u(t) = 311\cos(100\pi t)$ (en V). (1) Donner son amplitude, sa pulsation, sa fréquence, sa période et sa phase à l'origine. (2) Calculer sa valeur efficace. (3) Un courant $i(t) = 5\cos(100\pi t - \pi/3)$ (en A) traverse le même circuit : quel est le déphasage de $i$ par rapport à $u$ ? Lequel est en avance ?
\end{exercice}

\begin{corrige}[Exercice 1]
(1) Par identification avec $u(t) = U_m\cos(\omega t + \varphi)$ : $U_m = 311\ \text{V}$, $\omega = 100\pi \approx 314\ \text{rad/s}$, $\varphi = 0$. Donc $f = \omega/(2\pi) = 50\ \text{Hz}$ et $T = 1/f = 20\ \text{ms}$.
(2) $U_{\text{eff}} = U_m/\sqrt{2} = 311/\sqrt{2} \approx 220\ \text{V}$ : c'est bien la tension efficace du réseau.
(3) $\Delta\varphi = \varphi_i - \varphi_u = -\pi/3 - 0 = -\pi/3\ \text{rad}$. Le déphasage est négatif : le courant est \emph{en retard} de $\pi/3$ rad sur la tension, soit un décalage horaire $\Delta t = T/6 \approx 3{,}3\ \text{ms}$.
\end{corrige}

\begin{aretenir}[L'essentiel de la section]
\begin{itemize}
\item Phénomène périodique : $T$ (s), $f = 1/T$ (Hz), $\omega = 2\pi/T$ (rad/s).
\item Grandeur sinusoïdale : $y(t) = Y_m\cos(\omega t + \varphi)$ ; amplitude $Y_m$, phase $(\omega t + \varphi)$, phase à l'origine $\varphi$.
\item Valeur efficace (admise) : $Y_{\text{eff}} = Y_m/\sqrt{2}$.
\item Déphasage : $\Delta\varphi = \varphi_2 - \varphi_1 = 2\pi\,\Delta t/T$ ; signe $>$ 0 : avance, $<$ 0 : retard.
\item Sur un oscillogramme : $T = n\,b$, $Y_m = m\,k$, $\Delta t = d\,b$, puis $\Delta\varphi = 2\pi\Delta t/T$.
\end{itemize}
\end{aretenir}

\begin{savaistu}[Pourquoi 50 Hz au Cameroun ?]
La fréquence de 50 Hz adoptée au Cameroun (comme en Europe et en Afrique francophone) est un compromis historique : assez élevée pour éviter le scintillement visible des lampes, assez basse pour limiter les pertes dans les alternateurs des centrales comme Song Loulou ou Edéa. Les sismologues de l'IRGM, eux, travaillent avec des ondes de fréquences bien plus basses (de $0{,}1$ Hz à $20$ Hz) : la même boîte à outils mathématique — période, fréquence, phase — suffit à décrire ces mondes si différents.
\end{savaistu}

\coursec{Stroboscopie}

Dans la section précédente, nous avons appris à caractériser un phénomène périodique par son amplitude, sa période $T$, sa fréquence $f$ et sa phase. Mais une question pratique se pose immédiatement : comment \emph{mesurer} la fréquence d'un mouvement très rapide, par exemple la rotation de l'arbre d'un moteur de mototaxi qui peut atteindre plusieurs dizaines de tours par seconde ? L'œil humain est incapable de suivre un tel mouvement : au-delà d'une vingtaine d'images par seconde, tout se brouille. Il existe pourtant un instrument d'une simplicité géniale qui permet de « figer » le mouvement : le \cle{stroboscope}. Cette section lui est entièrement consacrée.

\courssub{Principe du stroboscope}

\begin{definition}[Stroboscope]
Un \cle{stroboscope} est une source lumineuse qui émet des \cle{éclairs} très brefs, séparés par des intervalles de temps égaux. La fréquence des éclairs, notée $f_e$, est réglable et connue avec précision. La durée d'un éclair est très inférieure à l'intervalle entre deux éclairs.
\end{definition}

Le principe repose sur une idée simple : entre deux éclairs, l'objet observé est dans le noir total, donc invisible. L'observateur ne voit l'objet qu'\emph{aux instants des éclairs}. Tout se passe comme si on prenait une photographie instantanée de l'objet à chaque éclair : le mouvement continu est remplacé par une succession de positions discrètes.

\begin{experience}[Observation d'un disque en rotation]
\textbf{Matériel :} un moteur électrique faisant tourner un disque blanc sur lequel on a tracé un rayon noir (ou collé une pastille) ; un stroboscope ; une salle obscurcie.

\textbf{Protocole :}
\begin{enumerate}
\item On met le disque en rotation à la fréquence $f$ (inconnue).
\item On éclaire le disque avec le stroboscope de fréquence $f_e$ réglable.
\item On fait varier $f_e$ et on observe l'aspect du rayon marqué.
\end{enumerate}

\textbf{Observations :} selon la valeur de $f_e$, le rayon paraît immobile (une ou plusieurs positions fixes), tourner lentement dans le sens réel, ou tourner lentement en sens inverse.
\end{experience}

\courssub{Immobilité apparente}

\begin{definition}[Immobilité apparente]
Il y a \cle{immobilité apparente} lorsque, éclairé par le stroboscope, l'objet en mouvement périodique paraît immobile : à chaque éclair, il occupe exactement la même position (ou un petit nombre de positions fixes).
\end{definition}

Cherchons la condition d'immobilité dans le cas du disque à un seul rayon marqué, tournant à la fréquence $f$ (période $T = 1/f$), éclairé à la fréquence $f_e$ (période des éclairs $T_e = 1/f_e$).

Entre deux éclairs consécutifs, il s'écoule la durée $T_e$. Pendant cette durée, le disque effectue un nombre de tours égal à :
\[ n = \frac{T_e}{T} = \frac{f}{f_e}. \]
Le rayon marqué réapparaît à la même position à chaque éclair si et seulement si $n$ est un \emph{entier} $k$ : le disque a alors fait exactement $k$ tours complets entre deux éclairs.

\begin{propriete}[Condition d'immobilité apparente]
Un objet animé d'un mouvement périodique de fréquence $f$, éclairé par un stroboscope de fréquence $f_e$, paraît immobile (en une seule position) si et seulement si :
\[ f = k\, f_e \qquad \text{soit} \qquad f_e = \frac{f}{k}, \quad k \in \mathbb{N}^{*}. \]
En particulier, pour $k = 1$ : $f_e = f$ (un tour exactement entre deux éclairs).
\end{propriete}

\begin{experience}[Démonstration guidée : pourquoi $f = k\,f_e$ ?]
Raisonnons pas à pas sur le disque à un rayon marqué.
\begin{enumerate}
\item À l'instant $t_0$ d'un premier éclair, le rayon est observé dans une position donnée, repérée par l'angle $\theta_0$.
\item Le disque tourne à la vitesse angulaire $\omega = 2\pi f$. À l'instant $t_0 + T_e$ de l'éclair suivant, le rayon a tourné de l'angle $\Delta\theta = \omega T_e = 2\pi f T_e = 2\pi \dfrac{f}{f_e}$.
\item Le rayon est revu \emph{dans la même position} si $\Delta\theta$ est un multiple entier de $2\pi$ : $2\pi \dfrac{f}{f_e} = 2\pi k$, soit $\boxed{f = k\,f_e}$.
\item Vérification dimensionnelle : $f$ et $f_e$ sont toutes deux en hertz (s$^{-1}$), $k$ est sans dimension. La relation est homogène.
\end{enumerate}
\textbf{Cas particulier $f_e = f/2$ ($k=2$) :} entre deux éclairs, le disque fait deux tours ; le rayon est revu à la même place. L'immobilité est donc \emph{aussi} observée. C'est le piège des sous-multiples, étudié plus loin.
\end{experience}

\begin{attention}[Le piège des sous-multiples]
L'immobilité apparente est observée pour $f_e = f$, mais \emph{aussi} pour $f_e = f/2$, $f_e = f/3$, etc. Conclure « $f = f_e$ » dès la première immobilité observée est une erreur classique au Baccalauréat : si l'immobilité apparaît pour $f_e = \SI{25}{Hz}$, la fréquence réelle peut valoir $25$, $50$, $75$~Hz\dots{} Seule la \textbf{plus grande} fréquence d'éclairs donnant l'immobilité est égale à $f$.
\end{attention}

\begin{methode}[Mesurer une fréquence sans ambiguïté]
\begin{enumerate}
\item Régler le stroboscope à sa fréquence maximale, puis \textbf{diminuer progressivement} $f_e$.
\item Repérer la \textbf{première} (donc la plus grande) valeur $f_{e1}$ pour laquelle l'objet paraît immobile en une seule position.
\item Conclure : $f = f_{e1}$.
\item Vérification : en continuant à diminuer, on retrouve l'immobilité pour $f_{e1}/2$, $f_{e1}/3$\dots{} ce qui confirme le résultat.
\end{enumerate}
\end{methode}

\courssub{Mouvement ralenti apparent}

Que se passe-t-il lorsque $f_e$ est \emph{voisine} de $f$, sans lui être égale ?

\begin{definition}[Mouvement ralenti apparent]
Lorsque $f_e$ est légèrement différente de $f$ (ou d'un sous-multiple $f/k$), l'objet paraît animé d'un mouvement lent : c'est le \cle{mouvement ralenti apparent}.
\begin{itemize}
\item Si $f_e \lesssim f$ (éclairs légèrement moins fréquents) : entre deux éclairs, l'objet fait un peu plus d'un tour ; il est revu légèrement \emph{en avance}. Le ralenti apparent se fait \textbf{dans le sens réel} du mouvement.
\item Si $f_e \gtrsim f$ : l'objet fait un peu moins d'un tour ; il est revu légèrement \emph{en arrière}. Le ralenti apparent se fait \textbf{en sens inverse} du mouvement réel.
\end{itemize}
\end{definition}

\begin{propriete}[Fréquence apparente]
Pour $f_e$ voisine de $f/k$, la fréquence du mouvement apparent est :
\[ f_a = \left| f - k\, f_e \right|. \]
Le sens apparent dépend du signe de $f - k f_e$ : positif $\Rightarrow$ sens réel ; négatif $\Rightarrow$ sens inverse.
\end{propriete}

Justifions cette formule sur le disque. Entre deux éclairs, le rayon tourne de $\Delta\theta = 2\pi f T_e$. Par rapport à la position précédente, le décalage angulaire apparent (ramené entre $-\pi$ et $\pi$) est $\delta = 2\pi(f T_e - k) = 2\pi\,\dfrac{f - k f_e}{f_e}$. Ce décalage se reproduit à chaque éclair, c'est-à-dire $f_e$ fois par seconde : la vitesse angulaire apparente est $\omega_a = 2\pi |f - k f_e|$, d'où $f_a = |f - k f_e|$.

\begin{exemplebox}[Ralenti apparent d'un ventilateur]
Un ventilateur tourne à $f = \SI{20,0}{Hz}$. On l'éclaire avec un stroboscope de fréquence $f_e = \SI{19,5}{Hz}$ ($k = 1$).
\[ f_a = |f - f_e| = |20,0 - 19,5| = \SI{0,5}{Hz}. \]
Le ventilateur paraît tourner lentement à $0,5$ tour par seconde, dans le sens réel (car $f_e < f$). Avec $f_e = \SI{20,5}{Hz}$, on obtiendrait la même fréquence apparente $\SI{0,5}{Hz}$ mais en sens inverse.
\end{exemplebox}

\begin{popfigure}
\begin{center}
\begin{tikzpicture}[scale=1.05]
% --- Disque 1 : fe = f ---
\begin{scope}[shift={(0,0)}]
\draw[thick,popblue,fill=popblueL] (0,0) circle (1.3);
\fill[popdark] (0,0) circle (0.05);
\foreach \a in {30} {\draw[line width=2pt,poporange] (0,0) -- (\a:1.25); \fill[poporange] (\a:1.1) circle (0.09);}
\draw[dashed,popmuted] (0,0) -- (90:1.25);
\node[below,align=center] at (0,-1.5) {\small $f_e = f$\\ \small immobilité};
\end{scope}
% --- Disque 2 : fe légèrement < f ---
\begin{scope}[shift={(4.6,0)}]
\draw[thick,popblue,fill=popblueL] (0,0) circle (1.3);
\fill[popdark] (0,0) circle (0.05);
\foreach \a in {30,55,80} {\draw[line width=1.4pt,poporange!80] (0,0) -- (\a:1.25); \fill[poporange!80] (\a:1.1) circle (0.08);}
\draw[-{Stealth},thick,poppurple] (30:0.7) arc (30:80:0.7);
\node[below,align=center] at (0,-1.5) {\small $f_e \lesssim f$\\ \small ralenti sens réel};
\end{scope}
% --- Disque 3 : fe = f/2 ---
\begin{scope}[shift={(9.2,0)}]
\draw[thick,popblue,fill=popblueL] (0,0) circle (1.3);
\fill[popdark] (0,0) circle (0.05);
\draw[line width=2pt,popgreen] (0,0) -- (30:1.25); \fill[popgreen] (30:1.1) circle (0.09);
\node[below,align=center] at (0,-1.5) {\small $f_e = f/2$\\ \small immobilité aussi !};
\end{scope}
% Stroboscope
\draw[fill=popgoldL,draw=popgold,thick] (-1.6,2.2) rectangle (-0.6,2.8);
\node at (-1.1,2.5) {\small S};
\foreach \x in {-0.5,-0.1,0.3} {\draw[-{Stealth},popgold] (\x,2.5) -- (\x+0.5,1.6);}
\node[popgold,right] at (0.4,2.2) {\small éclairs};
\end{tikzpicture}
\end{center}
\legende{Disque à un rayon marqué éclairé par un stroboscope (S). À gauche : $f_e = f$, le rayon est toujours revu à la même place. Au centre : $f_e$ légèrement inférieure à $f$, les positions successives (éclairs 1, 2, 3) avancent lentement dans le sens réel (flèche violette). À droite : $f_e = f/2$, le disque fait deux tours entre deux éclairs et paraît également immobile.}
\end{popfigure}

\courssub{Applications de la stroboscopie}

\begin{itemize}
\item \textbf{Mesure de fréquence de rotation} des moteurs (ateliers mécaniques, maintenance des groupes électrogènes, si précieux lors des délestages !) sans contact avec la machine.
\item \textbf{Réglage et contrôle des machines} : le ralenti apparent permet d'observer le fonctionnement d'un mécanisme rapide (vibration d'une pièce, défaut d'alignement d'un arbre).
\item \textbf{Cinéma et vidéo} : la caméra échantillonne le mouvement à $25$ images par seconde, comme un stroboscope : d'où les illusions de roues ou d'hélices tournant à l'envers.
\end{itemize}

\begin{exoresolu}[Fréquence de rotation d'un moteur]
\textbf{Énoncé.} Pour mesurer la fréquence de rotation $f$ de l'arbre d'un moteur, un technicien y fixe un disque portant un seul repère et l'éclaire avec un stroboscope. En diminuant $f_e$ depuis sa valeur maximale, la première immobilité (repère unique fixe) apparaît pour $f_{e1} = \SI{50}{Hz}$ ; l'immobilité réapparaît pour $f_{e2} = \SI{25}{Hz}$. (1) Déterminer $f$. (2) Que vaut $f_{e2}$ par rapport à $f$ ? (3) On règle ensuite $f_e = \SI{49}{Hz}$ : décrire le mouvement apparent.
\tcblower
\textbf{Solution.}
\begin{enumerate}
\item La plus grande fréquence d'éclairs donnant l'immobilité est $f_{e1} = \SI{50}{Hz}$. D'après la condition $f = k f_e$ avec $k = 1$ : $\boxed{f = \SI{50}{Hz}}$, soit $50 \times 60 = 3000~\text{tr/min}$.
\item $f_{e2} = 25 = f/2$ : c'est le sous-multiple d'ordre $k = 2$. Entre deux éclairs, le disque effectue deux tours complets ; le repère est revu à la même place. Cela confirme la mesure.
\item $f_e = \SI{49}{Hz} \lesssim f$ : mouvement ralenti apparent dans le sens réel, de fréquence
\[ f_a = |f - f_e| = |50 - 49| = \SI{1}{Hz}. \]
Le repère paraît faire un tour par seconde dans le sens réel de rotation.
\end{enumerate}
\end{exoresolu}

\begin{attention}[Rédaction au Baccalauréat]
Lorsqu'un énoncé donne plusieurs fréquences d'immobilité ($f_{e1} > f_{e2} > \dots$), la fréquence du mouvement est la \textbf{plus grande} : $f = f_{e1}$. Et on doit vérifier la cohérence : $f_{e1} = 2 f_{e2} = 3 f_{e3}$\dots{} Préciser toujours le sens (réel ou inverse) lorsqu'on décrit un ralenti apparent.
\end{attention}

\begin{savaistu}[Les hélices qui tournent à l'envers]
Sur les vidéos filmées à $25$ images par seconde (caméras, téléphones), les hélices d'avion ou les roues de voiture semblent parfois tourner lentement, voire à l'envers. C'est exactement l'effet stroboscopique : l'obturateur joue le rôle du stroboscope avec $f_e = \SI{25}{Hz}$. Si une pale d'hélice tourne à $f = \SI{26}{Hz}$, la fréquence apparente vaut $f_a = |26 - 25| = \SI{1}{Hz}$ dans le sens réel ; si $f = \SI{24}{Hz}$, l'hélice semble reculer à $\SI{1}{Hz}$ ! Le même phénomène explique les roues de charrettes « immobiles » dans les westerns.
\end{savaistu}

\begin{aretenir}[L'essentiel de la stroboscopie]
\begin{itemize}
\item Le stroboscope émet des éclairs brefs de fréquence réglable $f_e$ ; on ne voit l'objet qu'aux instants des éclairs.
\item \textbf{Immobilité apparente} : $f = k f_e$ ($k$ entier). La fréquence cherchée est la \textbf{plus grande} valeur de $f_e$ donnant l'immobilité.
\item \textbf{Ralenti apparent} : $f_a = |f - k f_e|$ ; sens réel si $f_e \lesssim f/k$, sens inverse si $f_e \gtrsim f/k$.
\item Attention aux sous-multiples : $f_e = f/2, f/3, \dots$ donnent aussi l'immobilité.
\end{itemize}
\end{aretenir}

La stroboscopie nous a donné un premier outil d'analyse des mouvements périodiques. Pour aller plus loin — notamment pour \emph{additionner} des grandeurs sinusoïdales de même fréquence, comme on le fera en électricité et en optique —, il nous faut un outil géométrique puissant : la représentation de Fresnel, objet de la section suivante.

\coursec{Représentation de Fresnel}

La stroboscopie nous a appris à « photographier » un mouvement périodique : en choisissant bien la fréquence des éclairs, on fige le mouvement ou on le ralentit. Mais figer un instant ne suffit pas pour \emph{calculer}. En électricité comme en mécanique, on doit souvent \textbf{additionner} des grandeurs sinusoïdales de même fréquence : deux tensions alternatives aux bornes de dipôles en série, deux vibrations qui se superposent sur une corde, deux ondes sonores qui arrivent à ton oreille. Additionner directement des cosinus déphasés est fastidieux. Le physicien français Augustin Fresnel a eu une idée géniale : remplacer chaque grandeur sinusoïdale par un \textbf{vecteur tournant}, ce qui transforme les calculs trigonométriques en simple géométrie. C'est l'objet de cette section.

\courssub{Principle : du vecteur tournant à la sinusoïde}

\textbf{Une observation géométrique}\par

Imagine une tige rigide de longueur $Y_m$, fixée en un point $O$, qui tourne dans un plan à vitesse angulaire constante $\omega$ (sens trigonométrique). À l'instant $t$, elle fait l'angle $(\omega t + \varphi)$ avec un axe horizontal $Ox$. Éclaire cette tige par le haut : son ombre sur l'axe $Ox$ oscille de $-Y_m$ à $+Y_m$. Or l'abscisse de l'extrémité $M$ de la tige vaut exactement :
\[ y(t) = Y_m \cos(\omega t + \varphi). \]
L'ombre du vecteur tournant décrit \emph{précisément} une grandeur sinusoïdale ! Réciproquement, toute grandeur sinusoïdale peut être vue comme la projection d'un vecteur tournant. C'est toute la représentation de Fresnel.

\begin{definition}[Vecteur de Fresnel associé à une grandeur sinusoïdale]
À toute grandeur sinusoïdale $y(t) = Y_m \cos(\omega t + \varphi)$, on associe un vecteur $\overrightarrow{OM}$, appelé \cle{vecteur de Fresnel}, tel que :
\begin{itemize}
\item sa \textbf{norme} est égale à l'amplitude : $\|\overrightarrow{OM}\| = Y_m$ ;
\item l'\textbf{angle} qu'il fait avec l'axe de référence $Ox$ est la phase instantanée : $(\overrightarrow{Ox}, \overrightarrow{OM}) = \omega t + \varphi$ ;
\item il \textbf{tourne} autour de $O$ dans le sens trigonométrique à la vitesse angulaire $\omega = \dfrac{2\pi}{T} = 2\pi f$.
\end{itemize}
La grandeur $y(t)$ est alors la \textbf{projection orthogonale} de $\overrightarrow{OM}$ sur l'axe $Ox$ :
\[ y(t) = \|\overrightarrow{OM}\| \cos(\omega t + \varphi). \]
\end{definition}

\begin{aretenir}[Construction de Fresnel à $t = 0$]
Comme tous les vecteurs de Fresnel d'un même problème tournent à la \emph{même} vitesse $\omega$, leurs positions relatives ne changent jamais. On les représente donc figés à l'instant $t = 0$ : le vecteur $\overrightarrow{OM}$ a pour norme $Y_m$ et fait l'angle $\varphi$ (phase à l'origine) avec $Ox$. Toute l'information de la grandeur sinusoïdale — amplitude et phase — est contenue dans ce seul vecteur.
\end{aretenir}

\begin{popfigure}
\begin{center}
\begin{tikzpicture}[scale=1.0]
% --- Cercle de Fresnel (gauche) ---
\begin{scope}[shift={(0,0)}]
\draw[popline] (0,0) circle (1.6);
\draw[->,popmuted] (-2.1,0) -- (2.3,0) node[below left] {$Ox$};
\draw[->,popmuted] (0,-2.1) -- (0,2.1);
% vecteur tournant
\draw[->,very thick,popblue] (0,0) -- (50:1.6) node[above right] {$M$};
\draw[poporange,thick,->] (0.7,0) arc (0:50:0.7) node[midway,right] {\scriptsize $\omega t+\varphi$};
% projection
\coordinate (P) at ({1.6*cos(50)},0);
\draw[dashed,popink] (50:1.6) -- (P);
\fill[popink] (P) circle (1.3pt) node[below] {\scriptsize $y(t)$};
\fill[popblue] (50:1.6) circle (1.5pt);
\draw[->,popgreen,thick] (120:1.9) arc (120:80:1.9) node[midway,above] {\scriptsize $\omega$};
\node[popmuted] at (0,-2.5) {\scriptsize Vecteur de Fresnel tournant};
\end{scope}
% --- Sinusoïde déroulée (droite) ---
\begin{scope}[shift={(6.2,0)}]
\draw[->,popmuted] (-0.3,0) -- (5.6,0) node[below left] {$t$};
\draw[->,popmuted] (0,-2.1) -- (0,2.1) node[left] {$y$};
\draw[popblue,thick,domain=0:5.2,samples=120] plot (\x,{1.6*cos(50 - 57.3*1.05*\x)});
% correspondance point courant
\fill[popink] (0,{1.6*cos(50)}) circle (1.5pt);
\draw[dashed,popink] (0,{1.6*cos(50)}) -- (-3.4,{1.6*cos(50)});
\draw[poporange] (0,1.6) node[above right] {\scriptsize $Y_m$};
\draw[<->,poporange] (5.0,0) -- (5.0,1.6);
\draw[<->,popgreen] (0.955,1.9) -- (1.91,1.9) node[midway,above] {\scriptsize $T$};
\node[popmuted] at (2.6,-2.5) {\scriptsize Sinusoïde déroulée $y(t)$};
\end{scope}
\end{tikzpicture}
\legende{Correspondance entre le vecteur de Fresnel tournant (à gauche) et la sinusoïde $y(t) = Y_m\cos(\omega t+\varphi)$ (à droite) : $y(t)$ est la projection de $\overrightarrow{OM}$ sur $Ox$.}
\end{center}
\end{popfigure}

\begin{exemplebox}[Lecture directe d'un vecteur de Fresnel]
Une tension alternative est représentée à $t=0$ par un vecteur de norme $5$ unités faisant l'angle $-\dfrac{\pi}{6}$ avec $Ox$, et tournant à $\omega = \SI{100\pi}{rad.s^{-1}}$ (fréquence du réseau : $f = \SI{50}{Hz}$, comme chez ENEO). Son équation horaire s'écrit immédiatement :
\[ u(t) = 5\cos\left(100\pi t - \frac{\pi}{6}\right) \quad (\text{en volts si l'échelle est } 1\ \text{unité} = \SI{1}{V}). \]
\end{exemplebox}

\courssub{Somme de deux grandeurs sinusoïdales de même pulsation}

\textbf{Le problème posé}\par

Soient deux grandeurs de \textbf{même pulsation} $\omega$ :
\[ y_1(t) = Y_{1m}\cos(\omega t + \varphi_1) \qquad \text{et} \qquad y_2(t) = Y_{2m}\cos(\omega t + \varphi_2). \]
Que vaut $y = y_1 + y_2$ ? L'addition directe des cosinus est inextricable. Mais observons les vecteurs de Fresnel $\overrightarrow{OM_1}$ et $\overrightarrow{OM_2}$ : tournant tous deux à la vitesse $\omega$, l'angle entre eux, $\varphi_2 - \varphi_1$, reste \textbf{constant}. Leur somme vectorielle $\overrightarrow{OM} = \overrightarrow{OM_1} + \overrightarrow{OM_2}$ tourne donc elle aussi à la vitesse $\omega$, sans se déformer. C'est le cœur du théorème suivant.

\begin{propriete}[Théorème de la somme — démonstration exigible]
La somme de deux grandeurs sinusoïdales de même pulsation $\omega$ est une grandeur sinusoïdale de même pulsation $\omega$. Son vecteur de Fresnel est la somme vectorielle des vecteurs de Fresnel :
\[ \overrightarrow{OM} = \overrightarrow{OM_1} + \overrightarrow{OM_2}. \]
\end{propriete}

\begin{experience}[Démonstration guidée du théorème]
\begin{enumerate}
\item Par définition, $y_1(t)$ est la projection de $\overrightarrow{OM_1}$ sur $Ox$, et $y_2(t)$ celle de $\overrightarrow{OM_2}$.
\item La projection est une opération \textbf{linéaire} : la projection d'une somme de vecteurs est la somme des projections. Donc la projection de $\overrightarrow{OM} = \overrightarrow{OM_1} + \overrightarrow{OM_2}$ sur $Ox$ vaut $y_1(t) + y_2(t) = y(t)$.
\item Les deux vecteurs tournent à la même vitesse angulaire $\omega$ : le parallélogramme $OM_1MM_2$ tourne \emph{en bloc}, sans se déformer. Le vecteur $\overrightarrow{OM}$ a donc une norme constante $Y_m$ et tourne à la vitesse $\omega$.
\item Sa projection sur $Ox$ est donc une grandeur sinusoïdale de pulsation $\omega$, d'amplitude $Y_m = \|\overrightarrow{OM}\|$ et de phase à l'origine $\varphi = (\overrightarrow{Ox}, \overrightarrow{OM})$. \hfill $\blacksquare$
\end{enumerate}
\end{experience}

\textbf{Détermination de l'amplitude et de la phase résultantes}\par

À $t = 0$, on trace $\overrightarrow{OM_1}$ (norme $Y_{1m}$, angle $\varphi_1$) et $\overrightarrow{OM_2}$ (norme $Y_{2m}$, angle $\varphi_2$), puis on complète le \textbf{parallélogramme} : la diagonale $\overrightarrow{OM}$ donne $Y_m$ (sa longueur) et $\varphi$ (son angle avec $Ox$).

\begin{popfigure}
\begin{center}
\begin{tikzpicture}[scale=1.1]
\draw[->,popmuted] (-0.5,0) -- (5.2,0) node[below left] {$Ox$};
\draw[->,popmuted] (0,-0.8) -- (0,3.6);
% OM1 : norme 3, angle 0°
\draw[->,very thick,popblue] (0,0) -- (3,0) node[midway,below] {$\overrightarrow{OM_1}$};
% OM2 : norme 2, angle 60°
\draw[->,very thick,poporange] (0,0) -- (60:2) node[midway,above left] {$\overrightarrow{OM_2}$};
% parallélogramme
\coordinate (M) at ($(3,0)+(60:2)$);
\draw[dashed,popmuted] (3,0) -- (M);
\draw[dashed,popmuted] (60:2) -- (M);
% somme
\draw[->,very thick,popink] (0,0) -- (M) node[midway,above right] {$\overrightarrow{OM}$};
\fill[popink] (M) circle (1.6pt) node[above] {$M$};
\fill[popblue] (3,0) circle (1.4pt) node[below right] {$M_1$};
\fill[poporange] (60:2) circle (1.4pt) node[above left] {$M_2$};
% angles
\draw[popblue,->] (1.1,0) arc (0:60:1.1) node[midway,right] {\scriptsize $\varphi_2 - \varphi_1 = \dfrac{\pi}{3}$};
\draw[popink,->] (2.1,0) arc (0:30:2.1) node[midway,right] {\scriptsize $\varphi$};
\end{tikzpicture}
\legende{Construction du parallélogramme : la somme $y = y_1 + y_2$ a pour vecteur de Fresnel la diagonale $\overrightarrow{OM}$. Ici $Y_{1m} = 3$, $Y_{2m} = 2$, déphasage $\pi/3$.}
\end{center}
\end{popfigure}

\begin{methode}[Construire la somme de deux grandeurs sinusoïdales de même pulsation]
\begin{enumerate}
\item \textbf{Identifier} les amplitudes $Y_{1m}$, $Y_{2m}$ et les phases à l'origine $\varphi_1$, $\varphi_2$. Vérifier que les deux grandeurs ont la \emph{même} pulsation $\omega$ (sinon la méthode ne s'applique pas !).
\item \textbf{Choisir une échelle} (par exemple $1\ \text{cm} \leftrightarrow 1$ unité) et tracer $\overrightarrow{OM_1}$ et $\overrightarrow{OM_2}$ à $t = 0$, au rapporteur.
\item \textbf{Compléter le parallélogramme} $OM_1MM_2$ (ou reporter $\overrightarrow{OM_2}$ au bout de $\overrightarrow{OM_1}$).
\item \textbf{Mesurer} la diagonale : sa longueur donne $Y_m$ (avec l'échelle), son angle avec $Ox$ donne $\varphi$.
\item \textbf{Conclure} : $y(t) = Y_m \cos(\omega t + \varphi)$.
\item \emph{Vérification par le calcul} (loi des cosinus) : $Y_m = \sqrt{Y_{1m}^2 + Y_{2m}^2 + 2\,Y_{1m}Y_{2m}\cos(\varphi_2 - \varphi_1)}$.
\end{enumerate}
\end{methode}

\courssub{Cas particuliers remarquables}

\begin{aretenir}[Trois cas à connaître par cœur]
Soit $\Delta\varphi = \varphi_2 - \varphi_1$ le déphasage entre les deux grandeurs.
\begin{itemize}
\item \textbf{Grandeurs en phase} ($\Delta\varphi = 0$) : les vecteurs sont colinéaires de même sens, les amplitudes \textbf{s'ajoutent} : $Y_m = Y_{1m} + Y_{2m}$.
\item \textbf{Grandeurs en opposition de phase} ($\Delta\varphi = \pi$) : les vecteurs sont colinéaires de sens contraires, les amplitudes \textbf{se retranchent} : $Y_m = |Y_{1m} - Y_{2m}|$. Si $Y_{1m} = Y_{2m}$, la somme est nulle : c'est le principe des casques audio « à réduction active de bruit ».
\item \textbf{Grandeurs en quadrature} ($\Delta\varphi = \pm \dfrac{\pi}{2}$) : les vecteurs sont perpendiculaires ; le parallélogramme est un rectangle et Pythagore donne $Y_m = \sqrt{Y_{1m}^2 + Y_{2m}^2}$.
\end{itemize}
\end{aretenir}

\begin{attention}[Piège classique : additionner les amplitudes de grandeurs déphasées]
Deux grandeurs d'amplitudes $3$ et $4$ n'ont une somme d'amplitude $7$ \textbf{que si elles sont en phase}. En général, $Y_m < Y_{1m} + Y_{2m}$ : c'est l'inégalité triangulaire appliquée au parallélogramme. Avant tout calcul, demande-toi toujours : \emph{quel est le déphasage ?} Rédiger « $Y_m = Y_{1m} + Y_{2m}$ » sans justifier que les grandeurs sont en phase est une erreur éliminatoire au Bac.
\end{attention}

\begin{exoresolu}[Somme de deux grandeurs en quadrature]
\textbf{Énoncé.} Deux vibrations sinusoïdales de même pulsation $\omega$ s'exercent en un point d'une membrane : $y_1(t) = \SI{3}{\milli\meter}\cos(\omega t)$ et $y_2(t) = \SI{4}{\milli\meter}\cos\left(\omega t + \dfrac{\pi}{2}\right)$. Déterminer l'amplitude $Y_m$ et la phase à l'origine $\varphi$ de la vibration résultante $y = y_1 + y_2$.
\tcblower
\textbf{Solution.} Les deux grandeurs ont la même pulsation $\omega$ : le théorème de Fresnel s'applique. À $t = 0$, $\overrightarrow{OM_1}$ a pour norme $\SI{3}{\milli\meter}$ et angle $0$ ; $\overrightarrow{OM_2}$ a pour norme $\SI{4}{\milli\meter}$ et angle $\dfrac{\pi}{2}$. Les deux vecteurs sont \textbf{perpendiculaires} (quadrature) : le triangle $OM_1M$ est rectangle en $M_1$. Le théorème de Pythagore donne :
\[ Y_m = \sqrt{Y_{1m}^2 + Y_{2m}^2} = \sqrt{3^2 + 4^2} = \sqrt{25} = \SI{5}{\milli\meter}. \]
La phase à l'origine est l'angle de la diagonale avec $Ox$ :
\[ \tan\varphi = \frac{Y_{2m}}{Y_{1m}} = \frac{4}{3} \quad \Rightarrow \quad \varphi = \arctan\left(\frac{4}{3}\right) \approx \num{0,93}\ \text{rad} \ (\approx 53^\circ). \]
D'où la vibration résultante :
\[ y(t) = \SI{5}{\milli\meter}\cos(\omega t + \num{0,93}). \]
\emph{Remarque} : on reconnaît le célèbre triangle $3$--$4$--$5$. On vérifie que $Y_m = \SI{5}{\milli\meter} < \SI{7}{\milli\meter} = 3 + 4$ : les amplitudes ne s'ajoutent pas car les grandeurs sont déphasées.
\end{exoresolu}

\courssub{Complément : le lien avec les nombres complexes (hors programme strict)}

\begin{savaistu}[Pour les élèves rapides : la notation complexe]
Au vecteur de Fresnel $\overrightarrow{OM}$ (norme $Y_m$, angle $\omega t + \varphi$) correspond le nombre complexe $\underline{Y} = Y_m\,e^{i(\omega t + \varphi)}$, et $y(t)$ n'est autre que sa partie réelle. Additionner des grandeurs sinusoïdales revient alors à additionner des nombres complexes — ce que tu sais faire en mathématiques ! C'est cette « notation complexe » que les ingénieurs utilisent chaque jour pour dimensionner les réseaux électriques alternatifs, comme ceux qu'ENEO et la SONATREL déploient entre les barrages de Lom Pangar, Nachtigal et les villes du Cameroun. Elle n'est pas exigible au Bac, mais elle te donne une longueur d'avance pour l'enseignement supérieur.
\end{savaistu}

\begin{exercice}[Pour s'entraîner]
Deux tensions sinusoïdales de pulsation $\omega = \SI{100\pi}{rad.s^{-1}}$ sont associées en série : $u_1(t) = \SI{6}{V}\cos(100\pi t)$ et $u_2(t) = \SI{6}{V}\cos\left(100\pi t + \dfrac{\pi}{3}\right)$. Par la construction de Fresnel (puis par la loi des cosinus), déterminer l'amplitude $U_m$ et la phase à l'origine $\varphi$ de la tension totale $u = u_1 + u_2$. Que vaudrait $U_m$ si les deux tensions étaient en phase ? En opposition de phase ?
\end{exercice}

\begin{corrige}[Exercice : somme de deux tensions]
Le déphasage vaut $\Delta\varphi = \dfrac{\pi}{3}$. La loi des cosinus donne :
\[ U_m = \sqrt{6^2 + 6^2 + 2\times 6\times 6\times\cos\frac{\pi}{3}} = \sqrt{36 + 36 + 36} = \sqrt{108} = 6\sqrt{3} \approx \SI{10,4}{V}. \]
Par symétrie (losange), la diagonale bissecte l'angle : $\varphi = \dfrac{\pi}{6}$ rad. D'où $u(t) = \SI{6\sqrt{3}}{V}\cos\left(100\pi t + \dfrac{\pi}{6}\right)$. En phase, on aurait $U_m = \SI{12}{V}$ ; en opposition de phase, $U_m = \SI{0}{V}$.
\end{corrige}

En résumé, la représentation de Fresnel transforme l'algèbre des cosinus en géométrie des vecteurs : un outil puissant que nous réutiliserons dès que deux vibrations ou deux ondes se superposeront — précisément la situation que nous étudierons avec les ondes progressives sinusoïdales.

\coursec{Ondes mécaniques progressives}

Dans la section précédente, nous avons appris à représenter une grandeur sinusoïdale par un vecteur de Fresnel : un outil puissant pour décrire un \emph{seul} oscillateur, vibrant en un \emph{seul} point. Mais que se passe-t-il lorsque la vibration ne reste pas confinée à sa source ? Lorsqu'un enfant jette un caillou dans le lac Municipal de Yaoundé, des rides circulaires s'éloignent du point d'impact ; lorsqu'un guitariste pince une corde, la vibration parcourt la corde de part en part ; lorsque le tonnerre gronde sur les collines de l'Ouest, c'est une perturbation de l'air qui a voyagé jusqu'à ton oreille. Nous quittons donc l'étude des systèmes oscillants isolés pour celle de la \cle{propagation} : comment une perturbation se transmet-elle de proche en proche dans un milieu matériel ? C'est l'objet de cette section.

\courssub{Qu'est-ce qu'une onde mécanique ?}

\courssub{Mise en évidence expérimentale}

\begin{experience}[Trois expériences de référence]
\textbf{Expérience 1 — La corde.} On tend horizontalement une corde de quelques mètres, fixée à une extrémité. On imprime à l'autre extrémité un mouvement brusque vers le haut puis vers le bas (une \emph{impulsion}). On observe une déformation qui se déplace le long de la corde à vitesse constante, tandis que chaque point de la corde ne fait que s'élever puis redescendre, sans avancer.

\textbf{Expérience 2 — La cuve à ondes.} Une pointe vibre à la surface de l'eau contenue dans une cuve peu profonde. Des rides circulaires naissent au contact de la pointe et s'élargissent. Un bouchon de liège flottant à la surface monte et descend, mais ne s'éloigne pas de la source.

\textbf{Expérience 3 — Le ressort.} On dispose un long ressort à spires souples sur une table lisse. On comprime brusquement quelques spires à une extrémité, puis on relâche : une zone de compression se propage le long du ressort, chaque spire oscillant autour de sa position d'équilibre dans la direction du ressort.

\textbf{Interprétation commune.} Dans les trois cas, une \emph{perturbation} créée en un point se transmet de proche en proche grâce aux interactions entre les éléments voisins du milieu (tension entre brins de la corde, forces de pression et de surface pour l'eau, forces élastiques entre spires). Aucun élément de matière n'accompagne la perturbation dans son déplacement : seule la déformation — et l'énergie qu'elle porte — voyage.
\end{experience}

\begin{definition}[Onde mécanique progressive]
On appelle \cle{onde mécanique progressive} le phénomène de propagation d'une perturbation dans un milieu matériel, de proche en proche, \textbf{sans transport de matière} mais \textbf{avec transport d'énergie}.

Le milieu matériel (corde, eau, air, ressort, roche...) est dit \emph{milieu de propagation}. Une onde mécanique ne peut pas se propager dans le vide : c'est ce qui la distingue des ondes électromagnétiques (lumière, ondes radio), que tu étudieras plus tard.
\end{definition}

\begin{attention}[La matière ne voyage pas !]
L'erreur la plus fréquente consiste à croire que la matière se déplace avec l'onde. C'est faux : le bouchon de liège sur la cuve à ondes monte et descend mais ne s'éloigne pas de la source ; chaque point de la corde revient à sa position initiale après le passage de l'impulsion. Ce qui se propage, c'est \emph{l'état vibratoire} (la perturbation) et l'\emph{énergie mécanique} qu'il transporte. Retiens l'image de la « ola » dans un stade : la vague fait le tour des tribunes, mais chaque spectateur reste sur son siège — il se contente de se lever puis de se rasseoir.
\end{attention}

\begin{savaistu}[Les ondes sismiques : messagères de l'intérieur de la Terre]
Lors d'un séisme, la rupture des roches en profondeur émet simultanément deux familles d'ondes mécaniques : les ondes \textbf{P} (dites \emph{primaires}), qui compriment et dilatent les roches dans la direction de propagation, et les ondes \textbf{S} (dites \emph{secondaires}), qui cisaillement les roches perpendiculairement à la propagation. Les ondes P sont plus rapides (environ \SI{6}{km/s} dans la croûte) que les ondes S (environ \SI{3,5}{km/s}) : elles arrivent donc \emph{en premier} aux stations sismologiques. En mesurant l'écart $\Delta t$ entre les temps d'arrivée des ondes P et S en plusieurs stations, les sismologues calculent la distance de chaque station à l'épicentre et le localisent par triangulation. Le Cameroun, traversé par la ligne volcanique du mont Cameroun, dispose de stations de surveillance sismique autour de ce volcan actif.
\end{savaistu}

\courssub{Ondes transversales et ondes longitudinales}

L'expérience du ressort diffère de celle de la corde sur un point essentiel : la \emph{direction} du mouvement des éléments du milieu par rapport à la direction de propagation. Cela conduit à une classification fondamentale.

\begin{definition}[Onde transversale]
Une onde est dite \cle{transversale} lorsque la perturbation (le déplacement des points du milieu) est \textbf{perpendiculaire} à la direction de propagation.

Exemples : onde le long d'une corde tendue, rides à la surface de l'eau (en première approximation), ondes sismiques S.
\end{definition}

\begin{definition}[Onde longitudinale]
Une onde est dite \cle{longitudinale} lorsque la perturbation (le déplacement des points du milieu) est \textbf{parallèle} à la direction de propagation. Le milieu présente alors des zones de \emph{compression} et des zones de \emph{dilatation} qui se propagent.

Exemples : onde le long d'un ressort comprimé, son dans l'air (les couches d'air se compriment et se dilatent), ondes sismiques P.
\end{definition}

\begin{popfigure}
\begin{center}
\begin{tikzpicture}[scale=0.9]
% Corde à trois instants
\foreach \y/\lab/\dx in {0/{instant $t_1$}/0, -2.2/{instant $t_2 > t_1$}/1.6, -4.4/{instant $t_3 > t_2$}/3.2}{
  \draw[popline, thick] (0,\y) -- (10.5,\y);
  \draw[popink, very thick, domain=0:2.2, samples=60] plot ({\x+\dx}, {\y + 1.1*exp(-((\x-1.1)^2)/0.35)});
  \fill[poporange] ({\dx+1.1}, {\y+1.1}) circle (2pt);
  \node[right, popdark] at (10.6,\y) {\lab};
}
\draw[-{Stealth[length=3mm]}, poporange, very thick] (1.1,1.0) -- (4.3,1.0) node[midway, above, popdark]{propagation};
\node[popdark, left] at (0,0.55) {impulsion};
\end{tikzpicture}
\end{center}
\legende{Propagation d'une impulsion transversale le long d'une corde : la déformation se déplace vers la droite, mais chaque point de la corde ne se déplace que verticalement.}
\end{popfigure}

\begin{popfigure}
\begin{center}
\begin{tikzpicture}[scale=0.85]
% Ressort au repos (ligne de référence)
\draw[popline, dashed] (0,0) -- (12,0);
\node[left, popmuted] at (0,0) {repos};
% Ressort perturbé : spires resserrées puis écartées
\begin{scope}[yshift=-2cm]
\draw[popink, thick, decorate, decoration={coil, aspect=0.4, segment length=4mm, amplitude=3mm}] (0,0) -- (2.5,0);
\draw[popink, thick, decorate, decoration={coil, aspect=0.4, segment length=1.6mm, amplitude=3mm}] (2.5,0) -- (5,0);
\draw[popink, thick, decorate, decoration={coil, aspect=0.4, segment length=4mm, amplitude=3mm}] (5,0) -- (7.5,0);
\draw[popink, thick, decorate, decoration={coil, aspect=0.4, segment length=8mm, amplitude=3mm}] (7.5,0) -- (10.5,0);
\draw[popink, thick, decorate, decoration={coil, aspect=0.4, segment length=4mm, amplitude=3mm}] (10.5,0) -- (12,0);
\draw[poporange, very thick, -{Stealth[length=3mm]}] (3.75,0.9) -- (6.75,0.9) node[midway, above, popdark]{propagation};
\draw[popblue, thick, decorate, decoration={brace, amplitude=5pt, mirror}] (2.5,-0.55) -- (5,-0.55) node[midway, below=6pt, popblue]{compression};
\draw[popgreen, thick, decorate, decoration={brace, amplitude=5pt, mirror}] (7.5,-0.55) -- (10.5,-0.55) node[midway, below=6pt, popgreen]{dilatation};
\end{scope}
\end{tikzpicture}
\end{center}
\legende{Onde longitudinale sur un ressort : une zone de compression (spires resserrées) et une zone de dilatation (spires écartées) se propagent ; chaque spire oscille horizontalement, dans la direction de propagation.}
\end{popfigure}

\courssub{Célérité d'une onde}

\begin{definition}[Célérité]
La \cle{célérité} $c$ d'une onde est la vitesse de propagation de la perturbation dans le milieu. Si la perturbation parcourt la distance $d$ pendant la durée $\Delta t$, alors :
\[ c = \dfrac{d}{\Delta t} \]
avec $d$ en mètres (\si{m}), $\Delta t$ en secondes (\si{s}) et $c$ en mètres par seconde (\si{m/s}).
\end{definition}

\begin{propriete}[Facteurs dont dépend la célérité]
La célérité d'une onde mécanique \textbf{dépend uniquement des propriétés du milieu de propagation} (rigidité, inertie, température), et non de la forme ni de l'amplitude de la perturbation (pour des perturbations faibles).

Pour une corde tendue, on admet la relation :
\[ c = \sqrt{\dfrac{T}{\mu}} \]
où $T$ est la tension de la corde (en newtons) et $\mu$ sa masse linéique (masse par unité de longueur, en \si{kg/m}).

\emph{Vérification dimensionnelle :} $[T] = \si{kg.m.s^{-2}}$ et $[\mu] = \si{kg.m^{-1}}$, donc $\left[\sqrt{T/\mu}\right] = \sqrt{\si{m^2.s^{-2}}} = \si{m/s}$ : la relation est homogène.

Conséquence physique : plus la corde est tendue, plus l'onde est rapide ; plus la corde est lourde (grande inertie), plus l'onde est lente. C'est pourquoi les cordes graves d'une guitare sont plus épaisses que les cordes aiguës.
\end{propriete}

\begin{exemplebox}[Deux ordres de grandeur à connaître]
\begin{itemize}
\item Son dans l'air à \SI{20}{\celsius} : $c \approx \SI{340}{m/s}$. Le son d'un coup de tonnerre parcourt environ \SI{1}{km} en \SI{3}{s} : d'où la méthode populaire pour estimer la distance d'un orage.
\item Son dans l'eau : $c \approx \SI{1500}{m/s}$ ; dans l'acier : $c \approx \SI{5000}{m/s}$. Le son se propage d'autant plus vite que le milieu est rigide.
\end{itemize}
\end{exemplebox}

\begin{exoresolu}[Célérité et retard le long d'une corde]
Une impulsion transversale parcourt \SI{12}{m} de corde tendue en \SI{0,40}{s}.

\textbf{1.} Calculer la célérité de l'onde sur cette corde.

\textbf{2.} La corde a une masse linéique $\mu = \SI{25}{g/m}$. En déduire sa tension $T$.

\textbf{3.} Un point $M$ de la corde est situé à $d = \SI{6,0}{m}$ de la source $S$. Avec quel retard $M$ reproduit-il le mouvement de $S$ ?
\tcblower
\textbf{1. Célérité.} Par définition :
\[ c = \dfrac{d}{\Delta t} = \dfrac{12}{0{,}40} = \SI{30}{m/s}. \]

\textbf{2. Tension de la corde.} De $c = \sqrt{T/\mu}$ on tire $T = \mu c^2$. Conversion : $\mu = \SI{25}{g/m} = \SI{0,025}{kg/m}$. Donc :
\[ T = 0{,}025 \times 30^2 = 0{,}025 \times 900 = \SI{22,5}{N} \approx \SI{23}{N}. \]

\textbf{3. Retard du point $M$.} Le temps mis par la perturbation pour aller de $S$ à $M$ :
\[ \theta = \dfrac{d}{c} = \dfrac{6{,}0}{30} = \SI{0,20}{s}. \]
Le point $M$ reproduit donc le mouvement de la source avec \SI{0,20}{s} de retard.
\end{exoresolu}

\courssub{Onde progressive à une dimension : la notion de retard}

Considérons une corde tendue le long d'un axe $(Ox)$. Une source $S$, placée en $O$, est mise en mouvement : son élongation (déplacement transversal) à l'instant $t$ est notée $y_S(t)$. La perturbation se propage le long de la corde à la célérité $c$. Que fait un point $M$ de la corde, situé à la distance $d = SM$ de la source ?

\begin{propriete}[Propagation du mouvement avec retard — démonstration guidée]
\textbf{Hypothèses :} le milieu est homogène, la propagation se fait sans amortissement notable, à la célérité constante $c$.

\textbf{Étape 1 — Chaîne de transmission.} Le point $M$ ne vibre que parce que la perturbation lui est transmise de proche en proche depuis $S$. L'état vibratoire de $M$ à l'instant $t$ est donc exactement celui qu'avait la source à un instant antérieur.

\textbf{Étape 2 — Durée du trajet.} La perturbation parcourt la distance $d = SM$ à la célérité $c$ ; la durée du trajet est :
\[ \theta = \dfrac{d}{c}. \]

\textbf{Étape 3 — Traduction mathématique.} L'état vibratoire de $M$ à l'instant $t$ est celui de $S$ à l'instant $t - \theta$ :
\[ \boxed{\,y_M(t) = y_S(t - \theta) \quad \text{avec} \quad \theta = \dfrac{d}{c}\,} \]
Autrement dit : \emph{le mouvement de $M$ reproduit fidèlement celui de $S$, avec le retard $\theta$}.
\end{propriete}

\begin{methode}[Exploiter la relation $y_M(t) = y_S(t - \theta)$]
\begin{enumerate}
\item Repérer la distance $d = SM$ entre la source et le point étudié, et la célérité $c$ (donnée ou calculée par $c = d'/\Delta t$ à partir d'un trajet connu).
\item Calculer le retard $\theta = d/c$, en secondes.
\item Écrire l'élongation de $M$ en remplaçant $t$ par $(t - \theta)$ dans l'expression de $y_S$.
\item Pour comparer deux points $M_1$ et $M_2$ : calculer leurs retards $\theta_1$ et $\theta_2$ ; celui qui a le plus petit retard vibre « en avance » sur l'autre.
\end{enumerate}
\end{methode}

\begin{attention}[Pièges fréquents sur le retard]
\begin{itemize}
\item \textbf{Signe.} C'est $y_M(t) = y_S(t - \theta)$, avec un \emph{moins} : $M$ est \emph{en retard} sur $S$. Écrire $y_S(t+\theta)$ pour $M$ est une erreur classique qui inverse la causalité.
\item \textbf{Homogénéité.} $\theta = d/c$ exige $d$ en mètres et $c$ en \si{m/s} pour obtenir des secondes. Convertis systématiquement (cm $\to$ m, km/h $\to$ m/s).
\item \textbf{Le retard dépend de la position.} Deux points à des distances différentes de la source ont des retards différents : $\theta$ n'est pas une propriété de l'onde seule, mais du couple (source, point).
\end{itemize}
\end{attention}

\begin{exoresolu}[Équation du mouvement d'un point du milieu]
L'extrémité $S$ d'une corde est animée d'un mouvement dont l'élongation est, en centimètres :
\[ y_S(t) = 3\sin\left(20\pi\, t\right) \quad (t \text{ en secondes}). \]
L'onde se propage le long de la corde à la célérité $c = \SI{5,0}{m/s}$.

\textbf{1.} Déterminer l'élongation $y_M(t)$ du point $M$ situé à $d = \SI{75}{cm}$ de $S$.

\textbf{2.} Que vaut $y_M$ à l'instant $t = \SI{0,20}{s}$ ?
\tcblower
\textbf{1.} Le retard de $M$ sur $S$ vaut :
\[ \theta = \dfrac{d}{c} = \dfrac{0{,}75}{5{,}0} = \SI{0,15}{s}. \]
Le mouvement de $M$ reproduit celui de $S$ avec ce retard :
\[ y_M(t) = y_S(t - 0{,}15) = 3\sin\left[20\pi\,(t - 0{,}15)\right] = 3\sin\left(20\pi\, t - 3\pi\right) \quad (\text{en cm}). \]
On simplifie l'argument : $\sin(\alpha - 3\pi) = \sin(\alpha - \pi - 2\pi) = \sin(\alpha - \pi) = -\sin\alpha$, d'où $y_M(t) = -3\sin(20\pi\, t)$ : le point $M$ vibre \emph{en opposition de phase} avec la source (nous y reviendrons dans la section suivante).

\textbf{2.} À $t = \SI{0,20}{s}$ :
\[ y_M(0{,}20) = 3\sin(20\pi \times 0{,}20 - 3\pi) = 3\sin(4\pi - 3\pi) = 3\sin(\pi) = \SI{0}{cm}. \]
Vérification par la source : à l'instant $t - \theta = \SI{0,05}{s}$, $y_S(0{,}05) = 3\sin(\pi) = 0$ : le point $M$ reproduit bien, à $t = \SI{0,20}{s}$, l'état qu'avait $S$ à $t = \SI{0,05}{s}$.
\end{exoresolu}

\begin{aretenir}[L'essentiel de la section]
\begin{itemize}
\item Une \cle{onde mécanique progressive} est la propagation d'une perturbation dans un milieu matériel, \textbf{sans transport de matière}, \textbf{avec transport d'énergie}.
\item Onde \cle{transversale} : perturbation $\perp$ propagation (corde). Onde \cle{longitudinale} : perturbation $\parallel$ propagation (ressort, son).
\item Célérité : $c = \dfrac{d}{\Delta t}$ ; elle ne dépend que du milieu. Pour une corde : $c = \sqrt{T/\mu}$.
\item Tout point $M$ à la distance $d$ de la source reproduit le mouvement de $S$ avec le retard $\theta = \dfrac{d}{c}$ : $y_M(t) = y_S(t - \theta)$.
\end{itemize}
\end{aretenir}

Nous avons décrit la propagation d'une perturbation quelconque, en particulier d'une impulsion brève. Dans la section suivante, nous étudierons le cas capital où la source vibre \emph{sinusoïdalement} : l'onde prend alors une structure doublement périodique — dans le temps avec la période $T$, et dans l'espace avec la longueur d'onde $\lambda = cT$ — qui fera l'objet de nombreuses questions au Baccalauréat.

\coursec{Onde progressive sinusoïdale : double périodicité}

Dans la section précédente, nous avons introduit la notion d'onde mécanique progressive : une perturbation qui se propage dans un milieu matériel sans transport de matière, caractérisée par sa célérité $c$. Nous distinguons les ondes transversales (perturbation perpendiculaire à la direction de propagation, comme sur une corde ou à la surface de l'eau du lac de Lagdo) des ondes longitudinales (perturbation parallèle à la propagation, comme le son dans l'air ou les ondes sismiques P).

Nous allons maintenant nous concentrer sur le cas particulier, mais fondamental, où la source est \emph{sinusoïdale}. C'est le modèle de base pour décrire une onde radio émise par l'antenne de la CRTV au mont Fébé, une note de balafon accordée, ou le signal électrique alternatif du réseau ENEO (50 Hz). La sinusoïdale est la « brique élémentaire » de tout signal périodique (théorème de Fourier), et sa propagation révèle une structure spatio-temporelle d'une grande beauté : la \textbf{double périodicité}.

\courssub{Source sinusoïdale et période temporelle}

Considérons une source $S$ (un haut-parleur, une membrane, l'extrémité d'une corde vibrante) qui impose au milieu, en sa position $x=0$, une élongation $y_S$ variant sinusoïdalement avec le temps.

\begin{definition}[Grandeur sinusoïdale, période, fréquence, pulsation]
Une grandeur $y$ est \textbf{sinusoïdale} si elle s'écrit sous la forme :
\[ y(t) = a \cos(\omega t + \varphi_0) \quad \text{ou} \quad y(t) = a \sin(\omega t + \varphi_0) \]
où :
\begin{itemize}
    \item $a > 0$ est l'\textbf{amplitude} (élongation maximale, en mètres \si{\meter} pour une onde mécanique).
    \item $\omega = \dfrac{2\pi}{T}$ est la \textbf{pulsation} (en \si{\radian\per\second}), liée à la période $T$.
    \item $T$ est la \textbf{période} (en secondes \si{\second}) : plus petit intervalle de temps après lequel le mouvement se reproduit à l'identique ($y(t+T)=y(t)$).
    \item $f = \dfrac{1}{T}$ est la \textbf{fréquence} (en hertz \si{\hertz}).
    \item $\varphi_0$ est la \textbf{phase à l'origine} (en radians), état du système à $t=0$.
\end{itemize}
L'analyse dimensionnelle confirme l'homogénéité : $\omega t$ et $\varphi_0$ sont des angles (sans dimension), $a$ a la dimension de $y$.
\end{definition}

L'équation horaire de la source $S$ s'écrit donc :
\[ y_S(t) = a \cos(\omega t + \varphi_S) \]
où $\varphi_S$ est la phase initiale de la source. En un point $M$ quelconque du milieu, la vibration sera également sinusoïdale, de même fréquence $f$ (imposée par la source), mais décalée dans le temps à cause de la propagation.

\courssub{Retard de propagation et périodicité spatiale}

L'onde met un certain temps pour aller de la source $S$ jusqu'au point $M$ situé à la distance $d = SM$. Ce temps est le \textbf{retard} $\theta$ :
\[ \theta = \frac{d}{c} \]
Le point $M$ se met à vibrer \emph{après} la source. Son élongation à l'instant $t$ est exactement celle qu'avait la source à l'instant $t-\theta$ :
\[ y_M(t) = y_S(t - \theta) = a \cos\bigl[\omega(t - \theta) + \varphi_S\bigr] \]

\begin{attention}[Sens du retard : piège classique]
\emph{Erreur fréquente :} écrire $y_M(t) = y_S(t + \theta)$.
\emph{Correct :} $y_M(t) = y_S(t - \theta)$.
\textbf{Raisonnement physique :} À l'instant $t$, le point $M$ « reçoit » l'information que la source a émise \emph{plus tôt}, à l'instant $t-\theta$. Le retard \emph{retarde} la phase : $\varphi_M = \varphi_S - \omega\theta$. Sur un oscillogramme, la courbe de $M$ est décalée vers la \emph{droite} (vers les temps plus grands) par rapport à celle de $S$.
\end{attention}

Remplaçons $\theta$ par $d/c$ et $\omega$ par $2\pi/T = 2\pi f$ :
\[ y_M(t) = a \cos\left[ \omega t - \omega\frac{d}{c} + \varphi_S \right] = a \cos\left( \omega t - \frac{2\pi d}{\lambda} + \varphi_S \right) \]
où nous avons introduit une nouvelle grandeur fondamentale : la \textbf{longueur d'onde} $\lambda$.

\begin{definition}[Longueur d'onde $\lambda$ et double périodicité]
La \textbf{longueur d'onde} $\lambda$ (en mètres \si{\meter}) est la distance parcourue par l'onde pendant une période $T$ :
\[ \lambda = c\,T = \frac{c}{f} \]
C'est la \textbf{périodicité spatiale} : à un instant $t_0$ fixé, si l'on « photographie » l'onde, le profil spatial $y(x)$ se répète tous les $\lambda$.
\begin{itemize}
    \item \textbf{Périodicité temporelle :} En un point $M$ fixé ($d$ constant), $y_M(t+T) = y_M(t)$. Période $T$.
    \item \textbf{Périodicité spatiale :} À un instant $t$ fixé, $y(x+\lambda, t) = y(x, t)$. Période spatiale $\lambda$.
\end{itemize}
\end{definition}

\begin{propriete}[Relation fondamentale $\lambda = cT$ — Démonstration]
\emph{Énoncé :} Pour une onde progressive sinusoïdale de célérité $c$, de période $T$ et de longueur d'onde $\lambda$, on a $\lambda = cT$.
\emph{Démonstration :}
Considérons un crest (creux ou ventre) de l'onde au niveau de la source $S$ à l'instant $t=0$. Ce crest se propage à la vitesse $c$. Au bout d'une période $T$, la source émet un nouveau crest identique (périodicité temporelle). Le premier crest a parcouru la distance $d = c \times T$. Par définition, la distance séparant deux crests successifs \emph{au même instant} est la longueur d'onde $\lambda$. Donc $\lambda = cT$. \qed
\emph{Analyse dimensionnelle :} $[c] = \si{\meter\per\second}$, $[T] = \si{\second}$, $[\lambda] = \si{\meter}$. L'égalité est homogène.
\end{propriete}

\begin{savaistu}[Le son dans l'air au Cameroun]
La célérité du son dans l'air dépend de la température : $c \approx 331 + 0,6\,\theta$ (avec $\theta$ en $^\circ\mathrm{C}$).
À Yaoundé ($\theta \approx \SI{25}{\celsius}$), $c \approx \SI{346}{\meter\per\second}$.
Pour la note \textbf{La}$_3$ du diapason ($f = \SI{440}{\hertz}$), la longueur d'onde est $\lambda = c/f \approx \SI{0,79}{\meter}$.
Pour les ultrasons d'une échographie médicale ($f = \SI{2}{\mega\hertz}$), qui se propagent dans les tissus du corps à $c \approx \SI{1540}{\meter\per\second}$, $\lambda \approx \SI{0,77}{\milli\meter}$ : assez fin pour « voir » un fœtus !
Pour les ondes radio FM (ex. Radio Environnement \SI{94,0}{\mega\hertz}), $\lambda \approx \SI{3,19}{\meter}$ (ondes électromagnétiques, $c=\SI{3e8}{\meter\per\second}$).
\end{savaistu}

\courssub{Équation du mouvement d'un point du milieu}

Nous avons maintenant tous les éléments pour écrire l'équation complète.

\begin{propriete}[Équation horaire d'un point M — Démonstration guidée]
\emph{Énoncé :} Pour une onde progressive sinusoïdale de célérité $c$, d'amplitude $a$, de pulsation $\omega$, de phase à l'origine $\varphi_S$ à la source $S$, l'élongation d'un point $M$ situé à la distance $d = SM$ est :
\[ y_M(t) = a \cos\left[ \omega\left(t - \frac{d}{c}\right) + \varphi_S \right] = a \cos\left( \omega t - \frac{2\pi d}{\lambda} + \varphi_S \right) \]
\emph{Étapes de la démonstration :}
\begin{enumerate}
    \item Écrire l'équation de la source : $y_S(t) = a \cos(\omega t + \varphi_S)$.
    \item Exprimer le retard de propagation : $\theta = d/c$.
    \item Le point $M$ vibre comme la source \emph{avait vibré} il y a un temps $\theta$ : $y_M(t) = y_S(t - \theta)$.
    \item Substituer $\theta$ et $\omega = 2\pi/T = 2\pi c/\lambda$ : $\omega\theta = \omega d/c = 2\pi d/\lambda$.
    \item On obtient la forme canonique : $y_M(t) = a \cos(\omega t - k d + \varphi_S)$ avec $k = 2\pi/\lambda$ (nombre d'onde).
\end{enumerate}
\emph{Remarque :} Le signe $-$ devant $kd$ correspond à une onde se propageant dans le sens $x$ croissant (de $S$ vers $M$). Pour une onde revenant (réfléchie), le signe serait $+$.
\end{propriete}

\begin{exemplebox}[Calcul de longueur d'onde et déphasage]
Une onde mécanique progressive sinusoïdale se propage sur une corde tendue avec une célérité $c = \SI{2,00}{\meter\per\second}$. La fréquence de l'excitateur est $f = \SI{5,00}{\hertz}$.
\begin{enumerate}
    \item Calculer la période $T$ et la longueur d'onde $\lambda$.
    \item Deux points $M_1$ et $M_2$ sont distants de $d = \SI{0,60}{\meter}$. Sont-ils en phase ou en opposition de phase ?
\end{enumerate}
\textbf{Solution :}
\begin{enumerate}
    \item $T = 1/f = \SI{0,200}{\second}$.
    $\lambda = cT = 2,00 \times 0,200 = \SI{0,400}{\meter}$.
    \item $d = \SI{0,60}{\meter} = 1,5 \times \lambda = 3\lambda/2$.
    Le déphasage est $\Delta\varphi = \dfrac{2\pi d}{\lambda} = \dfrac{2\pi \times 1,5\lambda}{\lambda} = 3\pi \equiv \pi \pmod{2\pi}$.
    Les deux points vibrent en \textbf{opposition de phase} (lorsque $M_1$ est au ventre haut, $M_2$ est au ventre bas).
\end{enumerate}
\end{exemplebox}

\courssub{Déphasage entre deux points : en phase, en opposition de phase}

Soient deux points $M_1$ et $M_2$ sur la trajectoire de l'onde, aux distances $d_1$ et $d_2$ de la source. Leurs élongations sont :
\[ y_1(t) = a \cos(\omega t - k d_1 + \varphi_S), \quad y_2(t) = a \cos(\omega t - k d_2 + \varphi_S) \]
Le déphasage $\Delta\varphi = \varphi_2 - \varphi_1$ (phase de $M_2$ moins phase de $M_1$) est :
\[ \Delta\varphi = (-k d_2) - (-k d_1) = k(d_1 - d_2) = \frac{2\pi}{\lambda}\,|d_1 - d_2| \]
En posant $d = |M_1M_2| = |d_1 - d_2|$ (distance entre les deux points le long de l'axe de propagation) :

\begin{propriete}[Déphasage spatial et conditions de phase]
Le déphasage entre deux points $M_1$ et $M_2$ distants de $d$ le long de la direction de propagation est :
\[ \Delta\varphi = \frac{2\pi d}{\lambda} \]
\begin{itemize}
    \item \textbf{Points en phase :} $\Delta\varphi = 2k\pi \;\Leftrightarrow\; d = k\lambda \quad (k \in \mathbb{N})$.
    Ils ont le même allongement et la même vitesse à tout instant.
    \item \textbf{Points en opposition de phase :} $\Delta\varphi = (2k+1)\pi \;\Leftrightarrow\; d = (2k+1)\frac{\lambda}{2} \quad (k \in \mathbb{N})$.
    Leurs allongements sont opposés à tout instant : $y_2(t) = -y_1(t)$.
\end{itemize}
\end{propriete}

\begin{methode}[Comparer l'état vibratoire de deux points M1 et M2]
\begin{enumerate}
    \item Identifier la longueur d'onde $\lambda$ (donnée ou calculée via $\lambda = c/f$).
    \item Mesurer ou calculer la distance $d = |M_1M_2|$ le long de l'axe de propagation.
    \item Calculer le rapport $d/\lambda$.
    \item En déduire le déphasage $\Delta\varphi = 2\pi (d/\lambda)$.
    \item Conclure :
    \begin{itemize}
        \item Si $d/\lambda$ est un entier $\to$ \textbf{en phase}.
        \item Si $d/\lambda$ est un demi-entier ($0,5; 1,5; 2,5\dots$) $\to$ \textbf{en opposition de phase}.
        \item Sinon $\to$ \textbf{déphasés} (ex: $d = \lambda/4 \Rightarrow \Delta\varphi = \pi/2$, quadrature de phase).
    \end{itemize}
\end{enumerate}
\end{methode}

\begin{attention}[Distance $d$ : le long de l'onde !]
La distance $d$ à utiliser dans $\Delta\varphi = 2\pi d/\lambda$ est la distance \emph{projetée sur la direction de propagation} (distance le long de l'axe de l'onde). Si $M_1$ et $M_2$ ne sont pas sur le même rayon (ex: ondes circulaires à la surface de l'eau), on utilise la différence de distances à la source : $d = |SM_2 - SM_1|$. Ne pas confondre avec la distance euclidienne « à vol d'oiseau » si les points ne sont pas alignés avec la source.
\end{attention}

\courssub{Lecture conjointe des graphes $y(t)$ et $y(x)$}

Au Baccalauréat, on vous présente souvent deux graphiques pour la même onde :
\begin{itemize}
    \item Le \textbf{chronogramme} $y(t)$ en un point fixe $M$ : on voit la période $T$.
    \item L'\textbf{instantané spatial} $y(x)$ à un instant fixe $t_0$ : on voit la longueur d'onde $\lambda$.
\end{itemize}
Ces deux graphiques sont des sinusoïdes de même amplitude $a$, mais l'abscisse change : temps $t$ (secondes) vs position $x$ (mètres). La célérité $c$ fait le lien : $\lambda = cT$.

\begin{popfigure}
\begin{tikzpicture}
% Premier graphique : Chronogramme
\begin{axis}[
    name=plot1,
    width=0.45\linewidth, height=5cm,
    axis lines=middle,
    xlabel={$t\,(\si{\second})$}, ylabel={$y\,(\si{\meter})$},
    domain=0:4, xmin=0, xmax=4.2,
    xtick={0,1,2,3,4},
    xticklabels={$0$, $T$, $2T$, $3T$, $4T$},
    title={Chronogramme $y_M(t)$ en un point $M$},
    ymin=-1.5, ymax=1.5,
    tick label style={font=\small},
    every axis plot post/.append style={popcurve, thick, samples=200},
    minor tick num=1, grid=both, grid style={popmuted, thin},
]
\addplot[domain=0:4] {cos(deg(2*pi*x))};
\draw[poporange, <->, thick] (axis cs:0,1.2) -- (axis cs:1,1.2) node[midway, above, font=\small, poporange] {$T$};
\draw[poporange, dashed] (axis cs:1, -1.5) -- (axis cs:1, 1.5);
\end{axis}

% Deuxième graphique : Instantané spatial
\begin{axis}[
    at={($(plot1.outer south east)+(1.5cm,0)$)},
    anchor=outer south west,
    width=0.45\linewidth, height=5cm,
    axis lines=middle,
    xlabel={$x\,(\si{\meter})$}, ylabel={$y\,(\si{\meter})$},
    domain=0:4, xmin=0, xmax=4.2,
    xtick={0,1,2,3,4},
    xticklabels={$0$, $\lambda$, $2\lambda$, $3\lambda$, $4\lambda$},
    title={Instantané spatial $y(x, t_0)$},
    ymin=-1.5, ymax=1.5,
    tick label style={font=\small},
    every axis plot post/.append style={popcurve, thick, samples=200},
    minor tick num=1, grid=both, grid style={popmuted, thin},
]
\addplot[domain=0:4] {cos(deg(2*pi*x))};
\draw[popgreen, <->, thick] (axis cs:0,1.2) -- (axis cs:1,1.2) node[midway, above, font=\small, popgreen] {$\lambda$};
\draw[popgreen, dashed] (axis cs:1, -1.5) -- (axis cs:1, 1.5);
\node[popblue, circle, fill, inner sep=1.5pt, label=above:{$M_1$}] at (axis cs:0,1) {};
\node[popblue, circle, fill, inner sep=1.5pt, label=above:{$M_2$}] at (axis cs:1,1) {};
\node[poppink, circle, fill, inner sep=1.5pt, label=below:{$M_3$}] at (axis cs:0.5,-1) {};
\draw[popblue, <->, dashed] (axis cs:0, -1.3) -- (axis cs:1, -1.3) node[midway, below, font=\small, popblue] {$d=\lambda$ (en phase)};
\draw[poppink, <->, dashed] (axis cs:0, -1.45) -- (axis cs:0.5, -1.45) node[midway, below, font=\small, poppink] {$d=\lambda/2$ (oppos.)};
\end{axis}
\end{tikzpicture}
\legende{Double périodicité : à gauche, vibration temporelle en un point (période $T$) ; à droite, forme de l'onde dans l'espace à un instant donné (période spatiale $\lambda$). Points $M_1$ et $M_2$ séparés de $\lambda$ : en phase. $M_1$ et $M_3$ séparés de $\lambda/2$ : en opposition de phase.}
\end{popfigure}

\begin{exoresolu}[Analyse complète d'une onde sur corde]
\textbf{Énoncé :} Une corde tendue horizontalement porte une onde progressive sinusoïdale se propageant vers la droite. La figure ci-contre (non fournie ici, on utilise les données) donne l'allure de l'onde à l'instant $t=0$. On sait que $c = \SI{10}{\meter\per\second}$ et que l'amplitude $a = \SI{0,02}{\meter}$.
\begin{center}
\begin{tikzpicture}[scale=0.8, domain=0:4]
\draw[->] (-0.5,0) -- (4.5,0) node[right] {$x\,(\si{\meter})$};
\draw[->] (0,-1.5) -- (0,1.5) node[above] {$y\,(\si{\centi\meter})$};
\draw[popcurve, thick, samples=200] plot (\x, {1*cos(deg(2*pi*\x))});
\foreach \x/\label in {0/0, 1/$\lambda$, 2/$2\lambda$, 3/$3\lambda$, 4/$4\lambda$}
    \draw (\x, -0.1) -- (\x, 0.1) node[below, font=\small] {\label};
\draw[dashed, popmuted] (0,1) -- (1,1);
\draw[dashed, popmuted] (0,-1) -- (1,-1);
\node[popblue, circle, fill, inner sep=2pt] at (0,1) {}; \node[above, font=\small] at (0,1.3) {$A$};
\node[popblue, circle, fill, inner sep=2pt] at (0.5,-1) {}; \node[below, font=\small] at (0.5,-1.3) {$B$};
\node[popblue, circle, fill, inner sep=2pt] at (1,1) {}; \node[above, font=\small] at (1,1.3) {$C$};
\end{tikzpicture}
\end{center}
Données lecture graphique : $\lambda = \SI{1,00}{\meter}$. Points $A(x=0)$, $B(x=\SI{0,50}{\meter})$, $C(x=\SI{1,00}{\meter})$.
\begin{enumerate}
    \item Déterminer la fréquence $f$ et la pulsation $\omega$.
    \item Écrire l'équation horaire $y_A(t)$ du point $A$ sachant qu'à $t=0$, $y_A = +a$ et la vitesse est nulle (donc $\varphi_A=0$).
    \item Écrire l'équation horaire $y_B(t)$ et $y_C(t)$.
    \item Préciser les relations de phase entre $A$ et $B$, $A$ et $C$, $B$ et $C$.
\end{enumerate}

\tcblower

\textbf{Solution détaillée :}
\begin{enumerate}
    \item $\lambda = \SI{1,00}{\meter}$, $c = \SI{10}{\meter\per\second}$.
    $T = \lambda/c = \SI{0,100}{\second}$.
    $f = 1/T = \SI{10,0}{\hertz}$.
    $\omega = 2\pi f = \SI{20\pi}{\radian\per\second} \approx \SI{62,8}{\radian\per\second}$.
    \item Point $A$ ($x=0$) : $y_A(0) = a = \SI{0,02}{\meter}$, $v_A(0)=0 \Rightarrow \varphi_A = 0$.
    $y_A(t) = \SI{0,02}{\meter} \cos(20\pi\, t)$.
    \item Point $B$ ($x=\SI{0,50}{\meter} = \lambda/2$) : retard $\theta_B = d/c = 0,50/10 = \SI{0,050}{\second} = T/2$.
    $\Delta\varphi_{B/A} = \omega\theta_B = 20\pi \times 0,05 = \pi$.
    $y_B(t) = \SI{0,02}{\meter} \cos(20\pi t - \pi) = -\SI{0,02}{\meter} \cos(20\pi t) = -y_A(t)$.
    Point $C$ ($x=\SI{1,00}{\meter} = \lambda$) : retard $\theta_C = 1/10 = \SI{0,100}{\second} = T$.
    $\Delta\varphi_{C/A} = 2\pi \equiv 0$.
    $y_C(t) = \SI{0,02}{\meter} \cos(20\pi t - 2\pi) = \SI{0,02}{\meter} \cos(20\pi t) = y_A(t)$.
    \item $A$ et $B$ : distance $\lambda/2 \Rightarrow$ \textbf{opposition de phase}.
    $A$ et $C$ : distance $\lambda \Rightarrow$ \textbf{en phase}.
    $B$ et $C$ : distance $\lambda/2 \Rightarrow$ \textbf{opposition de phase}.
\end{enumerate}
\end{exoresolu}

\courssub{Synthèse : les formules-clés à maîtriser}

\begin{aretenir}[Boîte à outils : Onde progressive sinusoïdale]
Pour une onde progressive sinusoïdale de célérité $c$, fréquence $f$, période $T$, longueur d'onde $\lambda$, amplitude $a$, nombre d'onde $k = 2\pi/\lambda$, pulsation $\omega = 2\pi f$ :
\begin{align*}
    &\textbf{Relations fondamentales :} & \lambda &= c T = \frac{c}{f}, \quad \omega = \frac{2\pi}{T}, \quad k = \frac{2\pi}{\lambda}, \quad \omega = c k. \\[4pt]
    &\textbf{Équation horaire en } M(d) : & y_M(t) &= a \cos\left[ \omega\left(t - \frac{d}{c}\right) + \varphi_S \right] = a \cos(\omega t - k d + \varphi_S). \\[4pt]
    &\textbf{Déphasage } M_1M_2 (d) : & \Delta\varphi &= \varphi_2 - \varphi_1 = -k(d_2 - d_1) = \frac{2\pi}{\lambda}\,|d_2 - d_1|. \\[4pt]
    &\textbf{En phase :} & d &= k\lambda \quad (k \in \mathbb{Z}). \\
    &\textbf{Opposition de phase :} & d &= (2k+1)\frac{\lambda}{2} \quad (k \in \mathbb{Z}).
\end{align*}
\textbf{Unités SI :} $[c]=\si{\meter\per\second}$, $[f]=\si{\hertz}$, $[T]=\si{\second}$, $[\lambda]=\si{\meter}$, $[\omega]=\si{\radian\per\second}$, $[k]=\si{\radian\per\meter}$, $[a]=\si{\meter}$.
\end{aretenir}

\begin{exercice}[Entraînement : Onde sonore et interférences]
Deux haut-parleurs $H_1$ et $H_2$, séparés de $d = \SI{1,70}{\meter}$, émettent le même son pur de fréquence $f = \SI{680}{\hertz}$ en phase. La célérité du son est $c = \SI{340}{\meter\per\second}$. Un auditeur se place au point $M$ tel que $MH_1 = \SI{3,40}{\meter}$ et $MH_2 = \SI{4,25}{\meter}$.
\begin{enumerate}
    \item Calculer $\lambda$ et $k$.
    \item Calculer la différence de marche $\Delta d = |MH_2 - MH_1|$.
    \item Déterminer le déphasage $\Delta\varphi$ entre les ondes arrivant en $M$.
    \item L'auditeur perçoit-il un maximum ou un minimum d'intensité ? Justifier.
\end{enumerate}
\end{exercice}

\begin{corrige}[Exercice : Onde sonore et interférences]
\begin{enumerate}
    \item $\lambda = c/f = 340/680 = \SI{0,500}{\meter}$. $k = 2\pi/\lambda = 4\pi \approx \SI{12,6}{\radian\per\meter}$.
    \item $\Delta d = |4,25 - 3,40| = \SI{0,85}{\meter}$.
    \item $\Delta\varphi = k\Delta d = 4\pi \times 0,85 = 3,4\pi = 2\pi + 1,4\pi \equiv 1,4\pi \pmod{2\pi}$.
    Ou via rapport : $\Delta d/\lambda = 0,85/0,50 = 1,7 = 1 + 0,7$. Ni entier, ni demi-entier.
    \item Ni maximum (phase), ni minimum (opposition). Interférence \textbf{intermédiaire}. L'intensité est entre le max et le min.
\end{enumerate}
\end{corrige}

Cette section achève la description mathématique de l'onde progressive sinusoïdale. La prochaine et dernière section (Synthèse) regroupera les méthodes types de résolution de problèmes (exploitation d'oscillogramme, stroboscopie, détermination de $\lambda$ et $c$, etc.) pour vous préparer efficacement aux épreuves du Baccalauréat.

\coursec{Synthèse et méthodes de résolution}

Dans la section précédente, nous avons établi le cœur de la leçon : une onde progressive sinusoïdale possède une \cle{double périodicité} — temporelle ($T$) et spatiale ($\lambda = cT$) — et le mouvement d'un point $M$ du milieu reproduit celui de la source avec le \cle{retard} $\theta = d/c$. Il est temps maintenant d'assembler toutes les pièces du puzzle. Cette section de synthèse te donne une vision globale de la chaîne « source $\to$ propagation $\to$ point $M$ », puis trois méthodes types qui couvrent la quasi-totalité des sujets de Baccalauréat sur ce chapitre, enfin un tableau récapitulatif et la liste des pièges classiques. Travaille-la avec soin : c'est ici que se jouent les points au Bac.

\courssub{Bilan : la chaîne complète source $\to$ propagation $\to$ point $M$}

Toute situation d'onde mécanique progressive sinusoïdale se décompose en quatre maillons logiques :

\begin{enumerate}
\item \textbf{La source} $S$ vibre sinusoïdalement : son élongation s'écrit $y_S(t) = a\cos(\omega t + \varphi)$, avec $a$ l'amplitude, $\omega = 2\pi f = \dfrac{2\pi}{T}$ la pulsation.
\item \textbf{Le milieu} transmet la perturbation \emph{sans transport de matière}, à la célérité $c$ qui ne dépend que du milieu (tension et masse linéique pour une corde, rigidité et densité pour l'air, etc.).
\item \textbf{Le point $M$}, situé à la distance $d = SM$ de la source, reproduit le mouvement de $S$ avec le retard $\theta = \dfrac{d}{c}$ :
\[ y_M(t) = y_S(t - \theta) = a\cos\left[\omega\left(t - \dfrac{d}{c}\right) + \varphi\right]. \]
\item \textbf{La comparaison de deux points} $M$ et $N$ repose sur le déphasage $\Delta\varphi = \dfrac{2\pi}{\lambda}\,|d_N - d_M|$ : les points vibrent \cle{en phase} si leur distance est un multiple entier de $\lambda$, \cle{en opposition de phase} si c'est un multiple impair de $\lambda/2$.
\end{enumerate}

\begin{popfigure}
\begin{tikzpicture}[font=\small, >=Stealth]
% Source
\draw[fill=popblueL] (0,0) circle (0.35);
\node[popblue] at (0,-0.75) {\textbf{Source $S$}};
\node[align=center] at (0,1.0) {$y_S(t) = a\cos(\omega t + \varphi)$};
% Milieu (corde ondulée)
\draw[poporange, thick, domain=0.4:6.6, samples=120] plot (\x, {0.45*sin(180*\x)});
\node[poporange] at (3.5,1.0) {milieu de propagation, célérité $c$};
% Point M
\draw[fill=poppurple] (5.5,0) circle (0.12);
\node[poppurple] at (5.5,-0.75) {\textbf{Point $M$}};
\node[align=center, poppurple] at (5.5,1.6) {$y_M(t) = a\cos\left[\omega\left(t - \dfrac{d}{c}\right) + \varphi\right]$};
% distance d
\draw[<->, popdark] (0.4,-1.3) -- (5.5,-1.3) node[midway, below] {$d = SM$, \quad retard $\theta = \dfrac{d}{c}$};
% flèche propagation
\draw[->, popgreen, very thick] (0.6,0.85) -- (2.4,0.85) node[midway, above] {$\vec{c}$};
\end{tikzpicture}
\legende{Schéma-bilan : la source impose sa vibration, le milieu la transporte à la célérité $c$, le point $M$ la reproduit avec le retard $\theta = d/c$.}
\end{popfigure}

\begin{aretenir}[L'idée directrice]
L'onde, c'est de la \emph{vibration qui voyage}. La source donne le « tempo » ($T$, $\omega$), le milieu donne la « vitesse de voyage » ($c$), et leur produit donne la « taille du motif » dans l'espace : $\lambda = cT$. Tout le reste — retards, déphasages, comparaisons de points — se déduit de ces trois relations.
\end{aretenir}

\courssub{Méthode type 1 : déterminer $c$, $T$, $\lambda$ à partir de données expérimentales}

Au laboratoire comme au Bac, les grandeurs ne sont jamais données directement : il faut les extraire d'un oscillogramme, d'une observation stroboscopique ou d'une mesure de retard.

\begin{methode}[Extraire $T$, $c$, $\lambda$ de données expérimentales]
\begin{enumerate}
\item \textbf{Période $T$} : sur un oscillogramme, compter le nombre de divisions horizontales d'un motif complet et multiplier par le balayage (en s/div). En stroboscopie, la plus grande fréquence d'immobilité apparente est la fréquence propre $f$ du phénomène, d'où $T = 1/f$.
\item \textbf{Célérité $c$} : mesurer la distance $d$ parcourue et le retard $\theta$ (décalage temporel entre deux traces sur l'oscillogramme, ou durée mesurée au chronomètre) : $c = \dfrac{d}{\theta}$.
\item \textbf{Longueur d'onde $\lambda$} : appliquer $\lambda = cT$, ou mesurer directement sur une photo (ou sous éclairage stroboscopique) la distance entre deux crêtes consécutives.
\item \textbf{Vérification croisée} : contrôler que $\lambda = cT = \dfrac{c}{f}$ est cohérent avec les trois mesures. Au Bac, cette vérification est souvent une question à part entière.
\end{enumerate}
\end{methode}

\begin{exemplebox}[Mesure de $\lambda$ par stroboscopie sur onde progressive]
On observe une onde progressive sur une corde (ou à la surface de l'eau dans une cuve à ondes) à l'aide d'un stroboscope réglé sur la fréquence de la source $f_e = \SI{50}{Hz}$. L'immobilité apparente permet de « photographier » l'onde : on mesure la distance entre deux crêtes successives, qui vaut $\SI{20}{cm}$.

\textbf{Exploitation :} La fréquence de l'onde est $f = f_e = \SI{50}{Hz}$, d'où $T = \dfrac{1}{50} = \SI{0,020}{s}$. La distance mesurée correspond à la longueur d'onde $\lambda = \SI{0,20}{m}$. La célérité est $c = \lambda f = 0{,}20 \times 50 = \SI{10}{m.s^{-1}}$.
\end{exemplebox}

\courssub{Méthode type 2 : établir l'équation horaire d'un point $M$ puis comparer deux points}

C'est l'exercice le plus fréquent au Baccalauréat. La démarche est toujours la même.

\begin{methode}[Fiche de résolution type Bac en 5 étapes]
\begin{enumerate}
\item \textbf{Convertir toutes les données en unités SI} : distances en mètres, temps en secondes, amplitudes en mètres (ou en cm, mais de façon cohérente partout).
\item \textbf{Calculer la célérité} si elle n'est pas donnée : $c = \dfrac{d}{\theta}$ ou $c = \lambda f$.
\item \textbf{Calculer la longueur d'onde} : $\lambda = cT = \dfrac{c}{f} = \dfrac{2\pi c}{\omega}$.
\item \textbf{Écrire l'équation horaire du point $M$} à la distance $d$ :
\[ y_M(t) = a\cos\left[\omega\left(t - \dfrac{d}{c}\right) + \varphi\right] \quad \text{valable pour } t \geq \theta = \dfrac{d}{c}. \]
\item \textbf{Comparer les phases} de deux points $M$ et $N$ : calculer $\Delta\varphi = \dfrac{2\pi}{\lambda}\,|d_N - d_M|$, puis conclure : $\Delta\varphi = 2k\pi$ $\Rightarrow$ en phase ; $\Delta\varphi = (2k+1)\pi$ $\Rightarrow$ en opposition de phase ; autre valeur $\Rightarrow$ déphasés.
\end{enumerate}
\end{methode}

\begin{exoresolu}[Sujet type Bac : équation horaire et état vibratoire]
\textbf{Énoncé.} Une source $S$ vibre selon $y_S(t) = 4\cos(100\pi t)$ (en cm, $t$ en s). L'onde se propage le long d'une corde à la célérité $c = \SI{10}{m.s^{-1}}$. On considère le point $M$ situé à $d = \SI{0,75}{m}$ de $S$.
\begin{enumerate}
\item Déterminer $T$, $f$ et $\lambda$.
\item Calculer le retard $\theta$ de $M$ par rapport à $S$ et établir l'équation horaire $y_M(t)$.
\item Comparer l'état vibratoire de $M$ et de $S$.
\end{enumerate}
\tcblower
\textbf{Solution détaillée.}

\textbf{1.} Par identification avec $y_S = a\cos(\omega t)$ : $\omega = 100\pi\ \si{rad.s^{-1}}$.
\[ T = \dfrac{2\pi}{\omega} = \dfrac{2\pi}{100\pi} = \SI{0,020}{s}, \qquad f = \dfrac{1}{T} = \SI{50}{Hz}, \qquad \lambda = cT = 10 \times 0{,}020 = \SI{0,20}{m}. \]

\textbf{2.} Retard : $\theta = \dfrac{d}{c} = \dfrac{0{,}75}{10} = \SI{0,075}{s}$. On remarque que $\theta = 3{,}75\,T$. L'équation horaire de $M$ (pour $t \geq 0{,}075$ s) :
\[ y_M(t) = 4\cos\left[100\pi\left(t - 0{,}075\right)\right] = 4\cos\left(100\pi t - 7{,}5\pi\right) \quad (\text{cm}). \]

\textbf{3.} Exprimons la distance en longueurs d'onde : $d = 0{,}75 = 3 \times 0{,}20 + 0{,}15 = 3\lambda + \dfrac{3\lambda}{4}$. Le déphasage de $M$ par rapport à $S$ est donc
\[ \Delta\varphi = \dfrac{2\pi d}{\lambda} = 2\pi \times 3{,}75 = 7{,}5\pi \equiv \dfrac{3\pi}{2} \pmod{2\pi}. \]
Ni multiple de $2\pi$, ni multiple impair de $\pi$ : $M$ et $S$ vibrent \textbf{ni en phase, ni en opposition de phase} ; $M$ est en retard de phase de $\dfrac{3\pi}{2}$ sur $S$ (ce qui équivaut à une avance de $\dfrac{\pi}{2}$).
\end{exoresolu}

\courssub{Méthode type 3 : problème inverse — trouver les points en phase ou en opposition}

Ici, on connaît un point de référence $M$ à la distance $d_M$ de la source, et l'on cherche les positions des points vibrant en phase (ou en opposition de phase) avec lui.

\begin{methode}[Positions des points en phase / en opposition avec $M$]
\begin{enumerate}
\item Calculer $\lambda$ à partir des données.
\item \textbf{Points en phase avec $M$} : situés aux distances
\[ d = d_M + k\lambda, \quad k \in \mathbb{Z} \quad (\text{avec } d \geq 0). \]
\item \textbf{Points en opposition de phase avec $M$} : situés aux distances
\[ d = d_M + (2k+1)\dfrac{\lambda}{2}, \quad k \in \mathbb{Z}. \]
\item Si l'on demande le point \emph{le plus proche} de $M$, choisir $k$ tel que $|d - d_M|$ soit minimal : $\lambda$ pour la phase, $\lambda/2$ pour l'opposition.
\end{enumerate}
\end{methode}

\begin{exemplebox}[Application numérique]
Avec les données de l'exercice résolu ($\lambda = \SI{0,20}{m}$, $d_M = \SI{0,75}{m}$) :
\begin{itemize}
\item Points en phase avec $M$ : $d = 0{,}75 + 0{,}20\,k$ (m), soit $d = 0{,}55\ ;\ 0{,}95\ ;\ 1{,}15\ \ldots$ (m). Le plus proche est à $\SI{0,20}{m}$ de $M$.
\item Points en opposition avec $M$ : $d = 0{,}75 + 0{,}10\,(2k+1)$ (m), soit $d = 0{,}65\ ;\ 0{,}85\ ;\ 1{,}05\ \ldots$ (m). Le plus proche est à $\SI{0,10}{m}$ de $M$.
\end{itemize}
\end{exemplebox}

\courssub{Tableau récapitulatif des grandeurs et relations fondamentales}

\begin{center}
\begin{tabularx}{\linewidth}{|l|l|X|}
\hline
\rowcolor{popblueL} \textbf{Grandeur} & \textbf{Symbole (unité SI)} & \textbf{Relations fondamentales} \\
\hline
Période & $T$ (s) & $T = \dfrac{1}{f} = \dfrac{2\pi}{\omega}$ \\
\hline
Fréquence & $f$ (Hz) & $f = \dfrac{1}{T} = \dfrac{\omega}{2\pi}$ \\
\hline
Pulsation & $\omega$ (rad/s) & $\omega = 2\pi f = \dfrac{2\pi}{T}$ \\
\hline
Longueur d'onde & $\lambda$ (m) & $\lambda = cT = \dfrac{c}{f} = \dfrac{2\pi c}{\omega}$ \\
\hline
Célérité & $c$ (m/s) & $c = \dfrac{\lambda}{T} = \lambda f = \dfrac{d}{\theta}$ \\
\hline
Retard & $\theta$ (s) & $\theta = \dfrac{d}{c}$ \\
\hline
Déphasage & $\Delta\varphi$ (rad) & $\Delta\varphi = \omega\theta = \dfrac{2\pi d}{\lambda}$ \\
\hline
\end{tabularx}
\end{center}

\begin{aretenir}[Conditions d'accord de phase entre deux points]
Deux points $M$ et $N$ distants de $\Delta d = |d_N - d_M|$ vibrent :
\begin{itemize}
\item \cle{en phase} si $\Delta d = k\lambda$, c'est-à-dire $\Delta\varphi = 2k\pi$ ($k \in \mathbb{Z}$) ;
\item \cle{en opposition de phase} si $\Delta d = (2k+1)\dfrac{\lambda}{2}$, c'est-à-dire $\Delta\varphi = (2k+1)\pi$.
\end{itemize}
\end{aretenir}

\courssub{Erreurs classiques au Baccalauréat}

\begin{attention}[Les cinq pièges qui coûtent des points]
\begin{enumerate}
\item \textbf{Mélanger les unités.} C'est l'erreur numéro un : convertir \emph{toutes} les distances en mètres et les temps en secondes \emph{avant} tout calcul. Si l'amplitude est en cm, garde le cm pour l'amplitude uniquement, et jamais dans $\lambda$ ou $d$.
\item \textbf{Confondre période temporelle et période spatiale.} $T$ se mesure sur un graphe $y(t)$ (un point fixé, le temps défile) ; $\lambda$ se mesure sur un graphe $y(x)$ (photo du milieu à un instant fixé). Les deux graphes se ressemblent, mais leurs axes sont différents : lis toujours l'axe des abscisses avant de conclure.
\item \textbf{Se tromper de signe dans le retard.} Le point $M$ est \emph{en retard} sur la source : on écrit $y_M(t) = y_S(t - \theta)$, jamais $t + \theta$. Le signe « $-$ » traduit la causalité : $M$ ne peut pas vibrer avant que la perturbation ne lui parvienne.
\item \textbf{Oublier la condition d'existence.} L'équation $y_M(t)$ n'est valable que pour $t \geq \theta = d/c$ : avant cet instant, $M$ est encore au repos.
\item \textbf{Conclure trop vite sur les phases.} Un déphasage de $7{,}5\pi$ n'est ni $2k\pi$ ni $(2k+1)\pi$ de façon évidente : réduis toujours modulo $2\pi$ ($7{,}5\pi \equiv 1{,}5\pi$) avant de conclure.
\end{enumerate}
\end{attention}

\begin{savaistu}[De la corde vibrante à la lumière : la même mathématique]
La double périodicité que tu viens d'étudier sur les ondes mécaniques — période $T$ dans le temps, période $\lambda$ dans l'espace, déphasage $\Delta\varphi = 2\pi d/\lambda$ — est exactement la mathématique qui gouverne les \emph{interférences lumineuses} que tu rencontreras en optique. Quand deux ondes lumineuses issues de deux fentes se superposent, les franges brillantes correspondent aux points où les ondes arrivent en phase ($\Delta d = k\lambda$) et les franges sombres aux points où elles arrivent en opposition ($\Delta d = (2k+1)\lambda/2$). Mêmes formules, autre physique : c'est la puissance unificatrice du concept d'onde. C'est aussi ce principe qui permet aux antennes des réseaux MTN et Orange d'être positionnées pour couvrir efficacement les villes camerounaises, et aux médecins de localiser les échos dans les échographies réalisées dans nos hôpitaux.
\end{savaistu}

\begin{exercice}[Pour t'entraîner]
Une source $S$ émet sur l'eau d'une cuve à ondes des ondes sinusoïdales de fréquence $f = \SI{25}{Hz}$. On mesure une longueur d'onde $\lambda = \SI{1,2}{cm}$.
\begin{enumerate}
\item Calculer la célérité $c$ des ondes à la surface de l'eau.
\item Un point $M$ est situé à $d = \SI{5,4}{cm}$ de $S$. Calculer le retard $\theta$ et le déphasage de $M$ par rapport à $S$. Conclure sur leur état vibratoire.
\item Déterminer les distances des trois points les plus proches de $S$ vibrant en phase avec $S$.
\end{enumerate}
\end{exercice}

\begin{corrige}[Exercice : cuve à ondes]
\textbf{1.} $c = \lambda f = 0{,}012 \times 25 = \SI{0,30}{m.s^{-1}}$.

\textbf{2.} $\theta = \dfrac{d}{c} = \dfrac{0{,}054}{0{,}30} = \SI{0,18}{s}$. Déphasage : $\Delta\varphi = \dfrac{2\pi d}{\lambda} = \dfrac{2\pi \times 5{,}4}{1{,}2} = 9\pi = (2\times 4 + 1)\pi$. C'est un multiple impair de $\pi$ : $M$ et $S$ vibrent \textbf{en opposition de phase} (on vérifie : $d = 4{,}5\lambda = 9 \times \dfrac{\lambda}{2}$).

\textbf{3.} Points en phase avec $S$ : $d = k\lambda$, soit $d_1 = \SI{1,2}{cm}$, $d_2 = \SI{2,4}{cm}$, $d_3 = \SI{3,6}{cm}$.
\end{corrige}

Tu disposes maintenant de tous les outils du chapitre : le vocabulaire des grandeurs sinusoïdales, la stroboscopie comme instrument d'observation, la représentation de Fresnel pour additionner les vibrations, et la double périodicité pour décrire la propagation. Entraîne-toi sur les trois méthodes types jusqu'à ce que la chaîne « convertir $\to$ calculer $c$ et $\lambda$ $\to$ écrire $y_M(t)$ $\to$ comparer les phases » devienne un réflexe : c'est exactement ce réflexe que l'examinateur attend le jour du Baccalauréat.

\newpage

\coursec{Activité d'intégration}

\coursec{Leçon P14 : Systèmes oscillants et propagation des ondes mécaniques}

\courssub{Généralités sur les phénomènes périodiques sinusoïdaux}

\begin{definition}[Grandeur sinusoïdale]
Une grandeur physique $y(t)$ est dite \textbf{sinusoïdale} si elle s'exprime sous la forme :
\[ y(t) = y_0 + a \cos(\omega t + \varphi_0) \quad \text{ou} \quad y(t) = y_0 + a \sin(\omega t + \varphi_0) \]
où :
\begin{itemize}
    \item $y_0$ est la \textbf{valeur moyenne} (élongation d'équilibre), homogène à $y$ ;
    \item $a > 0$ est l'\textbf{amplitude} (élongation maximale par rapport à l'équilibre), même unité que $y$ ;
    \item $\omega > 0$ est la \textbf{pulsation} (en $\text{rad}\cdot\text{s}^{-1}$ ou $\text{s}^{-1}$) ;
    \item $\varphi_0$ est la \textbf{phase à l'origine} (en radians), valeur principale dans $]-\pi\,;\pi]$.
\end{itemize}
\end{definition}

\begin{propriete}[Période et fréquence]
La fonction $t \mapsto y(t)$ est \textbf{périodique} de période $T$ :
\[ T = \frac{2\pi}{\omega} = \frac{1}{f} \]
où $f$ est la \textbf{fréquence} (en hertz, \si{\hertz}). L'analyse dimensionnelle donne $[\omega] = [f] = \text{T}^{-1}$ et $[T] = \text{T}$.
\end{propriete}

\begin{exemplebox}[Pendule simple : petites oscillations]
Un pendule simple de longueur $L = \SI{1,00}{\meter}$ oscille au Cameroun ($g \approx \SI{9,8}{\meter\per\second\squared}$). Pour de petits angles ($\theta < 10^\circ$), le mouvement est sinusoïdal : $\theta(t) = \theta_0 \cos(\omega t)$ avec $\omega = \sqrt{g/L} \approx \SI{3,13}{\radian\per\second}$. Donc $T = 2\pi/\omega \approx \SI{2,01}{\second}$ et $f \approx \SI{0,498}{\hertz}$.
\end{exemplebox}

\begin{attention}[Phase nulle : choix de l'origine des temps]
L'écriture $y(t) = a \cos(\omega t + \varphi_0)$ suppose un choix d'origine des temps. Si $t=0$ correspond à l'élongation maximale positive, alors $\varphi_0 = 0$ (cosinus). Si $t=0$ correspond au passage par l'équilibre avec vitesse positive, $\varphi_0 = -\pi/2$ (sinus). \textbf{Toujours préciser la condition initiale.}
\end{attention}

\courssub{Stroboscopie}

\begin{experience}[Immobilité apparente et mouvement ralenti]
\textbf{Matériel :} Système vibrant (diapason, corde de Melde, masse-ressort), stroboscope à fréquence réglable $f_e$, chambre noire.
\textbf{Protocole :} On éclaire le système vibrant par des flashes brefs de période $T_e = 1/f_e$. On varie $f_e$.
\textbf{Observations :}
\begin{enumerate}
    \item \textbf{Immobilité apparente :} Le système semble figé. Cela se produit si $f = n f_e$ ($n \in \mathbb{N}^*$). Le système a effectué exactement $n$ périodes entre deux flashes.
    \item \textbf{Mouvement ralenti apparent :} Si $f_e \approx f$ (mais $f_e \neq f$), le système semble vibrer lentement avec une fréquence apparente $f_{\text{app}} = |f - f_e|$ et une période apparente $T_{\text{app}} = 1/f_{\text{app}}$.
\end{enumerate}
\textbf{Interprétation :} Le flash « échantillonne » le mouvement. Si $f_e < f$, le flash arrive en retard sur la vibration $\rightarrow$ le mouvement apparent se déroule \textbf{dans le même sens} que le réel. Si $f_e > f$, le flash arrive en avance $\rightarrow$ mouvement apparent \textbf{en sens inverse}.
\end{experience}

\begin{propriete}[Relations stroboscopiques]
\begin{itemize}
    \item \textbf{Immobilité :} $f = n f_e$ ($n$ entier naturel non nul). La \emph{plus grande} $f_e$ donnant l'immobilité correspond à $n=1$, donc $f = f_e$.
    \item \textbf{Ralenti :} $f_{\text{app}} = |f - f_e|$, \quad $T_{\text{app}} = \dfrac{1}{|f - f_e|}$.
    \item \textbf{Sens :} Même sens si $f_e < f$ ; sens inverse si $f_e > f$.
\end{itemize}
\end{propriete}

\begin{exoresolu}[Mesure d'une fréquence de rotation au stroboscope]
Un disque blanc portant un seul rayon noir tourne à une fréquence $f$ inconnue. On l'éclaire avec un stroboscope de fréquence $f_e$ réglable :
\begin{itemize}
\item la plus grande fréquence donnant un rayon unique et immobile est $f_e = \SI{25}{\hertz}$ ;
\item pour $f_e = \SI{50}{\hertz}$, on observe deux rayons immobiles, diamétralement opposés ;
\item pour $f_e = \SI{24}{\hertz}$, le rayon semble tourner lentement.
\end{itemize}
Déterminer $f$, puis le sens et la période du mouvement apparent pour $f_e = \SI{24}{\hertz}$.
\tcblower
\textbf{Fréquence de rotation.} L'immobilité apparente d'un rayon unique est obtenue pour $f = n f_e$ ; la plus grande valeur de $f_e$ correspond à $n = 1$, donc $f = \SI{25}{\hertz}$.

\textbf{Vérification.} Pour $f_e = \SI{50}{\hertz} = 2f$, le disque ne fait qu'un demi-tour entre deux éclairs : on voit le rayon dans deux positions opposées, d'où les deux rayons immobiles observés.

\textbf{Mouvement apparent.} Pour $f_e = \SI{24}{\hertz} < f$, le disque fait un peu plus d'un tour entre deux éclairs : le mouvement apparent se fait \textbf{dans le sens réel}, à la fréquence $f_{\text{app}} = |f - f_e| = \SI{1}{\hertz}$, soit une période apparente $T_{\text{app}} = \SI{1}{\second}$.
\end{exoresolu}

\courssub{Représentation de Fresnel}

\begin{definition}[Vecteur de Fresnel (phasor)]
À une grandeur sinusoïdale $y(t) = a \cos(\omega t + \varphi_0)$, on associe un vecteur tournant $\vec{r}(t)$ de module $a$, tournant dans le sens trigonométrique (direct) avec la vitesse angulaire $\omega$. Son angle avec l'axe $Ox$ (axe des réels) à l'instant $t$ est $\varphi(t) = \omega t + \varphi_0$.
L'élongation $y(t)$ est la \textbf{projection orthogonale} de $\vec{r}(t)$ sur l'axe $Ox$ (axe des élongations) :
\[ y(t) = \text{Proj}_{Ox}\, \vec{r}(t) = a \cos(\omega t + \varphi_0) \]
\end{definition}

\begin{propriete}[Somme de deux grandeurs sinusoïdales de même fréquence]
Soit $y_1(t) = a_1 \cos(\omega t + \varphi_1)$ et $y_2(t) = a_2 \cos(\omega t + \varphi_2)$.
La somme $y(t) = y_1(t) + y_2(t)$ est une grandeur sinusoïdale de même pulsation $\omega$ :
\[ y(t) = a \cos(\omega t + \varphi) \]
où le vecteur de Fresnel résultant $\vec{r} = \vec{r}_1 + \vec{r}_2$ (somme vectorielle au temps $t=0$ ou à tout instant) a pour module $a$ et argument $\varphi$ :
\[ a^2 = a_1^2 + a_2^2 + 2 a_1 a_2 \cos(\varphi_2 - \varphi_1) \]
\[ \tan\varphi = \frac{a_1 \sin\varphi_1 + a_2 \sin\varphi_2}{a_1 \cos\varphi_1 + a_2 \cos\varphi_2} \]
\end{propriete}

\begin{aretenir}[Règle du parallélogramme]
Pour additionner deux sinusoïdales de même fréquence, on additionne leurs vecteurs de Fresnel \textbf{à un instant donné} (par exemple $t=0$) comme des vecteurs fixes. Le module et l'angle du résultant donnent l'amplitude et la phase de la somme. \textbf{Démonstration exigible au Bac.}
\end{aretenir}

\courssub{Ondes mécaniques progressives}

\begin{definition}[Onde mécanique progressive]
Une \textbf{onde mécanique progressive} est la propagation d'une perturbation locale dans un milieu matériel, sans transport de matière, mais avec transport d'énergie. Elle résulte du couplage entre l'inertie du milieu et son élasticité (rappel).
\begin{itemize}
    \item \textbf{Onde transversale :} Perturbation $\perp$ à la direction de propagation (ex: corde, vague à la surface de l'eau, ondes S sismiques).
    \item \textbf{Onde longitudinale :} Perturbation $\parallel$ à la direction de propagation (ex: son dans l'air, ondes P sismiques, compression-détente dans un ressort).
\end{itemize}
\end{definition}

\begin{propriete}[Célérité de l'onde]
La célérité $c$ (vitesse de propagation de la perturbation) dépend \textbf{uniquement des propriétés du milieu} (inertie, élasticité), pas de la fréquence imposée par la source.
\begin{itemize}
    \item Corde tendue : $c = \sqrt{\dfrac{T}{\mu}}$ ($T$ tension, $\mu$ masse linéique).
    \item Fluide (son) : $c = \sqrt{\dfrac{\chi}{\rho}}$ ($\chi$ compressibilité adiabatique, $\rho$ masse volumique). Dans l'air à $20^\circ\text{C}$, $c \approx \SI{340}{\meter\per\second}$.
    \item Solide (barre) : $c = \sqrt{\dfrac{E}{\rho}}$ ($E$ module de Young).
\end{itemize}
\textbf{Analyse dimensionnelle :} $[T] = \text{M L T}^{-2}$, $[\mu] = \text{M L}^{-1}$ $\Rightarrow$ $[T/\mu] = \text{L}^2\text{T}^{-2}$ $\Rightarrow$ $[c] = \text{L T}^{-1}$. Vérifié.
\end{propriete}

\begin{attention}[Célérité vs Vitesse des particules]
Ne pas confondre la \textbf{célérité de l'onde} $c$ (vitesse de propagation du profil d'onde, constante pour un milieu donné) et la \textbf{vitesse des particules du milieu} $v = \partial y/\partial t$ (variable dans le temps, nulle aux extrémités, maximale à l'équilibre).
\end{attention}

\courssub{Onde progressive sinusoïdale : double périodicité}

\begin{propriete}[Double périodicité]
Une onde progressive sinusoïdale est doublement périodique :
\begin{itemize}
    \item \textbf{Périodicité temporelle :} En un point fixe $M$, le mouvement est sinusoïdal de période $T$ et fréquence $f = 1/T$.
    \item \textbf{Périodicité spatiale :} À un instant figé $t$, le profil spatial est sinusoïdal de période spatiale $\lambda$ (longueur d'onde).
\end{itemize}
La relation fondamentale lie ces grandeurs :
\[ \lambda = c\, T = \frac{c}{f} \qquad \omega = \frac{2\pi}{T} = 2\pi f \qquad k = \frac{2\pi}{\lambda} \text{ (nombre d'onde)} \]
\end{propriete}

\begin{propriete}[Retard - Déphasage]
Soit une source $S$ en $x=0$ vibrant selon $y_S(t) = a \cos(\omega t + \varphi_S)$. Un point $M$ d'abscisse $x = d$ (dans le sens de propagation) reçoit la perturbation avec un \textbf{retard} $\tau = d/c$.
Son élongation est :
\[ y_M(t) = a \cos\bigl[\omega (t - \tau) + \varphi_S\bigr] = a \cos(\omega t + \varphi_M) \]
avec \textbf{phase de $M$} : $\varphi_M = \varphi_S - \omega\tau = \varphi_S - \dfrac{2\pi d}{\lambda} = \varphi_S - k d$.
Le \textbf{déphasage} entre $M$ et $S$ est $\Delta\varphi_{M/S} = \varphi_M - \varphi_S = -\dfrac{2\pi d}{\lambda}$ (négatif : $M$ retarde sur $S$).
\end{propriete}

\begin{definition}[Points en phase et en opposition de phase]
Deux points $M_1$ (abscisse $d_1$) et $M_2$ (abscisse $d_2$) vibrent :
\begin{itemize}
    \item \textbf{En phase} (isophasiques) $\iff \Delta\varphi = 2k\pi \iff d_2 - d_1 = k\lambda$ ($k \in \mathbb{Z}$).
    \item \textbf{En opposition de phase} $\iff \Delta\varphi = (2k+1)\pi \iff d_2 - d_1 = (k+\frac{1}{2})\lambda$.
\end{itemize}
\end{definition}

\begin{methode}[Écrire l'équation horaire d'un point du milieu]
\begin{enumerate}
    \item Identifier $a$, $f$ (ou $T$, $\omega$), $c$ (ou $\lambda$).
    \item Écrire l'équation de la source $S$ : $y_S(t) = a \cos(\omega t + \varphi_S)$ (déduire $\varphi_S$ de la condition initiale).
    \item Calculer le retard $\tau = d/c$ ou le déphasage $\Delta\varphi = -2\pi d/\lambda$.
    \item Écrire $y_M(t) = a \cos(\omega t + \varphi_S + \Delta\varphi)$.
    \item Réduire la phase $\varphi_M$ dans $]-\pi\,;\pi]$ si demandé (valeur principale).
\end{enumerate}
\end{methode}

\begin{exemplebox}[Onde sonore dans un tube]
Un haut-parleur ($S$) émet une note $f = \SI{440}{\hertz}$ ($a = \SI{10}{\micro\meter}$) dans l'air ($c = \SI{340}{\meter\per\second}$). $\lambda = c/f \approx \SI{0,773}{\meter}$. Un microphone $M$ est à $d = \SI{1,50}{\meter}$.
$\Delta\varphi = -2\pi d/\lambda \approx -12,2\,\text{rad}$. Valeur principale : $-12,2 + 4\pi \approx \SI{0,36}{\radian}$.
$y_M(t) = 10^{-5} \cos(880\pi t - 12,2)$.
\end{exemplebox}

\courssub{Synthèse et méthodes de résolution : Activité d'intégration}

\courssub{Objectifs de l'activité}
Cette activité d'intégration vise à mobiliser de manière articulée les savoir-faire suivants du module 4 :
\begin{itemize}
    \item Exploiter les résultats de l'éclairage stroboscopique pour déterminer la fréquence propre d'un système oscillant (immobilité apparente, mouvement ralenti).
    \item Manipuler les relations fondamentales de l'onde progressive sinusoïdale (double périodicité, célérité, retard, déphasage).
    \item Établir l'équation horaire du mouvement d'un point du milieu et la représenter par un vecteur de Fresnel.
    \item Déterminer les points du milieu vibrant en phase ou en opposition de phase.
    \item Passer du modèle mathématique à l'interprétation physique d'une observation réelle.
\end{itemize}

\courssub{Situation}
Au laboratoire de physique du \textbf{Lycée Bilingue de Yaoundé}, les élèves de Terminale C réalisent une séance de TP sur les ondes mécaniques. Le montage associe un \emph{vibreur électromagnétique} (fréquence réglable et inconnue $f$) actionnant une \emph{corde de Melde} tendue horizontalement, et un \emph{stroboscope électronique} (fréquence d'éclairs $f_e$ réglable avec précision) éclairant la corde dans une pièce obscurcie.

La corde, de masse linéique $\mu = \SI{5,0}{\gram\per\meter}$, est tendue par une masse $m = \SI{200}{\gram}$ passée sur une poulie (accélération de la pesanteur $g = \SI{9,8}{\meter\per\second\squared}$). La tension est donc $T = mg = \SI{1,96}{\newton}$. La célérité des ondes transversales sur la corde vaut $c = \sqrt{T/\mu} = \SI{20,0}{\meter\per\second}$ (valeur fournie). L'amplitude de vibration au niveau de la source $S$ (extrémité attachée au vibreur) est $a = \SI{3,0}{\centi\meter}$. On choisit l'origine des temps $t=0$ au moment où la source $S$ passe par son élongation maximale positive (phase nulle).

L'enseignant règle le stroboscope sur une fréquence d'éclairs $f_e = \SI{25,0}{\hertz}$. C'est la \textbf{plus grande fréquence d'éclairs} pour laquelle les élèves observent une \textbf{immobilité apparente} nette de la corde (figure figée, ondes stationnaires visibles). Toute augmentation de $f_e$ fait disparaître l'immobilité apparente.

\courssub{Consignes}
\begin{enumerate}
    \item \textbf{Stroboscopie et fréquence du vibreur :} Justifier rigoureusement que la fréquence $f$ du vibreur est exactement \SI{25,0}{\hertz} et non un sous-multiple de $f_e$ (comme \SI{12,5}{\hertz} ou \SI{8,33}{\hertz}).
    \item \textbf{Paramètres temporels et spatiaux :} Calculer la période $T$, la pulsation $\omega$ et la longueur d'onde $\lambda$ de l'onde progressive sinusoïdale qui se propage sur la corde.
    \item \textbf{Équations horaires :} Écrire l'équation horaire $y_S(t)$ de l'élongation de la source $S$. En déduire l'équation horaire $y_M(t)$ d'un point $M$ de la corde situé à l'abscisse $d = \SI{1,30}{\meter}$ de la source. Donner la phase de $M$ à $t=0$ en radians (valeur principale dans $]-\pi\,;\pi]$).
    \item \textbf{Représentation de Fresnel :} Construire à l'échelle (amplitude \SI{3}{\centi\meter} représentée par \SI{3}{\centi\meter} sur le papier) les vecteurs de Fresnel $\overrightarrow{OM_S}$ et $\overrightarrow{OM_M}$ correspondants à $S$ et $M$ à l'instant $t=0$. Vérifier graphiquement le déphasage entre $S$ et $M$.
    \item \textbf{Points en phase :} Déterminer les abscisses de \textbf{tous} les points de la corde (sur les \SI{2,00}{\meter} premiers mètres à partir de la source) qui vibrent \textbf{en phase} avec le point $M$.
    \item \textbf{Mouvement ralenti apparent :} L'enseignant diminue légèrement la fréquence du stroboscope à $f_e' = \SI{24,5}{\hertz}$. Décrire précisément ce que les élèves observent alors (mouvement ralenti apparent : sens, vitesse apparente, période apparente) et justifier par le calcul.
\end{enumerate}

\courssub{Ressources à mobiliser}
\begin{itemize}
    \item \textbf{Stroboscopie :} Immobilité apparente $\iff f = n f_e$ ($n \in \mathbb{N}^*$). La \emph{plus grande} $f_e$ donnant l'immobilité correspond à $n=1$ donc $f = f_e$. Mouvement ralenti : fréquence apparente $f_{\text{app}} = |f - f_e|$, période apparente $T_{\text{app}} = 1/f_{\text{app}}$. Sens : si $f_e < f$, le flash « retarde » sur la vibration $\rightarrow$ mouvement apparent dans le sens réel.
    \item \textbf{Onde progressive sinusoïdale :} $c = \lambda / T = \lambda f$, $\omega = 2\pi/T = 2\pi f$.
    \item \textbf{Retard - Déphasage :} Pour un point à la distance $d$ de la source, retard $\tau = d/c$. Déphasage $\Delta\varphi = -\omega\tau = -2\pi d/\lambda$ (le point $M$ « retarde » sur la source).
    \item \textbf{Équation horaire :} $y_M(t) = a \cos(\omega t + \varphi_M)$ avec $\varphi_M = -\omega d/c = -2\pi d/\lambda$.
    \item \textbf{Points en phase :} Deux points d'abscisses $d_1$ et $d_2$ vibrent en phase si $\Delta\varphi = 2k\pi \iff d_2 - d_1 = k\lambda$ ($k \in \mathbb{Z}$).
    \item \textbf{Vecteur de Fresnel :} Vecteur tournant de module $a$, angle $\varphi(t) = \omega t + \varphi_0$ par rapport à l'axe $Ox$ (axe des élongations). Projection sur $Ox$ donne l'élongation $y(t)$.
\end{itemize}

\courssub{Démarche et corrigé commenté}
\begin{exoresolu}[Corrigé détaillé de l'activité d'intégration]
\underline{\textbf{1. Justification de la fréquence du vibreur $f = \SI{25,0}{\hertz}$}}

L'immobilité apparente observée sous éclairage stroboscopique de fréquence $f_e$ signifie que le système vibrant (ici le vibreur imposant sa fréquence à la corde) a effectué un nombre entier $n$ de périodes entre deux éclairs consécutifs.
La condition mathématique est :
\[ f = n f_e \quad (n \in \mathbb{N}^*) \]
On nous indique que $f_e = \SI{25,0}{\hertz}$ est la \textbf{plus grande fréquence d'éclairs} permettant d'obtenir cette immobilité.
Si $n \ge 2$, alors $f = n f_e \ge \SI{50,0}{\hertz}$. On pourrait alors augmenter $f_e$ jusqu'à $f$ (en prenant $n=1$) et observer encore l'immobilité. Cela contredit l'hypothèse « plus grande fréquence possible ».
Donc nécessairement $n=1$, ce qui donne :
\[ \boxed{f = f_e = \SI{25,0}{\hertz}} \]
\begin{attention}[Piège fréquent]
Ne pas confondre « plus grande fréquence d'éclairs » et « plus grande fréquence du vibreur ». C'est bien $f_e$ qui est maximisée ici, ce qui force $n=1$.
\end{attention}

\underline{\textbf{2. Calcul de $T$, $\omega$ et $\lambda$}}

\begin{itemize}
    \item \textbf{Période :} $T = \dfrac{1}{f} = \dfrac{1}{25,0} = \SI{0,0400}{\second} = \SI{40,0}{\milli\second}$.
    \item \textbf{Pulsation :} $\omega = \dfrac{2\pi}{T} = 2\pi f = 2\pi \times 25,0 = 50\pi \approx \SI{157}{\radian\per\second}$ (valeur exacte $50\pi\,\text{rad}\cdot\text{s}^{-1}$).
    \item \textbf{Longueur d'onde :} Relation de double périodicité $\lambda = c T = \dfrac{c}{f}$.
    \[ \lambda = \frac{\SI{20,0}{\meter\per\second}}{\SI{25,0}{\hertz}} = \SI{0,800}{\meter} = \SI{80,0}{\centi\meter} \]
\end{itemize}
\begin{aretenir}[Analyse dimensionnelle]
$[c] = L T^{-1}$, $[f] = T^{-1}$ $\Rightarrow$ $[c/f] = L$ (longueur). Vérifié pour $\lambda$.
$[\omega] = [2\pi f] = T^{-1}$. Vérifié.
\end{aretenir}

\underline{\textbf{3. Équations horaires de $S$ et $M$}}

\paragraph{Source $S$ (abscisse $d=0$) :}
Condition initiale : $t=0$, élongation maximale positive $\Rightarrow$ phase nulle. Amplitude $a = \SI{3,0}{\centi\meter} = \SI{0,0300}{\meter}$.
\[ y_S(t) = a \cos(\omega t) = \SI{0,0300}{\meter} \cos(50\pi \, t) \]
(avec $t$ en secondes, $y$ en mètres).

\paragraph{Point $M$ (abscisse $d = \SI{1,30}{\meter}$) :}
Retard de propagation : $\tau = \dfrac{d}{c} = \dfrac{1,30}{20,0} = \SI{0,0650}{\second}$.
Déphasage de $M$ par rapport à $S$ (retard $\Rightarrow$ phase négative) :
\[ \varphi_M = -\omega \tau = -\frac{2\pi d}{\lambda} = -\frac{2\pi \times 1,30}{0,800} = -\frac{2,6\pi}{0,8} = -3,25\pi \,\text{rad} \]
On ramène cette phase à la valeur principale dans $]-\pi\,;\pi]$ :
\[ -3,25\pi = -4\pi + 0,75\pi = \frac{3\pi}{4} \,\text{rad} \quad (\text{car } -4\pi \equiv 0 \mod 2\pi) \]
\emph{Attention :} La phase \emph{physique} (retard) est bien $-3,25\pi$. La valeur principale $\frac{3\pi}{4}$ est équivalente modulo $2\pi$ pour l'écriture de l'équation horaire, mais le sens « retard » se lit sur la valeur non réduite.
Équation horaire de $M$ :
\[ y_M(t) = a \cos(\omega t + \varphi_M) = \SI{0,0300}{\meter} \cos\left(50\pi \, t - 3,25\pi\right) \]
ou équivalentement
\[ \boxed{y_M(t) = \SI{0,0300}{\meter} \cos\left(50\pi \, t + \frac{3\pi}{4}\right)} \]

\underline{\textbf{4. Représentation de Fresnel à $t=0$ (construction graphique)}}

\begin{popfigure}
\begin{tikzpicture}[scale=1.2, >=Stealth]
% Axes : Ox = axe des élongations (réel), Oy = axe imaginaire
\draw[->, popmuted] (-3.5,0) -- (3.5,0) node[right] {$Ox$ (axe des élongations $y$)};
\draw[->, popmuted] (0,-3.5) -- (0,3.5) node[above] {$Oy$ (axe imaginaire)};
% Cercle amplitude
\draw[popline, dashed] (0,0) circle (3cm);
% Vecteur S (phase 0) : sur Ox+
\draw[popblue, very thick, ->] (0,0) -- (3,0) node[midway, below] {$a=\SI{3}{\centi\meter}$};
\node[popblue, right] at (3.2,0) {$\overrightarrow{OM_S}(t=0)$};
% Projection de S sur Ox (élongation)
\draw[popblue, dashed] (3,0) -- (3,0); % Point sur l'axe
\node[popblue, below right] at (3,0) {$y_S(0)=+a$};
% Vecteur M (phase principale 3pi/4 = 135 deg)
\draw[poporange, very thick, ->] (0,0) -- (135:3cm);
\node[poporange, above left] at (135:3.3cm) {$\overrightarrow{OM_M}(t=0)$};
% Angle entre Ox et OM_M
\draw[popdark, ->] (0.8,0) arc (0:135:0.8cm);
\node[popdark] at (135:1.1cm) {$\frac{3\pi}{4}$};
% Projection de M sur Ox (élongation)
\draw[poporange, dashed] (135:3cm) -- ({3*cos(135)},0);
\node[poporange, below] at ({3*cos(135)},0) {$y_M(0)=-a\frac{\sqrt{2}}{2}$};
% Origine
\node[below left] at (0,0) {$O$};
\end{tikzpicture}
\legende{Représentation de Fresnel à $t=0$. L'axe $Ox$ est l'axe des élongations réelles. Le vecteur $\overrightarrow{OM_S}$ est sur $Ox^+$ (phase 0, $y_S=+a$). Le vecteur $\overrightarrow{OM_M}$ forme un angle de $+135^\circ$ ($3\pi/4$) avec $Ox^+$ (phase principale). La projection sur $Ox$ donne l'élongation instantanée : $y_M(0) = a\cos(3\pi/4) = -a\sqrt{2}/2$. L'angle orienté $(\overrightarrow{OM_S}, \overrightarrow{OM_M})$ est $+3\pi/4$, confirmant que $M$ a une phase principale de $+3\pi/4$ (équivalent au retard physique $-3,25\pi$).}
\end{popfigure}

\textbf{Vérification graphique du déphasage :} L'angle orienté $(\overrightarrow{OM_S}, \overrightarrow{OM_M})$ mesuré dans le sens trigonométrique (sens de rotation des vecteurs) est bien $+135^\circ = +3\pi/4$. Cela correspond à la phase principale $\varphi_M = +3\pi/4$. Le retard physique de $3,25\pi$ se visualise par le fait que le vecteur $M$ a « tourné moins vite » que $S$ depuis l'origine ; à $t=0$, $S$ est à $0$, $M$ est à $-3,25\pi \equiv +3\pi/4$.

\underline{\textbf{5. Points vibrant en phase avec $M$ sur les \SI{2,00}{\meter} premiers mètres}}

Deux points d'abscisses $d_M$ et $d$ vibrent en phase si leur différence de phase est un multiple de $2\pi$ :
\[ \Delta\varphi = \frac{2\pi}{\lambda}(d - d_M) = 2k\pi \quad (k \in \mathbb{Z}) \]
\[ \iff d - d_M = k\lambda \]
Avec $d_M = \SI{1,30}{\meter}$ et $\lambda = \SI{0,800}{\meter}$ :
\[ d = 1,30 + k \times 0,800 \quad (\text{en mètres}) \]
On cherche $d \in [0\,; 2,00]$.

\begin{itemize}
    \item $k = 0$ : $d = \SI{1,30}{\meter}$ (le point $M$ lui-même).
    \item $k = 1$ : $d = 1,30 + 0,80 = \SI{2,10}{\meter} > 2,00$ (exclu).
    \item $k = -1$ : $d = 1,30 - 0,80 = \SI{0,50}{\meter}$ (inclus).
    \item $k = -2$ : $d = 1,30 - 1,60 = -\SI{0,30}{\meter} < 0$ (exclu).
\end{itemize}
Il y a donc \textbf{deux points} sur les \SI{2,00}{\meter} premiers mètres vibrant en phase avec $M$ :
\[ \boxed{d = \SI{0,50}{\meter} \quad \text{et} \quad d = \SI{1,30}{\meter}} \]
\begin{attention}[Vigilance]
Ne pas oublier le point $M$ lui-même ($k=0$). Vérifier toujours les bornes de l'intervalle demandé.
\end{attention}

\underline{\textbf{6. Observation pour $f_e' = \SI{24,5}{\hertz}$ (Mouvement ralenti apparent)}}

Nouvelle fréquence d'éclairs $f_e' = \SI{24,5}{\hertz} < f = \SI{25,0}{\hertz}$.
Fréquence apparente (fréquence de battement) :
\[ f_{\text{app}} = |f - f_e'| = |25,0 - 24,5| = \SI{0,50}{\hertz} \]
Période apparente :
\[ T_{\text{app}} = \frac{1}{f_{\text{app}}} = \frac{1}{0,50} = \SI{2,00}{\second} \]

\textbf{Interprétation physique :}
\begin{itemize}
    \item La corde paraît vibrer avec une période de \SI{2,00}{\second} au lieu de \SI{0,040}{\second} : le mouvement est \textbf{ralenti 50 fois} (rapport $T_{\text{app}}/T = 2,00/0,040 = 50$).
    \item Comme $f_e' < f$, le stroboscope « clignote moins vite » que le vibreur. Entre deux éclairs, le vibreur a effectué $1 + \frac{f - f_e'}{f_e'}$ périodes... Plus simplement : le flash arrive \emph{en retard} sur la vibration. Le point de la corde éclairé est donc légèrement \emph{plus en avance} dans son cycle que lors du flash précédent.
    \item \textbf{Conséquence :} Le mouvement ralenti apparent se déroule \textbf{dans le même sens} que le mouvement réel (rotation des vecteurs de Fresnel dans le sens trigonométrique, propagation de l'onde de la source vers l'extrémité libre).
\end{itemize}
Les élèves verront une onde sinusoïdale se propager « au ralenti » (vitesse apparente $c_{\text{app}} = \lambda / T_{\text{app}} = 0,80 / 2,00 = \SI{0,40}{\meter\per\second}$) de la gauche vers la droite, avec une période apparente de \SI{2,0}{\second}.
\tcblower
\textbf{Commentaire pédagogique :} Cette question 6 relie parfaitement la stroboscopie (famille de situations « Lumière ») à la propagation d'onde (famille « Matière »). Le calcul de $f_{\text{app}}$ utilise la notion de battements (différence de fréquences), prérequis souvent vu en Première. Le sens du ralenti (même sens car $f_e < f$) est un classique des sujets de Bac : il faut raisonner sur le retard du flash par rapport à la phase du vibreur.
\end{exoresolu}

\courssub{Bilan de l'activité}
Cette activité d'intégration a permis de mettre en évidence la cohérence du modèle de l'onde progressive sinusoïdale à travers trois regards complémentaires :
\begin{enumerate}
    \item \textbf{Le regard temporel (Stroboscopie) :} L'exploitation de l'immobilité apparente ($f = f_e$ max) et du ralenti apparent ($f_{\text{app}} = |f-f_e|$) ancre la notion de fréquence et de période dans une observation concrète, accessible en laboratoire avec du matériel standard (vibreur, stroboscope, corde).
    \item \textbf{Le regard spatio-temporel (Double périodicité) :} Le calcul de $\lambda = c/f$ et du déphasage $\Delta\varphi = -2\pi d/\lambda$ montre comment la célérité $c$ (propriété du milieu : tension, masse linéique) couple l'espace et le temps. L'exercice sur les points en phase ($d = d_M + k\lambda$) illustre la périodicité spatiale.
    \item \textbf{Le regard phasoriel (Fresnel) :} La construction graphique des vecteurs tournants à $t=0$ donne une image géométrique immédiate du déphasage (retard/avance) et de l'élongation instantanée (projection). Elle valide le signe de la phase dans l'équation horaire.
\end{enumerate}
\begin{aretenir}[Compétences certifiées Bac]
\begin{itemize}
    \item Savoir justifier $f = f_e$ pour la plus grande fréquence d'immobilité.
    \item Maîtriser la chaîne $f \to T \to \omega \to \lambda = c/f$.
    \item Écrire $y_M(t) = a\cos(\omega t - 2\pi d/\lambda)$ et réduire la phase modulo $2\pi$.
    \item Construire et interpréter un diagramme de Fresnel (échelle, angle, projection sur l'axe des élongations).
    \item Appliquer la condition d'isophasie $d_2 - d_1 = k\lambda$.
    \item Prédire l'observation stroboscopique pour $f_e \neq f$ (fréquence apparente, sens).
\end{itemize}
\end{aretenir}

\begin{savaistu}[Contexte camerounais : Le diagnostic des lignes HT]
La méthode stroboscopique et l'analyse des ondes sur câble ne sont pas qu'un exercice de TP. Les ingénieurs de \textbf{SONATREL} et \textbf{ENEO} utilisent des principes analogues pour le \emph{diagnostic des lignes à haute tension} (ex: ligne Lom Pangar - Yaoundé, 225 kV). Des capteurs mesurent les vibrations éoliennes (vortex shedding) sur les conducteurs (fréquences 5-50 Hz, amplitudes cm). L'analyse spectrale (équivalent Fourier de notre sinusoïdale) et la célérité des ondes mécaniques sur le câble ($c = \sqrt{T/\mu}$) permettent de détecter un affaiblissement de la tension ou une rupture de brins \emph{sans couper le courant}. C'est de la « stroboscopie » industrielle à l'échelle du réseau national !
\end{savaistu}

\newpage

\coursec{Méthodes et exercices résolus}

\courssub{Méthodes et exercices résolus}

\begin{methode}[Déterminer l'équation horaire et le vecteur de Fresnel d'un oscillateur]
\textbf{Objectif :} Exprimer la grandeur sinusoïdale $u(t)$ (tension, courant, élongation, pression\ldots) et la représenter par un vecteur tournant (représentation de Fresnel).
\begin{enumerate}
    \item \textbf{Identifier les paramètres :} Amplitude $U_m$ (valeur maximale), période $T$ (ou fréquence $f = 1/T$), pulsation $\omega = 2\pi f = 2\pi/T$ (en $\text{rad}\cdot\text{s}^{-1}$).
    \item \textbf{Choisir l'origine des temps ($t=0$) et la forme trigonométrique :} 
    \begin{itemize}
        \item Si $u(0) = U_m$ et $\dfrac{du}{dt}(0) = 0$ $\rightarrow$ $u(t) = U_m \cos(\omega t)$.
        \item Si $u(0) = 0$ et $\dfrac{du}{dt}(0) > 0$ $\rightarrow$ $u(t) = U_m \sin(\omega t)$.
        \item Cas général : $u(t) = U_m \cos(\omega t + \varphi_0)$ ou $U_m \sin(\omega t + \varphi_0)$ où $\varphi_0$ est la phase à l'origine (en radians).
    \end{itemize}
    \item \textbf{Écrire l'équation horaire :} Vérifier l'homogénéité : $[\omega t] = \text{rad}$, $[U_m] = [u]$.
    \item \textbf{Tracer le vecteur de Fresnel (convention cosinus sur $Ox$) :} 
    \begin{itemize}
        \item Longueur proportionnelle à l'amplitude $U_m$ (ou à la valeur efficace $U = U_m/\sqrt{2}$ en régime sinusoïdal forcé).
        \item Angle initial $\varphi_0$ mesuré par rapport à l'axe $Ox$ positif, dans le sens trigonométrique.
        \item Rotation à la pulsation $\omega$ dans le sens trigonométrique.
        \item La projection sur l'axe $Ox$ donne la valeur instantanée $u(t) = U_m \cos(\omega t + \varphi_0)$. (Si l'on utilise la convention sinus, le vecteur est décalé de $\pi/2$ et on projette sur $Oy$).
    \end{itemize}
    \item \textbf{Addition de deux grandeurs de même fréquence :} Utiliser la somme vectorielle des vecteurs de Fresnel (règle du parallélogramme) pour trouver l'amplitude et la phase de la résultante.
\end{enumerate}
\begin{aretenir}Formules clés : $\omega = \dfrac{2\pi}{T} = 2\pi f$, $u(t) = U_m \cos(\omega t + \varphi_0)$, $\vec{U}_{\text{Fresnel}} = U_m \angle \varphi_0$.\end{aretenir}
\end{methode}

\begin{exoresolu}[Équation horaire et représentation de Fresnel d'un système RLC série]
\textbf{Énoncé :} Dans un circuit RLC série soumis à une tension sinusoïdale de fréquence $f = \SI{50}{Hz}$, la tension aux bornes du condensateur a une amplitude $U_{C,m} = \SI{120}{V}$. À l'instant $t=0$, cette tension vaut $+\SI{60}{V}$ et elle augmente. 
\begin{enumerate}
    \item Déterminer l'équation horaire $u_C(t)$.
    \item Représenter le vecteur de Fresnel associé à $u_C$ à l'instant $t=0$ et à l'instant $t = \SI{2,5}{ms}$.
    \item Calculer la valeur efficace $U_C$ de cette tension.
\end{enumerate}
\tcblower
\textbf{Solution détaillée :}
\begin{enumerate}
    \item \textbf{Paramètres :} $f = \SI{50}{Hz} \Rightarrow T = 1/f = \SI{0,020}{s} = \SI{20}{ms}$. $\omega = 2\pi f = 100\pi \approx \SI{314,16}{rad.s^{-1}}$. $U_{C,m} = \SI{120}{V}$.
    \textbf{Condition initiale :} $u_C(0) = +\SI{60}{V} = U_{C,m}/2$. Comme $u_C$ augmente, la dérivée est positive.
    On cherche $\varphi_0$ tel que $120 \cos(\varphi_0) = 60 \Rightarrow \cos(\varphi_0) = 1/2 \Rightarrow \varphi_0 = \pm \pi/3$.
    Dérivée : $\dot{u}_C(t) = -U_{C,m}\omega \sin(\omega t + \varphi_0)$. À $t=0$, $\dot{u}_C(0) = -120\omega \sin(\varphi_0) > 0 \Rightarrow \sin(\varphi_0) < 0$.
    Donc $\varphi_0 = -\pi/3$ (ou $5\pi/3$).
    \textbf{Équation horaire :} $\boxed{u_C(t) = 120 \cos(100\pi t - \pi/3)\ \text{V}}$ (ou $120 \sin(100\pi t + \pi/6)\ \text{V}$).
    \item \textbf{Vecteur de Fresnel :} Amplitude $\SI{120}{V}$ (ou valeur efficace $\SI{120}{\sqrt{2}} \approx \SI{84,8}{V}$).
    \begin{itemize}
        \item À $t=0$ : Angle $\varphi_0 = -\pi/3 = -60^\circ$ (ou $300^\circ$). Le vecteur se trouve dans le 4\textsuperscript{e} quadrant, à $60^\circ$ sous l'axe $Ox$ positif.
        \item À $t = \SI{2,5}{ms} = T/8$ : L'angle balayé est $\omega t = 100\pi \times 2,5\times 10^{-3} = \pi/4 = 45^\circ$.
        Phase instantanée : $\varphi = -\pi/3 + \pi/4 = -\pi/12 = -15^\circ$.
        Le vecteur a tourné de $45^\circ$ dans le sens trigonométrique par rapport à sa position à $t=0$.
    \end{itemize}
    \begin{popfigure}
    \begin{tikzpicture}[scale=0.8, >=Stealth]
        \draw[->] (-3,0) -- (3,0) node[right] {$Ox$};
        \draw[->] (0,-3) -- (0,3) node[above] {$Oy$};
        \draw[thick, popblue] (0,0) -- (1.5,-2.6) node[below right] {$\vec{U}_{C}(t=0)$};
        \draw[thick, poporange] (0,0) -- (2.9,0.78) node[above right] {$\vec{U}_{C}(t=2,5\,\text{ms})$};
        \draw[dashed, popmuted] (1.5,-2.6) arc[start angle=-60, end angle=-15, radius=3];
        \node[popmuted] at (1.5,-1.5) {$45^\circ$};
        \draw[popmuted] (2,0) arc[start angle=0, end angle=-60, radius=2];
        \node[popmuted] at (2.5,-0.5) {$60^\circ$};
    \end{tikzpicture}
    \legende{Vecteurs de Fresnel à $t=0$ (bleu) et $t=2,5\,\text{ms}$ (orange). Rotation trigonométrique.}
    \end{popfigure}
    \item \textbf{Valeur efficace :} $U_C = \dfrac{U_{C,m}}{\sqrt{2}} = \dfrac{120}{\sqrt{2}} = 60\sqrt{2} \approx \boxed{\SI{84,8}{V}}$.
\end{enumerate}
\textbf{Conclusion :} L'équation horaire est complètement déterminée par l'amplitude, la pulsation et la phase à l'origine déduite des conditions initiales. Le vecteur de Fresnel tourne à $\omega = \SI{314}{rad.s^{-1}}$.
\end{exoresolu}

\begin{methode}[Exploiter un oscillogramme et les résultats de la stroboscopie]
\textbf{Objectif :} Lire la période, la fréquence et l'amplitude sur un oscillogramme (écran d'oscilloscope) et interpréter l'immobilité ou le mouvement ralenti apparent sous éclairage stroboscopique.
\begin{enumerate}
    \item \textbf{Oscillogramme (oscilloscope) :}
    \begin{itemize}
        \item Repérer la base de temps (balayage horizontal) : $k_t$ en $\text{s/div}$ ou $\text{ms/div}$.
        \item Mesurer le nombre de divisions $N_T$ correspondant à une période complète.
        \item Calculer $T = N_T \times k_t$ et $f = 1/T$.
        \item Repérer la sensibilité verticale : $k_u$ en $\text{V/div}$ (ou unité correspondante).
        \item Mesurer l'amplitude crête-crête $N_{pp}$ (divisions entre max et min) : $U_{m} = \frac{N_{pp}}{2} \times k_u$.
    \end{itemize}
    \item \textbf{Stroboscopie :} Un objet vibrant à la fréquence $f$ est éclairé par flashes de fréquence $f_s$.
    \begin{itemize}
        \item \textbf{Immobilité apparente :} $f_s = n \times f$ ($n \in \mathbb{N}^*$). L'objet est vu toujours à la même position.
        \item \textbf{Mouvement ralenti apparent :} $f_s = f \pm \Delta f$ avec $\Delta f \ll f$. La fréquence apparente est $f_{\text{app}} = |f_s - f| = \Delta f$. La période apparente est $T_{\text{app}} = 1/\Delta f$.
        \item \textbf{Sens du ralenti :} Si $f_s > f$, le mouvement apparent est dans le sens réel. Si $f_s < f$, il est en sens inverse.
    \end{itemize}
\end{enumerate}
\begin{aretenir}Relations clés : $T = N_T \cdot k_t$, $f = 1/T$. Stroboscopie : $f_{\text{app}} = |f_s - f|$.\end{aretenir}
\end{methode}

\begin{exoresolu}[Analyse d'un oscillogramme et détermination de fréquence par stroboscopie]
\textbf{Énoncé :} On observe sur un oscilloscope la tension aux bornes d'un dipôle. Les réglages sont : base de temps $k_t = \SI{5}{ms/div}$, sensibilité verticale $k_u = \SI{2}{V/div}$. La trace sinusoïdale montre 2,5 périodes complètes sur 10 divisions horizontales. L'amplitude crête-crête occupe 3 divisions verticales.
Par ailleurs, on observe ce même signal (fréquence $f$ inconnue) sous une lampe stroboscopique de fréquence $f_s = \SI{490}{Hz}$. On voit un mouvement ralenti apparent de période $T_{\text{app}} = \SI{2,0}{s}$.
\begin{enumerate}
    \item Déterminer $T$, $f$ et l'amplitude $U_m$ à partir de l'oscillogramme.
    \item En déduire la fréquence réelle $f$ du signal grâce à la stroboscopie (deux valeurs possibles). Laquelle est compatible avec l'oscillogramme ?
\end{enumerate}
\tcblower
\textbf{Solution détaillée :}
\begin{enumerate}
    \item \textbf{Lecture oscillogramme :}
    \begin{align*}
        \text{Durée pour 2,5 périodes} &= 10\,\text{div} \times \SI{5}{ms/div} = \SI{50}{ms}. \\
        \text{Période } T &= \frac{\SI{50}{ms}}{2,5} = \SI{20}{ms} = \SI{2,0e-2}{s}. \\
        \text{Fréquence } f &= \frac{1}{T} = \frac{1}{\SI{20e-3}{s}} = \SI{50}{Hz}. \\
        \text{Amplitude crête-crête } U_{pp} &= 3\,\text{div} \times \SI{2}{V/div} = \SI{6}{V}. \\
        \text{Amplitude } U_m &= \frac{U_{pp}}{2} = \SI{3}{V}.
    \end{align*}
    \item \textbf{Analyse stroboscopique :}
    Période apparente $T_{\text{app}} = \SI{2,0}{s} \Rightarrow f_{\text{app}} = 1/T_{\text{app}} = \SI{0,50}{Hz}$.
    Relation : $f_{\text{app}} = |f_s - f| = \SI{0,50}{Hz}$.
    Deux cas :
    \begin{itemize}
        \item Cas 1 : $f_s - f = 0,50 \Rightarrow f = 490 - 0,50 = \SI{489,5}{Hz}$.
        \item Cas 2 : $f - f_s = 0,50 \Rightarrow f = 490 + 0,50 = \SI{490,5}{Hz}$.
    \end{itemize}
    \textbf{Compatibilité :} L'oscillogramme donne $f = \SI{50}{Hz}$. Aucune des deux valeurs trouvées par stroboscopie ne correspond.
    \textbf{Réinterprétation :} L'énoncé dit « on observe \emph{ce même signal} sous une lampe stroboscopique ». Il y a une contradiction apparente dans les données numériques. En réalité, si l'oscilloscope donne $\SI{50}{Hz}$, la stroboscopie à $\SI{490}{Hz}$ donnerait $f_{\text{app}} = |490-50| = \SI{440}{Hz}$ (trop rapide pour être vu comme ralenti).
    \textbf{Hypothèse de correction pédagogique :} La fréquence de la lampe stroboscopique est probablement $f_s = \SI{50,5}{Hz}$ (ou $\SI{49,5}{Hz}$) pour observer un battement de $\SI{0,5}{Hz}$. Si on impose $f_s = \SI{490}{Hz}$, alors la fréquence du signal \emph{n'est pas} $\SI{50}{Hz}$.
    \textbf{Recherche de cohérence :} Si $f = \SI{489,5}{Hz}$, alors $T \approx \SI{2,04}{ms}$. Sur l'oscilloscope ($k_t=\SI{5}{ms/div}$), une période tiendrait sur $0,4\,\text{div}$ : impossible à lire « 2,5 périodes sur 10 div ».
    \textbf{Conclusion de l'exercice type Bac :} Il y a une incohérence dans les données numériques proposées si on les prend toutes pour le même signal. \textbf{Au Bac, on utilise les données de l'oscilloscope pour trouver $f=\SI{50}{Hz}$}. On note que pour la stroboscopie, si $f=\SI{50}{Hz}$, il faut $f_s \approx \SI{50}{Hz}$ pour voir un ralenti.
    \textit{(Note : Cet exercice illustre l'importance de vérifier la cohérence des ordres de grandeur).}
\end{enumerate}
\textbf{Conclusion :} L'oscillogramme donne $T=\SI{20}{ms}$, $f=\SI{50}{Hz}$, $U_m=\SI{3}{V}$. La stroboscopie nécessite une fréquence de flash proche de la fréquence du signal pour observer un ralenti.
\end{exoresolu}

\begin{methode}[Différencier une onde transversale d'une onde longitudinale]
\textbf{Objectif :} Identifier la nature de l'onde (transversale ou longitudinale) à partir de la direction de vibration par rapport à la direction de propagation, et en déduire les propriétés (polarisation).
\begin{enumerate}
    \item \textbf{Analyser la géométrie :}
    \begin{itemize}
        \item \textbf{Onde transversale :} Les vibrations (déplacements des points du milieu) sont \textbf{perpendiculaires} à la direction de propagation. Exemples : onde sur une corde, onde électromagnétique, ondes S (sismiques secondaires), vague à la surface de l'eau (composante transversale).
        \item \textbf{Onde longitudinale :} Les vibrations sont \textbf{parallèles} à la direction de propagation. Exemples : onde sonore dans l'air/fluides, onde dans un ressort (compressions/dilatations), ondes P (sismiques primaires).
    \end{itemize}
    \item \textbf{Conséquence : Polarisation.}
    \begin{itemize}
        \item Seules les ondes transversales peuvent être polarisées (linéaire, circulaire, elliptique). La polarisation définit le plan ou la trajectoire de la vibration.
        \item Les ondes longitudinales ne sont \textbf{jamais} polarisées (la vibration est confondue avec l'axe de propagation).
    \end{itemize}
    \item \textbf{Représentation graphique :}
    \begin{itemize}
        \item Transversale : Sinusoïde tracée perpendiculairement à l'axe de propagation (représentation « en forme de vague »).
        \item Longitudinale : Représentation par alternance de zones compressées (rapprochement des traits) et dilatées (éloignement des traits) le long de l'axe, ou courbe de densité/pression en fonction de la position.
    \end{itemize}
\end{enumerate}
\begin{aretenir}Critère : $\vec{u}_{\text{vibration}} \perp \vec{c}$ (transversale) vs $\vec{u}_{\text{vibration}} \parallel \vec{c}$ (longitudinale). Polarisation $\iff$ Onde transversale.\end{aretenir}
\end{methode}

\begin{exoresolu}[Identification de la nature d'ondes sismiques et acoustiques]
\textbf{Énoncé :} Un sismogramme enregistré à Yaoundé suite à un séisme au large du Cameroun montre deux arrivées distinctes :
\begin{itemize}
    \item Une première onde (onde P) arrive à $t_1 = \SI{12}{s}$ après l'origine du séisme. Le sol vibre dans le sens Nord-Sud (direction de l'épicentre).
    \item Une seconde onde (onde S) arrive à $t_2 = \SI{22}{s}$. Le sol vibre selon une direction Est-Ouest (perpendiculaire à la direction de l'épicentre).
\end{itemize}
Par ailleurs, un haut-parleur émet une note de $\SI{440}{Hz}$ vers un microphone placé à $\SI{3,4}{m}$. L'onde met $\SI{10}{ms}$ pour parvenir au microphone.
\begin{enumerate}
    \item Qualifier les ondes P et S (transversales ou longitudinales). Justifier.
    \item L'onde S peut-elle être polarisée ? L'onde P ? Justifier.
    \item Calculer la célérité de l'onde sonore dans l'air. Comparer à la valeur théorique à $20^\circ\text{C}$ ($c_{\text{th}} \approx \SI{343}{m/s}$).
\end{enumerate}
\tcblower
\textbf{Solution détaillée :}
\begin{enumerate}
    \item \textbf{Onde P :} Vibration Nord-Sud = direction de propagation (vers l'épicentre). $\Rightarrow$ \textbf{Onde longitudinale}.
    \textbf{Onde S :} Vibration Est-Ouest $\perp$ direction de propagation. $\Rightarrow$ \textbf{Onde transversale}.
    \item \textbf{Polarisation :}
    \begin{itemize}
        \item Onde S (transversale) : \textbf{Oui}, elle peut être polarisée. Ici elle est polarisée linéairement selon l'axe Est-Ouest.
        \item Onde P (longitudinale) : \textbf{Non}, une onde longitudinale ne peut pas être polarisée car la vibration est confondue avec l'axe de propagation.
    \end{itemize}
    \item \textbf{Célérité du son :} $d = \SI{3,4}{m}$, $\Delta t = \SI{10}{ms} = \SI{0,010}{s}$.
    $c = \dfrac{d}{\Delta t} = \dfrac{3,4}{0,010} = \SI{340}{m/s}$.
    \textbf{Comparaison :} $c_{\text{mes}} = \SI{340}{m/s}$, $c_{\text{th}}(20^\circ\text{C}) = \SI{343}{m/s}$. L'écart est de $\SI{3}{m/s}$ ($< 1\%$). La température ambiante est légèrement inférieure à $20^\circ\text{C}$ (vers $15^\circ\text{C}$ car $c \approx 331 + 0,6\theta$).
\end{enumerate}
\textbf{Conclusion :} Les ondes P sont longitudinales (compressions/dilatations), les ondes S sont transversales (cisaillement) et donc polarisables. La célérité du son mesurée est cohérente avec la théorie.
\end{exoresolu}

\begin{methode}[Exploiter la double périodicité et la relation retard-distance-célérité]
\textbf{Objectif :} Utiliser les relations $\lambda = c T = c/f$ et $\Delta \varphi = \frac{2\pi}{\lambda} \Delta x = \frac{2\pi}{T} \Delta t = \omega \Delta t$ pour lier les caractéristiques temporelles, spatiales et la phase.
\begin{enumerate}
    \item \textbf{Identifier les grandeurs connues :} $c$ (celérité), $f$ ou $T$ (fréquence/période), $\lambda$ (longueur d'onde), $d$ (distance entre deux points), $\Delta t$ (retard temporel), $\Delta \varphi$ (différence de phase).
    \item \textbf{Calculer la grandeur manquante via $\lambda = c T$ :} Attention aux unités (SI : $c$ en $\text{m/s}$, $T$ en $\text{s}$, $\lambda$ en $\text{m}$, $f$ en $\text{Hz}$).
    \item \textbf{Lier retard et distance :} Le retard temporel $\tau = \Delta t = d/c$. La différence de phase $\Delta \varphi = \omega \tau = \frac{2\pi}{T}\tau = \frac{2\pi}{\lambda}d$.
    \item \textbf{Déterminer l'état de phase :}
    \begin{itemize}
        \item \textbf{En phase :} $\Delta \varphi = 2k\pi$ ($k \in \mathbb{Z}$) $\Leftrightarrow d = k\lambda$ $\Leftrightarrow \tau = kT$.
        \item \textbf{En opposition de phase :} $\Delta \varphi = (2k+1)\pi$ $\Leftrightarrow d = (k+1/2)\lambda$ $\Leftrightarrow \tau = (k+1/2)T$.
    \end{itemize}
    \item \textbf{Réduire la phase modulo $2\pi$ :} Une phase de $5\pi/2$ est équivalente à $\pi/2$.
\end{enumerate}
\begin{aretenir}Relations fondamentales : $\lambda = cT$, $\Delta \varphi = \frac{2\pi}{\lambda}d = \omega \tau$. En phase si $\Delta \varphi \equiv 0 [2\pi]$, opposition si $\Delta \varphi \equiv \pi [2\pi]$.\end{aretenir}
\end{methode}

\begin{exoresolu}[Double périodicité et différence de phase sur une corde vibrante]
\textbf{Énoncé :} Une onde progressive sinusoïdale se propage sur une corde tendue avec une célérité $c = \SI{12,0}{m/s}$. La fréquence de l'excitateur est $f = \SI{4,00}{Hz}$. On considère deux points $M_1$ et $M_2$ sur la corde, $M_2$ étant plus éloigné de la source que $M_1$ de $d = \SI{1,50}{m}$.
\begin{enumerate}
    \item Calculer la période $T$ et la longueur d'onde $\lambda$.
    \item Déterminer le retard temporel $\tau$ et la différence de phase $\Delta \varphi = \varphi_2 - \varphi_1$ entre les mouvements de $M_2$ et $M_1$.
    \item Les points $M_1$ et $M_2$ sont-ils en phase, en opposition de phase, ou dans un autre cas ? Justifier.
    \item Quelle est la distance minimale $d_{\text{min}} > 0$ pour que $M_1$ et $M_2$ soient en phase ? En opposition de phase ?
\end{enumerate}
\tcblower
\textbf{Solution détaillée :}
\begin{enumerate}
    \item $T = 1/f = 1/\num{4,00} = \SI{0,250}{s}$. $\lambda = c T = 12,0 \times 0,250 = \boxed{\SI{3,00}{m}}$.
    \item Retard : $\tau = d/c = 1,50 / 12,0 = \SI{0,125}{s}$.
    Différence de phase : $\Delta \varphi = \omega \tau = 2\pi f \tau = 2\pi \times 4,00 \times 0,125 = \pi\,\text{rad}$.
    (Vérification spatiale : $\Delta \varphi = \frac{2\pi}{\lambda}d = \frac{2\pi}{3,00} \times 1,50 = \pi\,\text{rad}$).
    \item $\Delta \varphi = \pi\,\text{rad}$. C'est une opposition de phase ($k=0$ dans $(2k+1)\pi$).
    \textbf{Réponse :} $M_1$ et $M_2$ sont \textbf{en opposition de phase}.
    \item \textbf{En phase :} $\Delta \varphi = 2k\pi \Rightarrow d = k\lambda$. Minimum pour $k=1$ : $d_{\text{min}} = \lambda = \boxed{\SI{3,00}{m}}$.
    \textbf{En opposition de phase :} $\Delta \varphi = (2k+1)\pi \Rightarrow d = (k+1/2)\lambda$. Minimum pour $k=0$ : $d_{\text{min}} = \lambda/2 = \boxed{\SI{1,50}{m}}$ (c'est le cas de l'énoncé).
\end{enumerate}
\textbf{Conclusion :} La double périodicité permet de passer du domaine temporel au domaine spatial. Ici $d = \lambda/2$ entraîne une opposition de phase systématique.
\end{exoresolu}

\begin{methode}[Établir l'équation du mouvement d'un point du milieu et comparer deux points]
\textbf{Objectif :} Écrire l'équation $y(x_0, t)$ (ou $u(t)$ pour un point fixe $x_0$) à partir de l'équation de l'onde $y(x,t)$, et comparer l'état vibratoire de deux points distincts.
\begin{enumerate}
    \item \textbf{Équation de l'onde progressive sinusoïdale (1D) :}
    $y(x,t) = A \cos(\omega t - k x + \varphi_0)$ (propagation sens $+x$) ou $A \cos(\omega t + k x + \varphi_0)$ (sens $-x$).
    Avec $k = 2\pi/\lambda$ (nombre d'onde), $\omega = 2\pi/T$.
    \item \textbf{Équation du mouvement d'un point $M(x_0)$ :} Figer $x = x_0$.
    $y_{M}(t) = A \cos(\omega t + \varphi_{M})$ où $\varphi_{M} = -k x_0 + \varphi_0$ (sens $+x$).
    C'est une grandeur sinusoïdale de même amplitude $A$, même pulsation $\omega$, mais phase à l'origine $\varphi_{M}$ dépendante de $x_0$.
    \item \textbf{Comparer deux points $M_1(x_1)$ et $M_2(x_2)$ :}
    \begin{align*}
        y_1(t) &= A \cos(\omega t + \varphi_1), \quad \varphi_1 = -k x_1 + \varphi_0 \\
        y_2(t) &= A \cos(\omega t + \varphi_2), \quad \varphi_2 = -k x_2 + \varphi_0 \\
        \Delta \varphi &= \varphi_2 - \varphi_1 = -k(x_2 - x_1) = -k \Delta x = -\frac{2\pi}{\lambda}\Delta x
    \end{align*}
    Le signe de $\Delta \varphi$ indique qui est en avance : si $\Delta \varphi > 0$ ($x_2 < x_1$ pour sens $+x$), $M_2$ est en avance sur $M_1$.
    \item \textbf{Représentation de Fresnel des deux points :} Deux vecteurs de même longueur $A$, tournant à même $\omega$, angle fixe $\Delta \varphi$ entre eux.
\end{enumerate}
\begin{aretenir}Équation onde : $y(x,t)=A\cos(\omega t - kx + \varphi_0)$. Point $x_0$ : $y(t)=A\cos(\omega t - kx_0 + \varphi_0)$. Différence de phase $\Delta\varphi = -k\Delta x$ (sens $+x$).\end{aretenir}
\end{methode}

\begin{exoresolu}[Équation de l'onde et mouvement de points sur une corde]
\textbf{Énoncé :} L'équation d'une onde progressive transversale se propageant sur une corde selon l'axe $Ox$ (vers $x>0$) est :
$y(x,t) = 0,040 \cos(20\pi t - 5\pi x + \pi/3)$ (unités SI : $y$ en m, $x$ en m, $t$ en s).
\begin{enumerate}
    \item Préciser l'amplitude $A$, la période $T$, la fréquence $f$, la longueur d'onde $\lambda$ et la célérité $c$ (valeur et direction).
    \item Écrire l'équation horaire du mouvement du point $M_1$ d'abscisse $x_1 = \SI{0,20}{m}$.
    \item Le point $M_2$ a pour abscisse $x_2 = \SI{0,50}{m}$. Déterminer la différence de phase $\Delta \varphi = \varphi_2 - \varphi_1$. En déduire le retard temporel $\tau$ de $M_2$ sur $M_1$.
    \item Représenter les vecteurs de Fresnel de $M_1$ et $M_2$ à l'instant $t=0$.
\end{enumerate}
\tcblower
\textbf{Solution détaillée :}
\begin{enumerate}
    \item Identification des paramètres dans $y = A \cos(\omega t - k x + \varphi_0)$ :
    $A = \SI{0,040}{m} = \SI{4,0}{cm}$.
    $\omega = 20\pi\,\text{rad/s} \Rightarrow T = 2\pi/\omega = \SI{0,100}{s}$, $f = \SI{10,0}{Hz}$.
    $k = 5\pi\,\text{rad/m} \Rightarrow \lambda = 2\pi/k = \SI{0,400}{m}$.
    $c = \omega/k = \lambda/T = (20\pi)/(5\pi) = \SI{4,00}{m/s}$. Direction : $+x$ (signe $-$ devant $kx$).
    \item Point $M_1(x_1=\SI{0,20}{m})$ :
    $\varphi_1 = -k x_1 + \pi/3 = -5\pi \times 0,20 + \pi/3 = -\pi + \pi/3 = -2\pi/3$.
    $y_1(t) = 0,040 \cos(20\pi t - 2\pi/3)\,\text{m}$.
    \item Point $M_2(x_2=\SI{0,50}{m})$ :
    $\varphi_2 = -5\pi \times 0,50 + \pi/3 = -2,5\pi + \pi/3 = -5\pi/2 + \pi/3 = -15\pi/6 + 2\pi/6 = -13\pi/6$.
    $\Delta \varphi = \varphi_2 - \varphi_1 = -13\pi/6 - (-4\pi/6) = -9\pi/6 = -3\pi/2$.
    Modulo $2\pi$ : $-3\pi/2 + 2\pi = \pi/2$.
    \textbf{Attention au sens :} $\Delta \varphi = -k\Delta x = -5\pi \times 0,30 = -1,5\pi \equiv \pi/2 [2\pi]$.
    $M_2$ est plus loin que $M_1$ dans le sens de propagation, donc $M_2$ est en \textbf{retard} sur $M_1$.
    Retard temporel $\tau = \frac{|\Delta \varphi|}{\omega} = \frac{3\pi/2}{20\pi} = \SI{0,075}{s} = \SI{75}{ms}$.
    (Vérif : $\tau = \Delta x / c = 0,30 / 4,00 = \SI{0,075}{s}$).
    \item \textbf{Vecteurs de Fresnel à $t=0$ :}
    Amplitude $\SI{4,0}{cm}$.
    $\vec{Y}_1$ : angle $\varphi_1 = -2\pi/3 = -120^\circ$.
    $\vec{Y}_2$ : angle $\varphi_2 \equiv \pi/2 = 90^\circ$ (ou $-270^\circ$).
    Différence de phase $\Delta \varphi = \varphi_2 - \varphi_1 = -150^\circ$ (ou $+210^\circ$). $M_2$ est en retard de $150^\circ$ sur $M_1$.
    \begin{popfigure}
    \begin{tikzpicture}[scale=1.2, >=Stealth]
        \draw[->] (-2.5,0) -- (2.5,0) node[right] {$Ox$};
        \draw[->] (0,-2.5) -- (0,2.5) node[above] {$Oy$};
        \draw[thick, popblue] (0,0) -- (-1,-1.73) node[below left] {$\vec{Y}_1\ (-120^\circ)$};
        \draw[thick, poporange] (0,0) -- (0,2) node[above] {$\vec{Y}_2\ (90^\circ)$};
        \draw[dashed, popmuted] (-1,-1.73) arc[start angle=-120, end angle=90, radius=2];
        \node[popmuted] at (-1.5, 0.5) {$210^\circ$};
        \draw[dashed, popgreen] (0,2) arc[start angle=90, end angle=-120, radius=2];
        \node[popgreen] at (1.5, 0.5) {$150^\circ$ (retard)};
    \end{tikzpicture}
    \legende{Vecteurs de Fresnel à $t=0$. $\vec{Y}_2$ est en retard de $150^\circ$ sur $\vec{Y}_1$ (angle trigonométrique négatif $\Delta\varphi=-150^\circ$).}
    \end{popfigure}
\end{enumerate}
\textbf{Conclusion :} L'équation de l'onde contient toute l'information. Le point $M_2$ (plus loin) vibre avec un retard de $\SI{75}{ms}$ (différence de phase $\Delta\varphi = -150^\circ$ modulo $2\pi$).
\end{exoresolu}

\begin{methode}[Résoudre un problème de synthèse : Onde progressive sur une corde / dans l'air]
\textbf{Objectif :} Mobiliser l'ensemble des savoir-faire (périodicité, célérité, équation, énergie, réflexion partielle) pour résoudre un problème type Bac.
\begin{enumerate}
    \item \textbf{Analyser le dispositif :} Source (fréquence imposée $f$), milieu (celérité $c$ ou $c=\sqrt{T/\mu}$ pour corde), conditions aux limites (fixe/libre $\rightarrow$ réflexion avec/sans inversion de phase).
    \item \textbf{Calculer les paramètres d'onde :} $T=1/f$, $\lambda=c/f$, $k=2\pi/\lambda$, $\omega=2\pi f$.
    \item \textbf{Écrire l'équation de l'onde incidente :} $y_i(x,t) = A_i \cos(\omega t - k x + \varphi_0)$. Choisir origine $x=0$ et $t=0$ pour simplifier $\varphi_0$.
    \item \textbf{Gérer la réflexion (si présent) :} 
    \begin{itemize}
        \item Extrémité fixe ($x=L$) : $y_r(L,t) = -y_i(L,t)$ $\rightarrow$ inversion de phase ($\pi$), nœud au point fixe.
        \item Extrémité libre : $y_r(L,t) = +y_i(L,t)$ $\rightarrow$ pas d'inversion, ventre au point libre.
    \end{itemize}
    Onde réfléchie : $y_r(x,t) = A_r \cos(\omega t + k x + \varphi_r)$ (sens $-x$).
    \item \textbf{Superposition (ondes stationnaires si réflexion totale) :} $y(x,t) = y_i + y_r$. Identifier nœuds ($A=0$) et ventres ($A=2A_i$).
    \item \textbf{Vérifier les unités et la cohérence :} Analyse dimensionnelle de chaque formule utilisée.
\end{enumerate}
\begin{aretenir}Corde : $c = \sqrt{T/\mu}$. Réflexion fixe : inversion $\pi$, nœud. Réflexion libre : pas d'inversion, ventre. Stationnaire : $y=2A\cos(kx)\cos(\omega t)$ (si réflexion fixe à $x=0$).\end{aretenir}
\end{methode}

\begin{exoresolu}[Synthèse : Onde sur corde, réflexion et ondes stationnaires]
\textbf{Énoncé :} Une corde de densité linéique $\mu = \SI{0,020}{kg/m}$ est tendue entre deux points fixes distants de $L = \SI{2,00}{m}$. La tension est $T = \SI{32}{N}$. L'extrémité $x=0$ est reliée à un vibreur imposant une oscillation $y(0,t) = 0,020 \cos(40\pi t)$ (SI). L'extrémité $x=L$ est fixe.
\begin{enumerate}
    \item Calculer la célérité $c$ des ondes sur la corde, la longueur d'onde $\lambda$ et le nombre d'onde $k$.
    \item Écrire l'expression de l'onde incidente $y_i(x,t)$ et de l'onde réfléchie $y_r(x,t)$ (amplitude réfléchie $A_r = \SI{0,015}{m}$).
    \item Donner l'expression de l'onde résultante $y(x,t)$. Montrer qu'il s'agit d'une onde stationnaire « imparfaite » (amplitude non nulle aux nœuds).
    \item Déterminer les positions des nœuds et des ventres de l'onde résultante.
\end{enumerate}
\tcblower
\textbf{Solution détaillée :}
\begin{enumerate}
    \item $c = \sqrt{T/\mu} = \sqrt{32 / 0,020} = \sqrt{1600} = \SI{40,0}{m/s}$.
    $\omega = 40\pi\,\text{rad/s} \Rightarrow f = \SI{20,0}{Hz}$, $T = \SI{0,0500}{s}$.
    $\lambda = c/f = 40,0 / 20,0 = \SI{2,00}{m}$.
    $k = 2\pi/\lambda = \pi\,\text{rad/m}$.
    \item \textbf{Onde incidente} (part de $x=0$ vers $+x$) : $y_i(x,t) = 0,020 \cos(40\pi t - \pi x)$.
    \textbf{Onde réfléchie} (revient vers $-x$, extrémité fixe à $x=L=2\,\text{m}$ $\rightarrow$ inversion de phase $\pi$) :
    À $x=L$, $y_r(L,t) = -y_i(L,t)$.
    $y_i(L,t) = 0,020 \cos(40\pi t - 2\pi) = 0,020 \cos(40\pi t)$.
    Donc $y_r(L,t) = -0,020 \cos(40\pi t) = 0,020 \cos(40\pi t + \pi)$.
    Forme générale onde réfléchie (sens $-x$) : $y_r(x,t) = A_r \cos(\omega t + k x + \varphi_r)$.
    À $x=L$ : $A_r \cos(40\pi t + 2\pi + \varphi_r) = 0,015 \cos(40\pi t + \varphi_r)$.
    Identification : $A_r = \SI{0,015}{m}$, $\varphi_r = \pi$.
    $y_r(x,t) = 0,015 \cos(40\pi t + \pi x + \pi) = -0,015 \cos(40\pi t + \pi x)$.
    \item \textbf{Onde résultante :} $y(x,t) = y_i + y_r = 0,020 \cos(40\pi t - \pi x) - 0,015 \cos(40\pi t + \pi x)$.
    Utiliser formules $\cos(A \mp B)$ :
    $y = (0,020 - 0,015)\cos(40\pi t)\cos(\pi x) + (0,020 + 0,015)\sin(40\pi t)\sin(\pi x)$.
    $y(x,t) = 0,005 \cos(40\pi t)\cos(\pi x) + 0,035 \sin(40\pi t)\sin(\pi x)$.
    Ce n'est pas une onde stationnaire pure (pas de factorisation $X(x)T(t)$ simple avec des zéros fixes pour tout $t$). C'est un mélange d'onde stationnaire et progressive (onde stationnaire « imparfaite » car réflexion partielle $A_r \neq A_i$).
    \item \textbf{Positions particulières :} Les « nœuds » (amplitude minimale) et « ventres » (amplitude maximale) ne sont plus fixes dans le temps car les coefficients de $\cos(\omega t)$ et $\sin(\omega t)$ ne s'annulent pas simultanément aux mêmes $x$.
    Cependant, on peut chercher les extrema de l'amplitude vibratoire en fonction de $x$.
    L'amplitude instantanée varie. Pour une onde stationnaire pure ($A_r=A_i$), nœuds pour $\cos(kx)=0 \Rightarrow x = (2n+1)\lambda/4 = (2n+1)\times 0,5\,\text{m}$.
    Ici, l'amplitude minimale (pseudo-nœuds) se trouve près de ces positions.
\end{enumerate}
\textbf{Conclusion :} La réflexion partielle ($A_r < A_i$) empêche la formation d'ondes stationnaires parfaites (nœuds immobiles d'amplitude nulle). L'énergie n'est pas totalement piégée.
\end{exoresolu}

\begin{attention}[Les 5 erreurs qui coûtent des points au Bac]
\begin{enumerate}
    \item \textbf{Confusion Période / Fréquence / Pulsation :} Oublier $T = 1/f$ ou $\omega = 2\pi f$. Écrire $\omega = 2\pi/T$ mais utiliser $f$ dans le calcul. \textbf{Réflexe :} Toujours écrire $\omega = 2\pi f = 2\pi/T$ explicitement avant de calculer.
    \item \textbf{Unités non homogènes (le piège n°1) :} Mélanger $\text{cm}$ et $\text{m}$, $\text{ms}$ et $\text{s}$, $\text{kHz}$ et $\text{Hz}$ dans une même formule ($\lambda = c/f$, $\Delta \varphi = 2\pi d/\lambda$). \textbf{Réflexe :} Convertir \textbf{tout} en SI (m, s, Hz, rad) dès la lecture de l'énoncé.
    \item \textbf{Phase en degrés vs radians :} Calculer $\Delta \varphi = 360^\circ \times d/\lambda$ puis l'utiliser dans $\cos(\omega t + \Delta \varphi)$ où $\omega$ est en $\text{rad/s}$. \textbf{Règle :} En physique, les arguments des fonctions trigonométriques sont \textbf{toujours en radians}. Convertir systématiquement : $\varphi_{\text{rad}} = \varphi_{\text{deg}} \times \pi/180$.
    \item \textbf{Sens du retard / avance de phase :} Dire « le point $M_2$ est en avance sur $M_1$ » alors que $M_2$ est plus loin de la source dans le sens de propagation. \textbf{Règle :} L'onde propage l'information de la source vers l'extérieur. Le point \textbf{lointain} reçoit l'information \textbf{plus tard} $\rightarrow$ il est en \textbf{retard} de phase ($\Delta \varphi < 0$ pour $\varphi_{\text{lointain}} - \varphi_{\text{proche}}$).
    \item \textbf{Onde longitudinale / transversale et polarisation :} Affirmer qu'une onde sonore (longitudinale) peut être polarisée, ou qu'une onde sur corde (transversale) ne peut pas l'être. \textbf{Rappel :} \textbf{Seules les ondes transversales sont polarisables.} (L'onde sonore dans l'air n'a pas de polarisation).
    \item \textbf{(Bonus) Stroboscopie :} Confondre $f_s = n f$ (immobilité) et $f_s = f \pm \Delta f$ (ralenti). Oublier la valeur absolue pour la fréquence apparente : $f_{\text{app}} = |f_s - f|$.
\end{enumerate}
\end{attention}

\newpage

\coursec{Exercices}

\coursec{Leçon P14 : Systèmes oscillants et propagation des ondes mécaniques}

\courssub{Généralités sur les phénomènes périodiques sinusoïdaux}

\begin{definition}[Phénomène périodique et grandeur sinusoïdale]
Un phénomène est \emph{périodique} s'il se reproduit à l'identique après un intervalle de temps constant appelé \cle{période} $T$ (en secondes, \si{\second}). La \cle{fréquence} $f$ (en hertz, \si{\hertz}) est l'inverse de la période : $f = \dfrac{1}{T}$.
Une grandeur physique $x(t)$ est \emph{sinusoïdale} si elle s'écrit sous la forme :
\[ x(t) = X_m \cos(\omega t + \varphi) \quad \text{ou} \quad x(t) = X_m \sin(\omega t + \varphi) \]
où :
\begin{itemize}
    \item $X_m$ est l'\cle{amplitude} (valeur maximale, même unité que $x$) ;
    \item $\omega = 2\pi f = \dfrac{2\pi}{T}$ est la \cle{pulsation} (en \si{\radian\per\second}) ;
    \item $\varphi$ est la \cle{phase à l'origine des dates} (en radians), valeur de l'argument à $t=0$.
\end{itemize}
\end{definition}

\begin{propriete}[Valeur efficace]
Pour une grandeur sinusoïdale d'amplitude $X_m$, la \cle{valeur efficace} (ou valeur quadratique moyenne) est :
\[ X_{\text{eff}} = \frac{X_m}{\sqrt{2}} \]
C'est la valeur d'une grandeur continue qui dissiperait la même puissance moyenne dans une charge résistive. \emph{Démonstration exigible au Bac : calcul de la moyenne quadratique sur une période.}
\end{propriete}

\begin{exemplebox}[Tension secteur au Cameroun]
La tension délivrée par le réseau ENEO est une tension sinusoïdale de valeur efficace $U_{\text{eff}} = 230\ \text{V}$ (norme nominale) et de fréquence $f = 50\ \text{Hz}$. Sa période est $T = 20\ \text{ms}$, sa pulsation $\omega = 100\pi \approx 314\ \text{rad/s}$, et son amplitude $U_m = 230\sqrt{2} \approx 325\ \text{V}$.
\end{exemplebox}

\begin{methode}[Écrire l'équation horaire d'une grandeur sinusoïdale]
\begin{enumerate}
    \item Identifier l'amplitude $X_m$ (lecture sur oscillogramme ou donnée).
    \item Déterminer la période $T$ (ou la fréquence $f$) et en déduire $\omega = 2\pi/T$.
    \item Choisir l'origine des dates ($t=0$) pour simplifier la phase $\varphi$ :
    \begin{itemize}
        \item Si $x(0) = X_m$ (maximum) $\Rightarrow$ fonction cosinus, $\varphi = 0$.
        \item Si $x(0) = 0$ et $x$ croît $\Rightarrow$ fonction sinus, $\varphi = 0$.
        \item Sinon, calculer $\varphi$ à partir de la condition initiale.
    \end{itemize}
    \item Écrire $x(t) = X_m \cos(\omega t + \varphi)$ (privilégier le cosinus par convention).
\end{enumerate}
\end{methode}

\begin{attention}[Pièges fréquents]
\begin{itemize}
    \item Confondre pulsation $\omega$ (rad/s) et fréquence $f$ (Hz). Vérifier : $\omega = 2\pi f$.
    \item Oublier le facteur $\sqrt{2}$ entre amplitude et valeur efficace.
    \item Ne pas mettre l'argument du cosinus en \emph{radians} (calculatrices en mode RAD).
    \item La phase à l'origine $\varphi$ est définie \emph{modulo $2\pi$} ; on la ramène généralement dans $]-\pi, \pi]$.
\end{itemize}
\end{attention}

\courssub{Stroboscopie}

\begin{experience}[Principe de la stroboscopie]
\textbf{Matériel :} Disque de Newton (ou roue dentée) muni d'un repère, moteur à vitesse variable, stroboscope à fréquence réglable $f_s$.
\textbf{Protocole :} On fait tourner le disque à fréquence $f$ inconnue. On éclaire par flashes brefs de fréquence $f_s$. On observe l'aspect du repère.
\textbf{Observations et interprétation :}
\begin{itemize}
    \item \textbf{Immobilité apparente (repère unique) :} $f_s = f$ (ou $f_s = nf$ avec $n$ entier, mais le repère paraît alors $n$ fois). La \emph{plus grande} fréquence $f_s$ donnant une immobilité unique correspond à $f_s = f$.
    \item \textbf{Mouvement ralenti apparent :} Si $f_s$ est proche de $f$, le repère semble tourner à la fréquence apparente $f_{\text{app}} = |f_s - f|$. Le sens apparent est le sens réel si $f_s > f$, inverse si $f_s < f$.
    \item \textbf{Immobilité multiple :} Si $f_s = k f$ ($k$ entier $>1$), on voit $k$ repères immobiles.
\end{itemize}
\end{experience}

\begin{propriete}[Relations stroboscopiques]
Soit un objet tournant à la fréquence $f$ (ou une onde de fréquence $f$), éclairé par un stroboscope de fréquence $f_s$.
\begin{itemize}
    \item Condition d'immobilité apparente (unique) : $f_s = f$.
    \item Fréquence apparente de rotation : $f_{\text{app}} = |f_s - f|$.
    \item Si $f_s = \frac{f}{n}$ ($n$ entier), on voit $n$ images fixes.
\end{itemize}
\end{propriete}

\begin{exoresolu}[Mesure de la vitesse d'un mototaxi à Douala]
Un mototaxi a une roue de rayon $R = 30\ \text{cm}$. On éclaire la roue (munie d'un repère) avec un stroboscope. La plus grande fréquence donnant une immobilité unique est $f_s = 12\ \text{Hz}$.
\tcblower
La fréquence de rotation de la roue est $f = f_s = 12\ \text{Hz}$.
Vitesse angulaire : $\omega = 2\pi f = 24\pi \approx 75,4\ \text{rad/s}$.
Vitesse linéaire du mototaxi (sans glissement) : $v = \omega R = 75,4 \times 0,30 \approx 22,6\ \text{m/s} \approx 81\ \text{km/h}$.
\end{exoresolu}

\courssub{Représentation de Fresnel}

\begin{definition}[Vecteur de Fresnel (faisceau tournant)]
À une grandeur sinusoïdale $x(t) = X_m \cos(\omega t + \varphi)$, on associe un \cle{vecteur de Fresnel} $\vec{X}$ de module $X_m$ qui tourne dans le sens trigonométrique (anti-horaire) avec la vitesse angulaire $\omega$. À l'instant $t$, ce vecteur forme un angle $\omega t + \varphi$ avec l'axe $Ox$. La grandeur $x(t)$ est la \emph{projection orthogonale} de $\vec{X}$ sur l'axe $Ox$ :
\[ x(t) = \vec{X} \cdot \vec{u}_x = X_m \cos(\omega t + \varphi) \]
\end{definition}

\begin{propriete}[Somme de deux grandeurs sinusoïdales de même fréquence]
La somme de deux grandeurs sinusoïdales de même pulsation $\omega$ est une grandeur sinusoïdale de même pulsation $\omega$.
Si $x_1(t) = X_{1m}\cos(\omega t + \varphi_1)$ et $x_2(t) = X_{2m}\cos(\omega t + \varphi_2)$, alors $x(t) = x_1(t) + x_2(t) = X_m\cos(\omega t + \varphi)$ où le vecteur de Fresnel $\vec{X}$ est la somme vectorielle :
\[ \vec{X} = \vec{X}_1 + \vec{X}_2 \]
On en déduit par le théorème d'Al-Kashi (ou projection sur axes) :
\begin{align*}
    X_m^2 &= X_{1m}^2 + X_{2m}^2 + 2 X_{1m} X_{2m} \cos(\varphi_2 - \varphi_1) \\
    \tan\varphi &= \frac{X_{1m}\sin\varphi_1 + X_{2m}\sin\varphi_2}{X_{1m}\cos\varphi_1 + X_{2m}\cos\varphi_2}
\end{align*}
\emph{La construction graphique (règle, compas, rapporteur) est exigible au Bac.}
\end{propriete}

\begin{exemplebox}[Somme de deux tensions alternatives]
$u_1(t) = 6\cos(100\pi t)$ V et $u_2(t) = 8\cos(100\pi t + \pi/2)$ V.
Vecteurs : $\vec{U}_1 = 6\ \text{V}$ sur $Ox$, $\vec{U}_2 = 8\ \text{V}$ sur $Oy$.
$\vec{U} = \vec{U}_1 + \vec{U}_2$ : module $U_m = \sqrt{6^2+8^2} = 10\ \text{V}$.
Phase $\varphi = \arctan(8/6) \approx 0,927\ \text{rad} \approx 53,1^\circ$.
$u(t) = 10\cos(100\pi t + 0,927)$ V.
\end{exemplebox}

\begin{popfigure}
\begin{tikzpicture}[scale=0.8, >=Stealth]
\draw[->] (-1,0) -- (7,0) node[right] {$Ox$};
\draw[->] (0,-1) -- (0,5) node[above] {$Oy$};
\draw[popblue, thick, ->] (0,0) -- (6,0) node[midway, below] {$\vec{U}_1 = 6\ \text{V}$};
\draw[poporange, thick, ->] (0,0) -- (0,4) node[midway, left] {$\vec{U}_2 = 8\ \text{V}$};
\draw[popgreen, very thick, ->] (0,0) -- (6,4) node[midway, above right] {$\vec{U} = 10\ \text{V}$};
\draw[dashed, popmuted] (6,0) -- (6,4) -- (0,4);
\draw[popgreen] (1.2,0) arc (0:33.69:1.2) node[midway, right] {$\varphi$};
\end{tikzpicture}
\legende{Construction de Fresnel pour la somme de deux tensions en quadrature de phase.}
\end{popfigure}

\courssub{Ondes mécaniques progressives}

\begin{definition}[Onde mécanique progressive]
Une \cle{onde mécanique progressive} est la propagation d'une perturbation locale dans un milieu matériel, sans transport de matière (les particules oscillent autour de leur position d'équilibre), mais avec transport d'énergie. La perturbation se propage à la \cle{celerité} $c$ (en \si{\meter\per\second}) qui dépend des propriétés du milieu (inertie, élasticité).
\end{definition}

\begin{propriete}[Ondes transversales et longitudinales]
\begin{itemize}
    \item \textbf{Onde transversale :} La perturbation (élongation) est \emph{perpendiculaire} à la direction de propagation.
    \emph{Exemples :} Onde sur une corde tendue, onde à la surface de l'eau (cuve à ondes), ondes sismiques de type S (secondaires), lumière (onde électromagnétique, non mécanique).
    \item \textbf{Onde longitudinale :} La perturbation (variation de pression ou de densité) est \emph{parallèle} à la direction de propagation.
    \emph{Exemples :} Onde sonore dans l'air (compressions/dilatations), onde dans un ressort (slinky), ondes sismiques de type P (primaires).
\end{itemize}
\end{propriete}

\begin{savaistu}[Ondes sismiques et barrage de Lom Pangar]
Les séismes génèrent des ondes P (longitudinales, rapides, $\approx 6-8\ \text{km/s}$) et des ondes S (transversales, plus lentes, $\approx 3-5\ \text{km/s}$). L'étude de leur propagation permet de sonder la structure interne de la Terre. Au Cameroun, la surveillance sismique est cruciale pour la sûreté des grands ouvrages comme les barrages de \textbf{Lom Pangar}, \textbf{Nachtigal} ou \textbf{Edéa}.
\end{savaistu}

\courssub{Onde progressive sinusoïdale : double périodicité}

\begin{definition}[Onde progressive sinusoïdale]
Une onde est \emph{sinusoïdale} si la source $S$ vibre selon une loi sinusoïdale $y_S(t) = A\cos(\omega t + \varphi_S)$. Tout point $M$ du milieu, situé à la distance $d = SM$ de la source, vibre avec le même amplitude $A$ (sans amortissement) et la même pulsation $\omega$, mais avec un \cle{retard temporel} $\theta = \dfrac{d}{c}$ :
\[ y_M(t) = A\cos\left[\omega\left(t - \frac{d}{c}\right) + \varphi_S\right] = A\cos(\omega t - k d + \varphi_S) \]
où $k = \dfrac{\omega}{c} = \dfrac{2\pi}{\lambda}$ est le \cle{nombre d'onde} (en \si{\radian\per\meter}) et $\lambda = \dfrac{c}{f} = cT$ est la \cle{longueur d'onde} (en \si{\meter}).
\end{definition}

\begin{propriete}[Double périodicité]
Une onde progressive sinusoïdale présente une double périodicité :
\begin{itemize}
    \item \textbf{Périodicité temporelle :} En un point fixe $M$, l'élongation se répète tous les $T$ : $y_M(t+T) = y_M(t)$.
    \item \textbf{Périodicité spatiale :} À un instant figé $t_0$, le profil de l'onde (élongation vs position) se répète tous les $\lambda$ : $y(x+\lambda, t_0) = y(x, t_0)$.
\end{itemize}
Relation fondamentale : $\boxed{\lambda = c\,T = \dfrac{c}{f}}$.
\end{propriete}

\begin{propriete}[Différence de phase entre deux points]
Pour deux points $M$ et $N$ distants de $d = MN$ le long de la direction de propagation, le retard temporel est $\theta = d/c$. La différence de phase (de $M$ vers $N$) est :
\[ \Delta\varphi = \varphi_N - \varphi_M = -\omega\theta = -k d = -\frac{2\pi d}{\lambda} \]
\begin{itemize}
    \item \textbf{En phase :} $\Delta\varphi \equiv 0 \pmod{2\pi} \iff d = k\lambda \quad (k \in \mathbb{Z})$.
    \item \textbf{En opposition de phase :} $\Delta\varphi \equiv \pi \pmod{2\pi} \iff d = \left(k+\frac{1}{2}\right)\lambda \quad (k \in \mathbb{Z})$.
\end{itemize}
\end{propriete}

\begin{methode}[Étudier l'état vibratoire de deux points d'une onde]
\begin{enumerate}
    \item Calculer la longueur d'onde $\lambda = c/f$ (ou la lire sur un profil spatial).
    \item Calculer la distance $d$ séparant les deux points le long de la propagation.
    \item Calculer le rapport $d/\lambda$.
    \item Conclure :
    \begin{itemize}
        \item Si $d/\lambda$ est un entier $\Rightarrow$ \textbf{en phase}.
        \item Si $d/\lambda$ est un demi-entier ($n+1/2$) $\Rightarrow$ \textbf{en opposition de phase}.
        \item Sinon $\Rightarrow$ \textbf{ni l'un ni l'autre} (différence de phase $\Delta\varphi = -2\pi d/\lambda$).
    \end{itemize}
\end{enumerate}
\end{methode}

\begin{attention}[Erreurs classiques sur les ondes]
\begin{itemize}
    \item Confondre la vitesse de propagation $c$ (de l'onde) et la vitesse des particules du milieu $v = dy/dt$ (vitesse de vibration).
    \item Croire qu'une onde transporte la matière : un bouchon sur l'eau monte et descend mais ne suit pas l'onde.
    \item Oublier le signe moins dans l'argument $\omega t - k x$ pour une onde se propageant dans le sens $x$ croissant.
    \item Mélanger période $T$ (temps) et longueur d'onde $\lambda$ (espace) dans les calculs de phase.
\end{itemize}
\end{attention}

\courssub{Synthèse et méthodes de résolution}

\begin{methode}[Résolution type d'un problème d'ondes (Bac)]
\begin{enumerate}
    \item \textbf{Identifier les données :} Fréquence $f$ (ou période $T$), célérité $c$, positions des points ($SM$, $MN$), équation de la source.
    \item \textbf{Calculer les grandeurs caractéristiques :} $\omega = 2\pi f$, $T = 1/f$, $\lambda = c/f$, $k = 2\pi/\lambda$.
    \item \textbf{Écrire l'équation de la source :} $y_S(t) = A\cos(\omega t + \varphi_S)$.
    \item \textbf{Écrire l'équation d'un point $M$ :} $y_M(t) = A\cos(\omega t - k\,SM + \varphi_S)$.
    \item \textbf{Comparer deux points $M$ et $N$ :} Calculer $\Delta\varphi = -k\,MN = -2\pi MN/\lambda$. Déduire l'état vibratoire.
    \item \textbf{Exploiter la stroboscopie :} $f_s = f$ pour immobilité ; $f_{\text{app}} = |f_s - f|$ pour mouvement ralenti.
    \item \textbf{Exploiter l'oscilloscope :} $T = (\text{nb div horizontales}) \times b$ ; $U_m = (\text{nb div verticales}/2) \times k$ (crête-à-crête) ou amplitude directe.
\end{enumerate}
\end{methode}

\courssub{Je vérifie mes connaissances}

\begin{exercice}[QCM : grandeurs sinusoïdales]
Pour chaque question, une seule réponse est exacte. Justifier brièvement chaque choix.
\begin{enumerate}
\item La tension $u(t) = 5\cos\left(100\pi t + \dfrac{\pi}{4}\right)$ (en volts) a pour fréquence :
\begin{enumerate}
\item[a)] $f = 100\ \text{Hz}$ \quad b)] $f = 50\ \text{Hz}$ \quad c)] $f = 5\ \text{Hz}$ \quad d)] $f = 314\ \text{Hz}$
\end{enumerate}
\item La valeur efficace de cette tension vaut :
\begin{enumerate}
\item[a)] $5\ \text{V}$ \quad b)] $5\sqrt{2}\ \text{V}$ \quad c)] $\dfrac{5}{\sqrt{2}}\ \text{V}$ \quad d)] $2{,}5\ \text{V}$
\end{enumerate}
\item La phase à l'origine des dates de $u(t)$ est :
\begin{enumerate}
\item[a)] $\pi\ \text{rad}$ \quad b)] $\dfrac{\pi}{4}\ \text{rad}$ \quad c)] $100\pi\ \text{rad}$ \quad d)] $0\ \text{rad}$
\end{enumerate}
\item Une onde mécanique se propageant le long d'une corde tendue est :
\begin{enumerate}
\item[a)] longitudinale \quad b)] transversale \quad c)] électromagnétique \quad d)] stationnaire
\end{enumerate}
\end{enumerate}
\end{exercice}

\begin{exercice}[Vrai ou faux justifié]
Répondre par \emph{vrai} ou \emph{faux} en justifiant chaque réponse par une phrase ou un calcul.
\begin{enumerate}
\item Deux points d'un milieu de propagation distants d'une longueur d'onde vibrent en opposition de phase.
\item La longueur d'onde $\lambda$ d'une onde sinusoïdale de fréquence $f = 200\ \text{Hz}$ se propageant à la célérité $c = 340\ \text{m/s}$ vaut $\lambda = 1{,}7\ \text{m}$.
\item Si un disque en rotation paraît immobile sous un stroboscope de fréquence $f_s = 50\ \text{Hz}$, alors la fréquence de rotation du disque est nécessairement $50\ \text{Hz}$.
\item La somme de deux grandeurs sinusoïdales de même fréquence est une grandeur sinusoïdale de même fréquence.
\item Une onde mécanique transporte de la matière sur de grandes distances.
\end{enumerate}
\end{exercice}

\begin{exercice}[Définitions et vocabulaire]
\begin{enumerate}
\item Définir : période, fréquence, amplitude, phase à l'origine des dates d'une grandeur sinusoïdale.
\item Définir une onde mécanique progressive. Citer deux exemples d'ondes transversales et deux exemples d'ondes longitudinales observables au laboratoire ou dans la vie courante.
\item Qu'appelle-t-on \cle{double périodicité} d'une onde progressive sinusoïdale ? Donner les deux grandeurs concernées et la relation qui les lie.
\item Deux points $M$ et $N$ d'un milieu sont distants de $d$. Donner, en fonction de $\lambda$, la condition pour qu'ils vibrent en phase, puis celle pour qu'ils vibrent en opposition de phase.
\end{enumerate}
\end{exercice}

\begin{exercice}[Calculs directs]
\begin{enumerate}
\item Un oscillateur harmonique a pour équation horaire $x(t) = 4\cos\left(20\pi t - \dfrac{\pi}{3}\right)$ (en cm). Déterminer son amplitude, sa pulsation, sa période, sa fréquence et sa phase à l'origine.
\item Un son émis par un haut-parleur a pour fréquence $f = 500\ \text{Hz}$ et se propage dans l'air à $c = 340\ \text{m/s}$. Calculer sa longueur d'onde.
\item Une onde met $\theta = 0{,}25\ \text{s}$ pour parcourir la distance $SM = 2{,}0\ \text{m}$ le long d'une corde. Calculer sa célérité.
\item Un moteur tourne à $1\,500\ \text{tr/min}$. Exprimer sa fréquence de rotation en hertz.
\end{enumerate}
\end{exercice}

\courssub{Je m'entraîne}

\begin{exercice}[Exploitation d'un oscillogramme]
On visualise à l'oscilloscope la tension délivrée par un générateur basse fréquence (GBF). Les réglages sont : balayage horizontal $b = 5\ \text{ms/div}$ ; sensibilité verticale $k = 2\ \text{V/div}$. La courbe observée est une sinusoïde dont une période occupe $4$ divisions horizontales et dont la hauteur crête-à-crête occupe $6$ divisions verticales. À la date $t = 0$, la courbe passe par son maximum.
\begin{enumerate}
\item Déterminer la période $T$ et la fréquence $f$ de la tension.
\item Déterminer l'amplitude $U_m$ et la valeur efficace $U_{\text{eff}}$.
\item Écrire l'équation horaire $u(t)$ de cette tension.
\item Représenter le vecteur de Fresnel associé à $u(t)$ à la date $t = 0$ (échelle : $1\ \text{cm} \leftrightarrow 1\ \text{V}$).
\end{enumerate}
\end{exercice}

\begin{exercice}[Stroboscopie : disque d'un moteur]
Un disque solidaire de l'arbre d'un moteur porte un unique repère peint en blanc. On l'éclaire à l'aide d'un stroboscope dont on fait varier la fréquence $f_s$ des éclairs.
\begin{enumerate}
\item La plus grande fréquence des éclairs pour laquelle le repère paraît immobile et unique est $f_s = 50\ \text{Hz}$. En déduire la fréquence de rotation $f$ du moteur, puis sa vitesse de rotation en tr/min.
\item Que observe-t-on si l'on règle le stroboscope sur $f_s = 25\ \text{Hz}$ ? Justifier.
\item On règle maintenant le stroboscope sur $f_s = 51\ \text{Hz}$. Décrire le mouvement apparent du repère et calculer sa fréquence apparente de rotation.
\end{enumerate}
\end{exercice}

\begin{exercice}[Somme de deux tensions par la construction de Fresnel]
Deux tensions sinusoïdales de même fréquence sont appliquées en série :
\[ u_1(t) = 6\cos(100\pi t) \quad \text{et} \quad u_2(t) = 8\cos\left(100\pi t + \frac{\pi}{2}\right) \quad (\text{en volts}). \]
\begin{enumerate}
\item Construire à l'échelle ($1\ \text{cm} \leftrightarrow 1\ \text{V}$) les vecteurs de Fresnel $\vec{V}_1$ et $\vec{V}_2$ associés à $u_1$ et $u_2$, puis le vecteur $\vec{V} = \vec{V}_1 + \vec{V}_2$.
\item Déduire de la construction, puis vérifier par le calcul, l'amplitude $U_m$ de la tension somme $u(t) = u_1(t) + u_2(t)$.
\item Déterminer la phase à l'origine $\varphi$ de $u(t)$ et écrire son équation horaire.
\item Quelle serait l'amplitude de la somme si les deux tensions étaient en opposition de phase ?
\end{enumerate}
\end{exercice}

\begin{exercice}[Onde le long d'une corde]
Une source $S$, placée à l'extrémité d'une corde tendue, est animée d'un mouvement sinusoïdal d'équation $y_S(t) = 2\cos(40\pi t)$, où $y_S$ est en cm et $t$ en s. L'onde se propage le long de la corde avec la célérité $c = 8\ \text{m/s}$. On néglige tout amortissement.
\begin{enumerate}
\item Déterminer la période $T$, la fréquence $f$ et la longueur d'onde $\lambda$ de l'onde.
\item Établir l'équation horaire $y_M(t)$ du mouvement d'un point $M$ de la corde situé à la distance $d = SM = 0{,}60\ \text{m}$ de la source.
\item Comparer les états vibratoires de $S$ et de $M$ : vibrent-ils en phase, en opposition de phase, ou ni l'un ni l'autre ? Justifier.
\item Déterminer les positions des points de la corde qui vibrent en phase avec $S$, puis celles des points qui vibrent en opposition de phase avec $S$.
\end{enumerate}
\end{exercice}

\courssub{Je me prépare au Bac}

\begin{exercice}[Type Bac : étude d'une onde sur une corde vibrante]
Au laboratoire de physique d'un lycée de Yaoundé, un élève de Terminale C attache une corde élastique de longueur $L = 2{,}0\ \text{m}$ à la lame d'un vibreur alimenté par un GBF. La fréquence du vibreur est réglée à $f = 100\ \text{Hz}$. L'extrémité $S$ de la corde, assimilée à la source, a pour équation horaire $y_S(t) = 3\cos(200\pi t)$ (en cm). On néglige l'amortissement.

\textbf{Partie A : célérité de l'onde}

L'élève éclaire la corde avec un stroboscope. La plus grande fréquence des éclairs pour laquelle la corde paraît immobile et unique est $f_s = 100\ \text{Hz}$.
\begin{enumerate}
\item Justifier que cette observation confirme la fréquence du vibreur.
\item L'élève mesure la distance entre deux points consécutifs de la corde vibrant en phase : il trouve $4{,}0\ \text{cm}$. Que représente cette distance ? En déduire la célérité $c$ de l'onde le long de la corde.
\end{enumerate}

\textbf{Partie B : états vibratoires}
\begin{enumerate}[resume]
\item Établir l'équation horaire $y_M(t)$ d'un point $M$ de la corde situé à $SM = 30\ \text{cm}$.
\item Déterminer l'état vibratoire de $M$ par rapport à $S$.
\item Déterminer le nombre et les positions des points de la corde qui vibrent en opposition de phase avec la source $S$.
\end{enumerate}

\textbf{Partie C : problème inverse}

On modifie la tension de la corde, ce qui change la célérité de l'onde, la fréquence restant égale à $100\ \text{Hz}$. L'élève constate que deux points $M$ et $N$ de la corde, distants de $MN = 45\ \text{cm}$, vibrent désormais en opposition de phase.
\begin{enumerate}[resume]
\item Montrer que les valeurs possibles de la longueur d'onde vérifient $\lambda = \dfrac{2 \times 45}{2k+1}$ (en cm), avec $k \in \mathbb{N}$.
\item En déduire les valeurs possibles de la célérité de l'onde pour $k = 0$, $k = 1$ et $k = 2$.
\end{enumerate}
\end{exercice}

\begin{exercice}[Type Bac : cuve à ondes et mesure à l'oscilloscope]
Dans le cadre d'un TP, un groupe d'élèves de Terminale C du lycée de Douala étudie la propagation d'ondes à la surface de l'eau d'une cuve à ondes. Un vibreur de fréquence réglable frappe la surface de l'eau en un point $S$. Deux pointes $M_1$ et $M_2$, alignées avec $S$, sont placées à $SM_1 = 6{,}0\ \text{cm}$ et $SM_2 = 10{,}5\ \text{cm}$. Chaque pointe est reliée à une voie d'un oscilloscope bicourbe, qui délivre une tension proportionnelle à l'élongation du point de la surface.

\textbf{Partie A : réglage de la fréquence}

La fréquence du vibreur est fixée à $f = 25\ \text{Hz}$.
\begin{enumerate}
\item Calculer la période $T$ de l'onde.
\item Rappeler la définition de la longueur d'onde $\lambda$ et la relation qui la lie à $c$ et $T$.
\end{enumerate}

\textbf{Partie B : exploitation de l'oscillogramme}

Sur l'écran (balayage $b = 10\ \text{ms/div}$), les deux sinusoïdes observées sont décalées horizontalement de $1{,}5$ division.
\begin{enumerate}[resume]
\item Déterminer le retard temporel $\theta$ du passage de l'onde en $M_2$ par rapport à $M_1$.
\item En déduire la célérité $c$ de l'onde à la surface de l'eau.
\item Calculer la longueur d'onde $\lambda$ et vérifier la relation $\lambda = cT$.
\item Les points $M_1$ et $M_2$ vibrent-ils en phase ou en opposition de phase ? Justifier à l'aide de la valeur de $\lambda$.
\end{enumerate}

\textbf{Partie C : stroboscopie}
\begin{enumerate}[resume]
\item On éclaire la cuve avec un stroboscope de fréquence $f_s = 25\ \text{Hz}$. Que observe-t-on ? Justifier.
\item On règle le stroboscope sur $f_s = 24\ \text{Hz}$. Décrire qualitativement le phénomène observé et calculer la fréquence apparente des ondes.
\end{enumerate}
\end{exercice}

\begin{exercice}[Type Bac : synthèse oscillations et ondes — le haut-parleur du bal du lycée]
Pour animer la fête de fin d'année d'un lycée de Bafoussam, on utilise un haut-parleur alimenté par un amplificateur. La tension aux bornes du haut-parleur est $u(t) = 10\cos(2000\pi t)$ (en volts). Le son émis se propage dans l'air à la célérité $c = 340\ \text{m/s}$.

\textbf{Partie A : la tension électrique}
\begin{enumerate}
\item Déterminer l'amplitude $U_m$, la valeur efficace $U_{\text{eff}}$, la période $T$ et la fréquence $f$ de la tension.
\item Représenter le vecteur de Fresnel associé à $u(t)$ à $t = 0$ (échelle : $1\ \text{cm} \leftrightarrow 2\ \text{V}$).
\item Une seconde voie de l'amplificateur délivre $u'(t) = 10\cos\left(2000\pi t + \dfrac{\pi}{2}\right)$ (en volts). Par la construction de Fresnel, déterminer l'amplitude et la phase à l'origine de la tension somme $u + u'$. Vérifier les résultats par le calcul.
\end{enumerate}

\textbf{Partie B : l'onde sonore}
\begin{enumerate}[resume]
\item Le son émis par le haut-parleur est-il une onde transversale ou longitudinale ? Justifier.
\item Calculer la longueur d'onde $\lambda$ du son émis.
\item Deux élèves $A$ et $B$ sont placés sur un même rayon issu du haut-parleur, à $d_A = 5{,}10\ \text{m}$ et $d_B = 6{,}80\ \text{m}$ de celui-ci. Comparer leurs états vibratoires (en phase ou en opposition de phase). Justifier.
\item Calculer le retard avec lequel l'élève $B$ perçoit le son par rapport à l'élève $A$.
\end{enumerate}

\textbf{Partie C : éclairage stroboscopique de la membrane}

Pour observer la membrane du haut-parleur, on l'éclaire avec un stroboscope.
\begin{enumerate}[resume]
\item Quelle fréquence d'éclairs faut-il choisir pour que la membrane paraisse immobile ? Justifier.
\item La fréquence du son étant trop élevée pour un stroboscope classique, proposer une méthode expérimentale alternative pour mesurer la fréquence de vibration de la membrane.
\end{enumerate}
\end{exercice}

\newpage

\coursec{Corrigés des exercices}

\begin{corrige}[Exercice 1 — QCM : grandeurs sinusoïdales]
\begin{enumerate}
\item \textbf{Réponse b) :} $f = \SI{50}{Hz}$. En effet, $\omega = \SI{100\pi}{rad/s}$ et $f = \dfrac{\omega}{2\pi} = \dfrac{100\pi}{2\pi} = \SI{50}{Hz}$.
\item \textbf{Réponse c) :} $U_{\text{eff}} = \dfrac{U_m}{\sqrt{2}} = \dfrac{\SI{5}{V}}{\sqrt{2}} \approx \SI{3,54}{V}$.
\item \textbf{Réponse b) :} l'argument du cosinus à $t=0$ est $\varphi = \dfrac{\pi}{4}\ \text{rad}$.
\item \textbf{Réponse b) :} le long d'une corde tendue, l'élongation des points de la corde est perpendiculaire à la direction de propagation : l'onde est \emph{transversale}.
\end{enumerate}
\end{corrige}

\begin{corrige}[Exercice 2 — Vrai ou faux justifié]
\begin{enumerate}
\item \textbf{Faux.} Deux points distants de $d = \lambda$ vibrent \emph{en phase}, car $d = k\lambda$ avec $k=1$ (entier). L'opposition de phase correspond à $d = \left(k+\tfrac12\right)\lambda$.
\item \textbf{Vrai.} $\lambda = \dfrac{c}{f} = \dfrac{\SI{340}{m/s}}{\SI{200}{Hz}} = \SI{1,7}{m}$.
\item \textbf{Faux.} L'immobilité apparente est obtenue dès que $f_s = \dfrac{f}{n}$ ($n$ entier) : la fréquence de rotation peut être $\SI{50}{Hz}$, mais aussi $\SI{100}{Hz}$, $\SI{150}{Hz}$, etc. Seule la \emph{plus grande} fréquence d'éclairs donnant une immobilité unique est égale à $f$.
\item \textbf{Vrai.} C'est une propriété fondamentale : la somme de deux grandeurs sinusoïdales de même pulsation $\omega$ est sinusoïdale de même pulsation $\omega$ (le vecteur de Fresnel somme tourne à la vitesse $\omega$).
\item \textbf{Faux.} Une onde mécanique transporte de l'\emph{énergie} sans transport de matière : les particules du milieu oscillent autour de leur position d'équilibre (un bouchon sur l'eau monte et descend sans suivre l'onde).
\end{enumerate}
\end{corrige}

\begin{corrige}[Exercice 3 — Définitions et vocabulaire]
\begin{enumerate}
\item \textbf{Période} $T$ : plus petite durée au bout de laquelle le phénomène se reproduit identique à lui-même (en \SI{}{s}). \textbf{Fréquence} $f = 1/T$ : nombre de périodes par seconde (en \SI{}{Hz}). \textbf{Amplitude} $X_m$ : valeur maximale de la grandeur sinusoïdale. \textbf{Phase à l'origine des dates} $\varphi$ : valeur de l'argument $(\omega t + \varphi)$ à la date $t = 0$ (en \SI{}{rad}).
\item Une \textbf{onde mécanique progressive} est la propagation d'une perturbation dans un milieu matériel, avec transport d'énergie et sans transport de matière. \emph{Transversales :} onde sur une corde tendue, ondes à la surface de l'eau (cuve à ondes). \emph{Longitudinales :} onde sonore dans l'air, onde de compression le long d'un ressort (slinky).
\item La \textbf{double périodicité} désigne : la périodicité \emph{temporelle} (en un point fixe, le mouvement se répète tous les $T$) et la périodicité \emph{spatiale} (à date fixe, le profil se répète tous les $\lambda$). Relation : $\lambda = cT = \dfrac{c}{f}$.
\item En phase : $d = k\lambda$ avec $k \in \mathbb{Z}$. En opposition de phase : $d = \left(k + \dfrac{1}{2}\right)\lambda$ avec $k \in \mathbb{Z}$.
\end{enumerate}
\end{corrige}

\begin{corrige}[Exercice 4 — Calculs directs]
\begin{enumerate}
\item Par identification avec $x(t) = X_m\cos(\omega t + \varphi)$ : amplitude $X_m = \SI{4}{cm}$ ; pulsation $\omega = 20\pi \approx \SI{62,8}{rad/s}$ ; période $T = \dfrac{2\pi}{\omega} = \dfrac{2\pi}{20\pi} = \SI{0,10}{s}$ ; fréquence $f = \dfrac{1}{T} = \SI{10}{Hz}$ ; phase à l'origine $\varphi = -\dfrac{\pi}{3}\ \text{rad}$.
\item $\lambda = \dfrac{c}{f} = \dfrac{\SI{340}{m/s}}{\SI{500}{Hz}} = \SI{0,68}{m}$.
\item $c = \dfrac{SM}{\theta} = \dfrac{\SI{2,0}{m}}{\SI{0,25}{s}} = \SI{8,0}{m/s}$.
\item $f = \dfrac{1500}{60} = \SI{25}{Hz}$ (25 tours par seconde).
\end{enumerate}
\end{corrige}

\begin{corrige}[Exercice 5 — Exploitation d'un oscillogramme]
\begin{enumerate}
\item Période : $T = 4\ \text{div} \times \SI{5}{ms/div} = \SI{20}{ms} = \SI{2,0e-2}{s}$. Fréquence : $f = \dfrac{1}{T} = \dfrac{1}{\SI{0,020}{s}} = \SI{50}{Hz}$.
\item La hauteur crête-à-crête vaut $6$ div, donc l'amplitude correspond à $3$ div : $U_m = 3 \times \SI{2}{V} = \SI{6}{V}$. Valeur efficace : $U_{\text{eff}} = \dfrac{U_m}{\sqrt{2}} = \dfrac{\SI{6}{V}}{\sqrt{2}} \approx \SI{4,24}{V}$.
\item Pulsation : $\omega = 2\pi f = 100\pi \approx \SI{314}{rad/s}$. À $t=0$, la courbe passe par un maximum : $u(0) = U_m$, donc $\varphi = 0$ (fonction cosinus). D'où :
\[ u(t) = \SI{6}{V}\cos(100\pi t). \]
\item À $t = 0$, le vecteur de Fresnel $\vec{U}$ fait l'angle $\varphi = 0$ avec $Ox$ : il est porté par l'axe $Ox$, de module $\SI{6}{V}$, soit un segment de $\SI{6}{cm}$ dirigé selon $Ox$ (échelle \SI{1}{cm/V}).
\end{enumerate}
\textbf{Point de vigilance :} ne pas confondre hauteur crête-à-crête ($2U_m$) et amplitude ($U_m$) : ici $6$ div crête-à-crête donnent $U_m = \SI{6}{V}$ et non $\SI{12}{V}$.
\end{corrige}

\begin{corrige}[Exercice 6 — Stroboscopie : disque d'un moteur]
\begin{enumerate}
\item La plus grande fréquence d'éclairs donnant une immobilité unique est égale à la fréquence de rotation : $f = f_s = \SI{50}{Hz}$. Vitesse de rotation : $50 \times 60 = \SI{3000}{tr/min}$.
\item Avec $f_s = \SI{25}{Hz} = \dfrac{f}{2}$ : entre deux éclairs, le disque effectue exactement $2$ tours, donc le repère est éclairé toujours à la même position. On observe \textbf{un repère immobile unique} (immobilité apparente), mais ce n'est pas la plus grande fréquence possible.
\item Avec $f_s = \SI{51}{Hz}$, légèrement supérieure à $f = \SI{50}{Hz}$ : entre deux éclairs, le repère n'a pas tout à fait achevé son tour ; il paraît reculer très lentement. On observe un \textbf{mouvement ralenti apparent en sens inverse} du sens réel, de fréquence apparente :
\[ f_{\text{app}} = |f_s - f| = |\SI{51}{Hz} - \SI{50}{Hz}| = \SI{1}{Hz}. \]
Le repère semble effectuer un tour (en sens inverse) en $\SI{1}{s}$.
\end{enumerate}
\end{corrige}

\begin{corrige}[Exercice 7 — Somme de deux tensions par la construction de Fresnel]
\begin{enumerate}
\item À $t = 0$ : $\vec{V}_1$ a pour module $\SI{6}{V}$ (soit $\SI{6}{cm}$) et fait l'angle $\varphi_1 = 0$ avec $Ox$ ; $\vec{V}_2$ a pour module $\SI{8}{V}$ (soit $\SI{8}{cm}$) et fait l'angle $\varphi_2 = \dfrac{\pi}{2}$ avec $Ox$ (porté par $Oy$). On construit $\vec{V} = \vec{V}_1 + \vec{V}_2$ par la règle du parallélogramme (ici un rectangle).
\item Les deux vecteurs étant perpendiculaires, le théorème de Pythagore donne :
\[ U_m = \sqrt{U_{1m}^2 + U_{2m}^2} = \sqrt{6^2 + 8^2} = \sqrt{100} = \SI{10}{V}. \]
Sur la construction, on mesure bien un segment de $\SI{10}{cm}$.
\item Phase à l'origine :
\[ \tan\varphi = \frac{U_{2m}\sin\varphi_2}{U_{1m}\cos\varphi_1} = \frac{8}{6} \implies \varphi = \arctan\left(\frac{4}{3}\right) \approx \SI{0,927}{rad} \approx 53,1^\circ. \]
Équation horaire : $u(t) = \SI{10}{V}\cos\left(100\pi t + 0,927\right)$.
\item Si les deux tensions étaient en opposition de phase ($\varphi_2 - \varphi_1 = \pi$), les vecteurs de Fresnel seraient colinéaires de sens opposés :
\[ U_m = |U_{2m} - U_{1m}| = |8 - 6| = \SI{2}{V}. \]
\end{enumerate}
\end{corrige}

\begin{corrige}[Exercice 8 — Onde le long d'une corde]
\begin{enumerate}
\item Pulsation : $\omega = 40\pi\ \text{rad/s}$. Période : $T = \dfrac{2\pi}{\omega} = \dfrac{2\pi}{40\pi} = \SI{0,050}{s} = \SI{50}{ms}$. Fréquence : $f = \dfrac{1}{T} = \SI{20}{Hz}$. Longueur d'onde : $\lambda = cT = \SI{8}{m/s} \times \SI{0,050}{s} = \SI{0,40}{m} = \SI{40}{cm}$.
\item Le point $M$ reproduit le mouvement de $S$ avec le retard $\theta = \dfrac{d}{c} = \dfrac{\SI{0,60}{m}}{\SI{8}{m/s}} = \SI{0,075}{s}$ :
\[ y_M(t) = y_S(t - \theta) = \SI{2}{cm}\cos\left[40\pi\left(t - 0,075\right)\right] = \SI{2}{cm}\cos(40\pi t - 3\pi). \]
(On peut aussi écrire $y_M(t) = \SI{2}{cm}\cos(40\pi t - \pi)$ car $-3\pi \equiv -\pi \pmod{2\pi}$.)
\item Différence de phase entre $M$ et $S$ : $\Delta\varphi = -\dfrac{2\pi d}{\lambda} = -\dfrac{2\pi \times 0,60}{0,40} = -3\pi \equiv \pi \pmod{2\pi}$.
Vérification : $\dfrac{d}{\lambda} = \dfrac{0,60}{0,40} = 1,5 = 1 + \dfrac{1}{2}$, demi-entier. Donc $M$ et $S$ vibrent \textbf{en opposition de phase}.
\item Points en phase avec $S$ : $d = k\lambda = 0,40\,k$ (en m), soit $d = 0\ ;\ 0,40\ ;\ 0,80\ ;\ 1,20\ \text{m}\ldots$
Points en opposition de phase avec $S$ : $d = \left(k+\dfrac{1}{2}\right)\lambda = 0,20 + 0,40\,k$ (en m), soit $d = 0,20\ ;\ 0,60\ ;\ 1,00\ \text{m}\ldots$ On retrouve bien $M$ à $0,60\ \text{m}$.
\end{enumerate}
\end{corrige}

\begin{corrige}[Exercice 9 — Type Bac : étude d'une onde sur une corde vibrante]
\textbf{Partie A}
\begin{enumerate}
\item La plus grande fréquence d'éclairs donnant une immobilité apparente unique est égale à la fréquence du phénomène observé. Or $f_s = \SI{100}{Hz} = f$ : l'observation confirme la fréquence du vibreur.
\item La distance entre deux points consécutifs vibrant en phase est la \textbf{longueur d'onde} : $\lambda = \SI{4,0}{cm} = \SI{0,040}{m}$. Célérité :
\[ c = \lambda f = \SI{0,040}{m} \times \SI{100}{Hz} = \SI{4,0}{m/s}. \]
\end{enumerate}
\textbf{Partie B}
\begin{enumerate}
\item Retard de $M$ par rapport à $S$ : $\theta = \dfrac{SM}{c} = \dfrac{\SI{0,30}{m}}{\SI{4,0}{m/s}} = \SI{0,075}{s}$. D'où :
\[ y_M(t) = \SI{3}{cm}\cos\left[200\pi(t - 0,075)\right] = \SI{3}{cm}\cos(200\pi t - 15\pi) = \SI{3}{cm}\cos(200\pi t - \pi). \]
\item $\dfrac{SM}{\lambda} = \dfrac{\SI{30}{cm}}{\SI{4,0}{cm}} = 7,5 = 7 + \dfrac{1}{2}$ : demi-entier. Le point $M$ vibre \textbf{en opposition de phase} avec $S$ (cohérent avec la phase $-\pi$ trouvée ci-dessus).
\item Points en opposition de phase avec $S$ : $SM = \left(k+\dfrac{1}{2}\right)\lambda = 2 + 4k$ (en cm), avec $SM \leq L = \SI{200}{cm}$.
$2 + 4k \leq 200 \iff k \leq 49,5$, donc $k$ varie de $0$ à $49$ : il y a \textbf{50 points}, situés à $SM = 2\ ;\ 6\ ;\ 10\ ;\ \ldots\ ;\ 198\ \text{cm}$.
\end{enumerate}
\textbf{Partie C}
\begin{enumerate}
\item $M$ et $N$ vibrent en opposition de phase, donc $MN = \left(k+\dfrac{1}{2}\right)\lambda$ avec $k \in \mathbb{N}$ :
\[ \SI{45}{cm} = \frac{2k+1}{2}\,\lambda \implies \lambda = \frac{2 \times 45}{2k+1} = \frac{90}{2k+1} \quad (\text{en cm}). \]
\item Célérité $c = \lambda f$ avec $f = \SI{100}{Hz}$ :
\begin{itemize}
\item $k = 0$ : $\lambda = \SI{90}{cm} = \SI{0,90}{m}$, $c = \SI{0,90}{m} \times \SI{100}{Hz} = \SI{90}{m/s}$ ;
\item $k = 1$ : $\lambda = \dfrac{90}{3} = \SI{30}{cm} = \SI{0,30}{m}$, $c = \SI{30}{m/s}$ ;
\item $k = 2$ : $\lambda = \dfrac{90}{5} = \SI{18}{cm} = \SI{0,18}{m}$, $c = \SI{18}{m/s}$.
\end{itemize}
\end{enumerate}
\textbf{Barème indicatif (10 pts) :} Partie A : 2 pts ; Partie B : 4 pts (équation horaire 1,5 pt) ; Partie C : 4 pts (démonstration 2 pts).
\textbf{Points de vigilance :} citer la condition $d = \left(k+\tfrac12\right)\lambda$ avant de conclure ; convertir les cm en m pour le calcul de $c$ ; dans la Partie C, l'opposition de phase ne fixe pas une valeur unique de $\lambda$ : toutes les valeurs $\dfrac{90}{2k+1}$ sont possibles.
\end{corrige}

\begin{corrige}[Exercice 10 — Type Bac : cuve à ondes et mesure à l'oscilloscope]
\textbf{Partie A}
\begin{enumerate}
\item $T = \dfrac{1}{f} = \dfrac{1}{\SI{25}{Hz}} = \SI{0,040}{s} = \SI{40}{ms}$.
\item La longueur d'onde $\lambda$ est la plus petite distance séparant deux points du milieu vibrant en phase (période spatiale de l'onde) ; c'est aussi la distance parcourue par l'onde pendant une période : $\lambda = cT = \dfrac{c}{f}$.
\end{enumerate}
\textbf{Partie B}
\begin{enumerate}
\item Retard : $\theta = 1,5\ \text{div} \times \SI{10}{ms/div} = \SI{15}{ms} = \SI{0,015}{s}$.
\item Distance $M_1M_2 = SM_2 - SM_1 = \SI{10,5}{cm} - \SI{6,0}{cm} = \SI{4,5}{cm} = \SI{0,045}{m}$. Célérité :
\[ c = \frac{M_1M_2}{\theta} = \frac{\SI{0,045}{m}}{\SI{0,015}{s}} = \SI{3,0}{m/s}. \]
\item $\lambda = cT = \SI{3,0}{m/s} \times \SI{0,040}{s} = \SI{0,12}{m} = \SI{12}{cm}$. Vérification : $\dfrac{c}{f} = \dfrac{\SI{3,0}{m/s}}{\SI{25}{Hz}} = \SI{0,12}{m}$ : la relation $\lambda = cT$ est bien vérifiée.
\item $\dfrac{M_1M_2}{\lambda} = \dfrac{4,5}{12} = 0,375$ : ni entier ni demi-entier. Les points $M_1$ et $M_2$ ne vibrent \textbf{ni en phase ni en opposition de phase}. Différence de phase : $\Delta\varphi = -\dfrac{2\pi \times 4,5}{12} = -\dfrac{3\pi}{4}\ \text{rad}$.
\end{enumerate}
\textbf{Partie C}
\begin{enumerate}
\item Avec $f_s = \SI{25}{Hz} = f$ : entre deux éclairs, l'onde avance d'exactement une longueur d'onde ; le profil observé est identique à chaque éclair. La surface de l'eau paraît \textbf{immobile} (rides fixes).
\item Avec $f_s = \SI{24}{Hz}$, légèrement inférieure à $f$ : entre deux éclairs, l'onde a eu le temps d'avancer d'un peu plus d'une longueur d'onde ; les rides paraissent avancer lentement dans le sens réel de propagation. Fréquence apparente :
\[ f_{\text{app}} = |f_s - f| = |\SI{24}{Hz} - \SI{25}{Hz}| = \SI{1}{Hz}. \]
On observe un \textbf{mouvement ralenti apparent} des rides à $\SI{1}{Hz}$ dans le sens réel.
\end{enumerate}
\textbf{Barème indicatif (10 pts) :} Partie A : 2 pts ; Partie B : 5 pts ; Partie C : 3 pts.
\textbf{Points de vigilance :} le retard se lit sur le décalage \emph{horizontal} de l'oscillogramme ($\theta = 1,5 \times b$) ; la distance à utiliser pour $c$ est $M_1M_2$ et non $SM_2$ ; pour conclure sur l'état vibratoire, comparer $M_1M_2$ à $\lambda$ et citer les critères.
\end{corrige}

\begin{corrige}[Exercice 11 — Type Bac : synthèse oscillations et ondes — le haut-parleur du bal du lycée]
\textbf{Partie A}
\begin{enumerate}
\item Amplitude : $U_m = \SI{10}{V}$. Valeur efficace : $U_{\text{eff}} = \dfrac{10}{\sqrt{2}} \approx \SI{7,07}{V}$. Pulsation $\omega = 2000\pi\ \text{rad/s}$, donc $f = \dfrac{\omega}{2\pi} = \SI{1000}{Hz} = \SI{1,0}{kHz}$ et $T = \dfrac{1}{f} = \SI{1,0e-3}{s} = \SI{1,0}{ms}$.
\item À $t = 0$, $\varphi = 0$ : le vecteur de Fresnel $\vec{U}$ est porté par $Ox$, de module $\SI{10}{V}$, soit un segment de $\SI{5}{cm}$ selon $Ox$ (échelle $\SI{1}{cm} \leftrightarrow \SI{2}{V}$).
\item $\vec{U}$ : module $\SI{10}{V}$, angle $0$ ; $\vec{U}'$ : module $\SI{10}{V}$, angle $\dfrac{\pi}{2}$. Les deux vecteurs sont perpendiculaires et de même module. Par Pythagore :
\[ U_{\text{somme},m} = \sqrt{10^2 + 10^2} = 10\sqrt{2} \approx \SI{14,1}{V}. \]
Phase : $\tan\varphi = \dfrac{10\sin(\pi/2)}{10\cos 0} = 1$, avec $\varphi$ dans le premier quadrant : $\varphi = \dfrac{\pi}{4}\ \text{rad}$. D'où :
\[ u(t) + u'(t) = 10\sqrt{2}\cos\left(2000\pi t + \frac{\pi}{4}\right) \quad (\text{en volts}). \]
\end{enumerate}
\textbf{Partie B}
\begin{enumerate}
\item Le son dans l'air est une onde \textbf{longitudinale} : les couches d'air subissent des compressions et dilatations \emph{parallèles} à la direction de propagation.
\item $\lambda = \dfrac{c}{f} = \dfrac{\SI{340}{m/s}}{\SI{1000}{Hz}} = \SI{0,34}{m} = \SI{34}{cm}$.
\item Distance entre les deux élèves : $AB = d_B - d_A = \SI{6,80}{m} - \SI{5,10}{m} = \SI{1,70}{m}$. Rapport :
\[ \frac{AB}{\lambda} = \frac{1,70}{0,34} = 5 \quad (\text{entier}). \]
Donc $AB = 5\lambda$ : les deux élèves perçoivent le son \textbf{en phase}.
\item Retard : $\theta = \dfrac{AB}{c} = \dfrac{\SI{1,70}{m}}{\SI{340}{m/s}} = \SI{5,0e-3}{s} = \SI{5,0}{ms}$. (Vérification : $\theta = 5T$, cohérent avec $AB = 5\lambda$.)
\end{enumerate}
\textbf{Partie C}
\begin{enumerate}
\item Pour que la membrane paraisse immobile et unique, il faut $f_s = f = \SI{1000}{Hz}$ (plus grande fréquence d'immobilité) ; toute sous-fréquence $f_s = \dfrac{1000}{n}\ \text{Hz}$ donne aussi une immobilité apparente.
\item Un stroboscope classique plafonne vers quelques centaines de hertz. Méthodes alternatives :
\begin{itemize}
\item brancher la tension du GBF (ou d'un microphone placé devant le haut-parleur) sur un \textbf{oscilloscope} et mesurer la période $T$ sur l'oscillogramme, puis $f = 1/T$ ;
\item utiliser un \textbf{microphone relié à une interface d'acquisition} (ordinateur, carte Sysam) et mesurer la fréquence du signal enregistré.
\end{itemize}
\end{enumerate}
\textbf{Barème indicatif (12 pts) :} Partie A : 4 pts (construction de Fresnel 1,5 pt) ; Partie B : 5 pts ; Partie C : 3 pts.
\textbf{Points de vigilance :} justifier le caractère longitudinal du son par la nature de la perturbation (pression) ; pour la question 6, comparer $AB$ (et non $d_A$ ou $d_B$ séparément) à $\lambda$ ; pour la question 3, préciser le quadrant de $\varphi$ (ici $\varphi = \pi/4$ et non $5\pi/4$ car les deux projections sont positives).
\end{corrige}

\newpage

\coursec{Fiche bilan}

\begin{popfigure}
\begin{tikzpicture}[mindmap, grow cyclic, every node/.style={concept, font=\small},
  concept color=popblue!60,
  level 1/.append style={level distance=3.4cm, sibling angle=60, font=\small},
  level 2/.append style={level distance=2.1cm, sibling angle=45, font=\scriptsize},
  scale=0.82, transform shape]
\node[concept color=popblue!70, font=\bfseries] {Oscillations et ondes mécaniques}
  child[concept color=poporange!70] { node {Phénomènes périodiques}
    child { node {Amplitude $X_m$, période $T$} }
    child { node {$f=\dfrac{1}{T}$, $\omega=2\pi f$} }
    child { node {Phase $\varphi$} }
  }
  child[concept color=poppink!70] { node {Stroboscopie}
    child { node {$f_e=f$ : immobilité} }
    child { node {$f_e\lesssim f$ : ralenti} }
  }
  child[concept color=poppurple!70] { node {Fresnel}
    child { node {Vecteur $(X_m,\varphi)$} }
    child { node {Somme même fréquence} }
  }
  child[concept color=popgreen!70] { node {Onde mécanique}
    child { node {Transversale} }
    child { node {Longitudinale} }
    child { node {$c=\dfrac{d}{\theta}$} }
  }
  child[concept color=popgold!80] { node {Onde sinusoïdale}
    child { node {$\lambda=cT$} }
    child { node {$y_M(t)=y_S\!\left(t-\dfrac{x}{c}\right)$} }
    child { node {Phase : $\Delta\varphi=\dfrac{2\pi d}{\lambda}$} }
  };
\end{tikzpicture}
\legende{Carte mentale de la leçon : systèmes oscillants et propagation des ondes mécaniques.}
\end{popfigure}

\begin{bilan}
\textbf{1. Grandeurs sinusoïdales}
\begin{itemize}
  \item Une grandeur \cle{sinusoïdale} s'écrit $x(t)=X_m\cos(\omega t+\varphi)$ : $X_m$ = \cle{amplitude} (m), $\omega$ = \cle{pulsation} (rad/s), $\varphi$ = \cle{phase à l'origine} (rad).
  \item \cle{Période} $T$ (s) : plus petite durée de répétition ; \cle{fréquence} $f=\dfrac{1}{T}$ (Hz) ; $\omega=2\pi f=\dfrac{2\pi}{T}$.
  \item Vérification dimensionnelle : $[\omega t]=\text{rad}$, sans dimension ; $[X_m]=[x]$.
\end{itemize}

\textbf{2. Stroboscopie}
\begin{itemize}
  \item Stroboscope : source d'éclairs brefs de fréquence $f_e$ réglable.
  \item $f_e=\dfrac{f}{k}$ ($k\in\mathbb{N}^*$) : \cle{immobilité apparente}.
  \item $f_e$ légèrement inférieure à $f$ : \cle{mouvement ralenti apparent} dans le sens réel ; $f_e$ légèrement supérieure : ralenti en sens inverse.
\end{itemize}

\textbf{3. Représentation de Fresnel}
\begin{itemize}
  \item À $x(t)=X_m\cos(\omega t+\varphi)$ on associe le vecteur $\overrightarrow{OM}$ de norme $X_m$ faisant l'angle $(\omega t+\varphi)$ avec l'axe, tournant à la vitesse $\omega$.
  \item La somme de deux grandeurs \textbf{de même fréquence} est sinusoïdale de même fréquence : on additionne les vecteurs de Fresnel.
\end{itemize}

\textbf{4. Ondes mécaniques progressives}
\begin{itemize}
  \item \cle{Onde mécanique} : propagation d'une perturbation dans un milieu matériel, \textbf{sans transport de matière} mais avec transport d'énergie.
  \item \cle{Transversale} : perturbation $\perp$ direction de propagation (corde, surface de l'eau). \cle{Longitudinale} : perturbation $\parallel$ (ressort, son).
  \item \cle{Célérité} : $c=\dfrac{d}{\theta}$ où $\theta$ est le retard du point $M$ situé à la distance $d$ de la source.
\end{itemize}

\textbf{5. Onde progressive sinusoïdale : double périodicité}
\begin{center}
\begin{tabularx}{\linewidth}{|l|X|X|}\hline
\rowcolor{popblueL} \textbf{Relation} & \textbf{Expression} & \textbf{Signification} \\ \hline
Périodicité temporelle & $y(t+T)=y(t)$ & chaque point vibre avec la période $T$ \\ \hline
Périodicité spatiale & $y(x+\lambda)=y(x)$ & $\lambda=cT$ : longueur d'onde (m) \\ \hline
Retard & $\theta=\dfrac{d}{c}$ & $M$ reproduit le mouvement de $S$ avec le retard $\theta$ \\ \hline
Équation de $M$ & $y_M(t)=a\cos\left[\omega\left(t-\dfrac{x}{c}\right)+\varphi\right]$ & si $y_S(t)=a\cos(\omega t+\varphi)$ \\ \hline
Déphasage & $\Delta\varphi=\dfrac{2\pi d}{\lambda}$ & $d=k\lambda$ : en phase ; $d=(2k+1)\dfrac{\lambda}{2}$ : opposition \\ \hline
\end{tabularx}
\end{center}

\textbf{6. Méthodes express}
\begin{itemize}
  \item \textbf{Oscillogramme} : $T=$ (nombre de divisions) $\times$ (balayage) ; $X_m=$ (divisions verticales) $\times$ (sensibilité).
  \item \textbf{Comparer deux points} : calculer $d$, puis $\dfrac{d}{\lambda}$ : entier $\Rightarrow$ en phase ; demi-entier $\Rightarrow$ opposition de phase.
  \item \textbf{Trouver $\lambda$ sur une photo} : distance entre deux crêtes consécutives ; puis $c=\dfrac{\lambda}{T}=\lambda f$.
\end{itemize}
\end{bilan}

\begin{aretenir}[Checklist avant le Bac]
\begin{itemize}
  \item[$\square$] Je sais écrire une grandeur sinusoïdale $x(t)=X_m\cos(\omega t+\varphi)$ et identifier $X_m$, $\omega$, $\varphi$.
  \item[$\square$] Je connais par cœur $f=\dfrac{1}{T}$ et $\omega=2\pi f=\dfrac{2\pi}{T}$ avec les unités SI.
  \item[$\square$] Je sais exploiter un oscillogramme (balayage, sensibilité) pour mesurer $T$ et $X_m$.
  \item[$\square$] Je sais prédire l'aspect stroboscopique : immobilité si $f_e=f/k$, ralenti sinon.
  \item[$\square$] Je sais construire le vecteur de Fresnel et additionner deux grandeurs de même fréquence.
  \item[$\square$] Je distingue onde transversale et onde longitudinale avec un exemple chacune.
  \item[$\square$] Je sais qu'une onde transporte de l'énergie sans transporter de matière.
  \item[$\square$] Je maîtrise $\theta=\dfrac{d}{c}$ et $\lambda=cT=\dfrac{c}{f}$ (avec analyse dimensionnelle).
  \item[$\square$] Je sais écrire $y_M(t)=a\cos\left[\omega\left(t-\dfrac{x}{c}\right)+\varphi\right]$ à partir de $y_S(t)$.
  \item[$\square$] Je sais comparer deux points : $d=k\lambda$ (en phase), $d=(2k+1)\lambda/2$ (opposition de phase).
  \item[$\square$] Je vérifie toujours l'homogénéité des formules et j'exprime les résultats avec unité et chiffres significatifs.
  \item[$\square$] Je pense à convertir : distances en mètres, durées en secondes, fréquences en hertz.
\end{itemize}
\end{aretenir}

\findecours

\end{document}
